\documentclass[12pt,a4paper,oneside,openany]{book}

% Компилировать через XeLaTeX.
\usepackage{fontspec}
\usepackage{polyglossia}
\setmainlanguage{russian}
\setotherlanguage{english}
\setmainfont{DejaVu Serif}[
    Ligatures=TeX,
    Scale=1.02
]
\setsansfont{DejaVu Sans}[
    Ligatures=TeX,
    Scale=MatchLowercase
]
\setmonofont{DejaVu Sans Mono}[
    Scale=MatchLowercase
]

% Математика
\usepackage{amsmath,amssymb,amsthm,mathtools}
\allowdisplaybreaks

% Страница
\usepackage[
    left=24mm,
    right=21mm,
    top=24mm,
    bottom=26mm,
    headsep=8mm
]{geometry}

% Единое оформление текста
\usepackage{indentfirst}
\usepackage{enumitem}
\usepackage{microtype}
\usepackage{xcolor}
\usepackage{setspace}
\onehalfspacing
\setlength{\parindent}{1.25em}
\setlength{\parskip}{0pt}
\emergencystretch=2em

\definecolor{logicblue}{HTML}{1F4E79}
\definecolor{logicteal}{HTML}{2F6F73}
\definecolor{logicgray}{HTML}{5F6872}
\definecolor{logicline}{HTML}{D8DEE6}

\setlist{
    topsep=0.45\baselineskip,
    itemsep=0.2\baselineskip,
    parsep=0pt,
    leftmargin=*
}

% Оглавление
\usepackage{tocloft}
\setcounter{tocdepth}{2}
\setcounter{secnumdepth}{3}
\renewcommand{\cftchapfont}{\bfseries\color{logicblue}}
\renewcommand{\cftchappagefont}{\bfseries\color{logicblue}}
\renewcommand{\cftchapleader}{\cftdotfill{\cftdotsep}}
\setlength{\cftbeforechapskip}{0.55em}

% Заголовки
\usepackage{titlesec}
\titleformat{\chapter}[display]
  {\normalfont\sffamily\bfseries\color{logicblue}\raggedright}
  {\Large\MakeUppercase{\chaptername}\ \thechapter}
  {0.6em}
  {\Huge}
  [\vspace{0.4em}{\titlerule[0.8pt]}]
\titlespacing*{\chapter}{0pt}{-8pt}{28pt}

\titleformat{name=\chapter,numberless}[display]
  {\normalfont\sffamily\bfseries\color{logicblue}\raggedright}
  {}
  {0pt}
  {\Huge}
  [\vspace{0.4em}{\titlerule[0.8pt]}]
\titlespacing*{name=\chapter,numberless}{0pt}{-8pt}{24pt}

\titleformat{\section}
  {\normalfont\Large\sffamily\bfseries\color{logicteal}}
  {\thesection}
  {0.75em}
  {}
\titlespacing*{\section}{0pt}{1.4\baselineskip}{0.55\baselineskip}

\titleformat{\subsection}
  {\normalfont\large\sffamily\bfseries\color{logicblue}}
  {\thesubsection}
  {0.7em}
  {}
\titlespacing*{\subsection}{0pt}{1.1\baselineskip}{0.35\baselineskip}

\titleformat{\subsubsection}
  {\normalfont\normalsize\sffamily\bfseries\color{logicgray}}
  {\thesubsubsection}
  {0.65em}
  {}

% Колонтитулы
\usepackage{fancyhdr}
\pagestyle{fancy}
\fancyhf{}
\fancyhead[L]{\sffamily\small\color{logicgray}Математическая логика}
\fancyhead[R]{\sffamily\small\color{logicgray}\nouppercase{\leftmark}}
\fancyfoot[C]{\sffamily\small\color{logicgray}\thepage}
\renewcommand{\headrulewidth}{0.4pt}
\renewcommand{\headrule}{\hbox to\headwidth{\color{logicline}\leaders\hrule height \headrulewidth\hfill}}
\setlength{\headheight}{25pt}

% Теоремные окружения
\theoremstyle{definition}
\newtheorem{definition}{Определение}[chapter]
\newtheorem{example}{Пример}[chapter]

\theoremstyle{plain}
\newtheorem{theorem}{Теорема}[chapter]
\newtheorem{lemma}{Лемма}[chapter]
\newtheorem{proposition}{Утверждение}[chapter]
\newtheorem{corollary}{Следствие}[chapter]

\theoremstyle{remark}
\newtheorem{remark}{Замечание}[chapter]

% Команды
\newcommand{\N}{\mathbb{N}}
\newcommand{\Z}{\mathbb{Z}}
\newcommand{\Q}{\mathbb{Q}}
\newcommand{\R}{\mathbb{R}}
\newcommand{\modelsrel}{\models}
\newcommand{\proves}{\vdash}

% Аргумент #1 позволяет менять ширину, по умолчанию она 1.2ex
\newcommand{\blank}[1][1.2ex]{%
  \mbox{%
    \vrule width 0.1ex height 0.4ex depth 0pt % Левый бортик (0.3ex — очень мало)
    \vrule width #1 height 0.1ex depth 0pt    % Нижняя линия
    \vrule width 0.1ex height 0.4ex depth 0pt % Правый бортик
  }%
}

\title{\sffamily\bfseries\color{logicblue}Полный курс математической логики}
\author{Ямщиков Герман M3233}
\date{от 31 августа 2026}

\usepackage[unicode=true]{hyperref}
\hypersetup{
    colorlinks=true,
    linkcolor=logicblue,
    citecolor=logicblue,
    urlcolor=logicteal,
    pdftitle={Полный курс математической логики},
    pdfauthor={Yamshchikov German M3233}
}

\begin{document}

\maketitle
\tableofcontents
\chapter*{Гл. 1-6 + 10 покрывют курс  "Математическая логика"}
\addcontentsline{toc}{chapter}{Гл. 1-6 + 10 покрывют курс Д.Штукенберга "Математическая логика"}

\chapter{Формальные языки}

\section{Формализация математики}
\section{Язык и метаязык}
\section{Синтаксис формул}
\section{Термы и подстановки}
\section{Свободные и связанные переменные}
\section{Формальные теории}

\section{Задачник к главе}

\chapter{Пропозициональная логика}

\section{Исчисление высказываний}
\section{Семантика и таблицы истинности}
\section{Аксиоматические системы}
\section{Теорема о дедукции}
\section{Корректность}
\section{Полнота}

\section{Задачник к главе}

\chapter{Логика первого порядка}

\section{Предикаты и кванторы}
\section{Язык первого порядка}
\section{Семантика первого порядка}
\section{Формальные выводы}
\section{Полнота Гёделя}
\section{Компактность}
\section{Теорема Лёвенгейма--Сколема}

\section{Задачник к главе}

\chapter{Формальная арифметика}

\section{Арифметика Робинсона}
\section{Арифметика Пеано}
\section{Представимые функции}
\section{Гёделева нумерация}
\section{Диагональный метод}

\section{Задачник к главе}

\chapter{Границы формальных систем}

\section{Первая теорема Гёделя}
\section{Вторая теорема Гёделя}
\section{Теорема Тарского}
\section{Теорема Черча}
\section{Нестандартная арифметика}

\section{Задачник к главе}


\chapter{Теория доказательств}

\section{Гильбертовы системы}
\section{Натуральный вывод}
\section{Секвенциальное исчисление}
\section{Устранимость сечения}
\section{Нормализация}
\section{Конструктивная логика}
\section{Интуиционистская логика}
\section{Семантика Крипке}

\section{Задачник к главе}

\chapter*{Глава 7 покрывает курс Р.Попкова "Между алгеброй и логикой"}
\addcontentsline{toc}{chapter}{Глава 7 покрывает курс Р.Попкова "Между алгеброй и логикой"}

\chapter{Теория моделей}

\section{Основы теории моделей}

\subsection{Язык (Сигнатура)}

Язык $L$ (также называемый словарем или сигнатурой) определяется набором символов трех типов:
\begin{enumerate}
    \item \textbf{Множество символов отношений}: Каждому символу $R$ сопоставлено натуральное число $n$ — его \textit{арность} или \textit{местность}.
    \item \textbf{Множество символов функций}: Каждому символу $f$ сопоставлена арность $n$.
    \item \textbf{Множество символов констант}: Можно рассматривать как функции арности 0.
\end{enumerate}

\subsection{\texorpdfstring{$L$}{L}-структуры}

$L$-структура $\mathfrak{A}$ состоит из множества $A$ (называемого \textbf{носителем}) и интерпретаций всех символов языка $L$:
\begin{itemize}
    \item Для каждого символа отношения $R$ арности $n$ интерпретация — это подмножество:
    \[ R^{\mathfrak{A}} \subseteq A^n \]
    \item Для каждого символа функции $f$ арности $n$ интерпретация — это отображение:
    \[ f^{\mathfrak{A}}: A^n \to A \]
    \item Для каждого символа константы $c$ интерпретация — это элемент носителя:
    \[ c^{\mathfrak{A}} \in A \]
\end{itemize}
\textbf{Замечание.}
Верхний индекс $\mathfrak A$ указывает, что рассматривается интерпретация соответствующего символа языка в структуре $\mathfrak A$.

Следует различать символ языка и его интерпретацию.
Например, символ отношения $R$ является лишь частью сигнатуры и сам по себе не имеет смысла.
После выбора структуры $\mathfrak A$ он получает конкретную интерпретацию $R^{\mathfrak A}$.


Для символа отношения $R$ арности $n$ интерпретация
\[
R^{\mathfrak A}\subseteq A^n
\]
представляет собой множество всех $n$-кортежей элементов носителя, на которых отношение $R$ истинно.

Например, пусть носитель
\[
A=\{1,2,3\},
\]
а символ отношения $R$ двуместный ($n=2$). Тогда
\[
A^2=A\times A=
\{(1,1),(1,2),(1,3),(2,1),(2,2),(2,3),(3,1),(3,2),(3,3)\}.
\]

Если интерпретировать $R$ как отношение ``меньше'', то
\[
R^{\mathfrak A}
=
\{(1,2),(1,3),(2,3)\}.
\]

То есть из всех пар элементов носителя выбираются только те, для которых выполняется отношение $<$.
\subsection{Примеры}

\subsubsection{Кольца}
Язык колец (ассоциативных, коммутативных с единицей):
\[ L_{\text{ring}} = \langle +^2, \cdot^2, -^1, 0^0, 1^0 \rangle \]

Примеры структур для этого языка:
\begin{itemize}
    \item Целые числа: $\mathbb{Z}_{\text{ring}} = \langle \mathbb{Z}; L_{\text{ring}} \rangle = \langle \mathbb{Z}, +^{\mathbb{Z}}, \cdot^{\mathbb{Z}}, -^{\mathbb{Z}}, 0^{\mathbb{Z}}, 1^{\mathbb{Z}} \rangle$
    \item Вещественные числа: $\mathbb{R}_{\text{ring}} = \langle \mathbb{R}; L_{\text{ring}} \rangle = \langle \mathbb{R}, +^{\mathbb{R}}, \cdot^{\mathbb{R}}, -^{\mathbb{R}}, 0^{\mathbb{R}}, 1^{\mathbb{R}} \rangle$
\end{itemize}
Примеры интерпретаций сложения: \\
1) $+^{\mathbb{Z}}: \mathbb{Z} \times \mathbb{Z} \to \mathbb{Z}$.\\
2) $+^{\mathbb{R}}: \mathbb{R} \times \mathbb{R} \to \mathbb{R}$.

\subsubsection{Порядок}
Язык бинарного отношения порядка: $L_{<} = \langle <^2 \rangle$. \\
Структура $\mathbb{Z}_{<} = \langle \mathbb{Z}, <^{\mathbb{Z}} \rangle$, где интерпретация отношения:
\[ <^{\mathbb{Z}} : \{ (a, b) \mid a, b \in \mathbb{Z}, a < b \} \]

\subsubsection{Группы}
\begin{itemize}
    \item Язык групп (мультипликативный): $L_{\text{gp}} = \langle \cdot^2, (\blank^{-1})^1, 1^0 \rangle$
    \item Язык аддитивных групп: $L_{\text{adgp}} = \langle +, -\blank, 0 \rangle$
\end{itemize}
\subsection{Операции над языками и мощностями}

\subsubsection{Обогащение и обеднение}
Связь между структурами разных языков определяется через вложение сигнатур:
\begin{itemize}
    \item \textbf{Обеднение}: Если язык $L$ является подмножеством языка $L'$ ($L \subseteq L'$), то $L$-структура, полученная из $L'$-структуры путем забывания интерпретаций лишних символов, называется \textit{обеднением} $L^{-}$. 
    \begin{itemize}
        \item  $L_{adgp}$ —  обеднение $L_{ring}$.
    \end{itemize}
    \item \textbf{Обогащение}: Наоборот, структура в более богатом языке $L'$, полученная добавлением новых интерпретаций к существующей $L$-структуре, называется \textit{обогащением} $L^{+}$.
    \begin{itemize}
        \item  $L_{ring}$ —  обогащение $L_{adgp}$.
        \item Язык упорядоченных колец: $L_{o-ring} = L_{ring} \cup \{ <^2 \}$.
    \end{itemize}
\end{itemize}

\begin{itemize}
    \item $L_{s-ring} = \langle +^2, \cdot^2, 0^0, 1^0 \rangle$ — язык полуколец.
    \item $L_{mon} = \langle \cdot^2, 1^0 \rangle$ — язык моноидов.
    \item $L_{graph} = \langle R^2 \rangle$ — язык графов (один бинарный символ отношения).
\end{itemize}


\begin{enumerate}
    \item \textbf{Мощность носителя} $|\mathfrak{A}|$: Это обычная мощность множества $A$ - носителя.
    \item \textbf{Мощность языка (Сигнатурная мощность)} $\|L\|$:
    \[ \|L\| = \max(\omega, |L|) \]
    где $|L|$ — количество символов в языке, а $\omega$  — мощность счетного множества (натуральных чисел).
\end{enumerate}


Даже если в языке всего один символ (например, язык графов), количество различных логических формул, которые можно составить, всегда будет не меньше счетного (за счет использования бесконечного запаса переменных $x_1, x_2, \dots$). Поэтому $\|L\|$ замеряет выразительную силу языка: если символов в языке больше, чем $\omega$, то количество формул совпадает с $|L|$, если меньше — то оно равно $\omega$.

\subsection{Пример произвольной структуры}
Рассмотрим структуру $\mathfrak{A}$ в языке $L = \langle R^3, f^1, c^0 \rangle$:
\begin{itemize}
    \item \textbf{Носитель}: $A = \{0, 1, 2, 3, 4\}$.
    \item \textbf{Константа}: $c^{\mathfrak{A}} = 3 \in A$.
    \item \textbf{Функция}: $f^{\mathfrak{A}}: x \mapsto x + 1 \pmod{5}$.
    \item \textbf{Отношение}: $R^{\mathfrak{A}} \subseteq A^3$, где
    \[ R^{\mathfrak{A}} = \{ (x, y, z) \mid \text{ровно два из } x, y, z \text{ равны} \} \]
\end{itemize}

\subsection{Вложения}

Пусть $\mathfrak{A}$ и $\mathfrak{B}$ — $L$-структуры. 
\textbf{Вложение} $\mathfrak{A}$ в $\mathfrak{B}$ — это инъективная функция $\pi: A \to B$, такая что:

\begin{enumerate}
    \item \textbf{Для предикатов:} Для любого $R^n$ и любых $a_1, \dots, a_n \in A$:
    \[ (a_1, \dots, a_n) \in R^{\mathfrak{A}} \iff (\pi(a_1), \dots, \pi(a_n)) \in R^{\mathfrak{B}} \]
    \item \textbf{Для функций:} Для любого $f^n$ и любых $a_1, \dots, a_n \in A$:
    \[ \pi(f^{\mathfrak{A}}(a_1, \dots, a_n)) = f^{\mathfrak{B}}(\pi(a_1), \dots, \pi(a_n)) \]
    \item \textbf{Для констант:} $\pi(c^{\mathfrak{A}}) = c^{\mathfrak{B}}$ где $c^{\mathfrak{A}} \in A,  c^{\mathfrak{B}} \in B$ \\
    Смысл: в структуре $\mathfrak{B}$ появляется копия структуры $\mathfrak{A}$.
\end{enumerate}

\textbf{Пример вложения (логарифм):}
$\mathfrak{A} = \langle \mathbb{R}^{>0}, \cdot \rangle$, $\mathfrak{B} = \langle \mathbb{R}, + \rangle$.
Вложение $\pi: x \mapsto \ln(x)$. Условие вложения: $\pi(x \cdot y) = \pi(x) + \pi(y)$, что соответствует свойству $\ln(xy) = \ln(x) + \ln(y)$.

\textbf{Пример вложения колец:}
$\mathbb{R}_{ring} \hookrightarrow \mathbb{C}_{ring}$ через $\pi: a \mapsto (a, 0)$.
\begin{itemize}
    \item  $\mathbb{Z}_{ring} \hookrightarrow \mathbb{Q}_{ring} \hookrightarrow \mathbb{R}_{ring} \hookrightarrow \mathbb{C}_{ring}$.
\end{itemize}


\textbf{Изоморфизм:}

$\mathfrak{A}$ и $\mathfrak{B}$ изоморфны, если существует вложение
\[
\pi:A\to B,
\]
которое \textbf{обратимо}, то есть существует вложение
\[
\sigma:B\to A
\]
такое, что для любого $a\in A$ и любого $b\in B$
\[
\sigma(\pi(a))=a,
\qquad
\pi(\sigma(b))=b.
\]

Обозначение:
\[
\mathfrak{A}\cong\mathfrak{B}.
\]

Смысл: структуры отличаются только обозначениями элементов. Все отношения, функции и константы сохраняются.

\textbf{Пример:}

Рассмотрим структуры
\[
\mathfrak{A}=\langle \{1,2,3\},< \rangle,
\qquad
\mathfrak{B}=\langle \{a,b,c\},R \rangle,
\]
где
\[
1<2<3,
\qquad
aRbRc.
\]

Отображение
\[
\pi:1\mapsto a,\quad 2\mapsto b,\quad 3\mapsto c
\]
является изоморфизмом.

Следовательно,
\[
\mathfrak{A}\cong\mathfrak{B}.
\]

\textbf{Пример (автоморфизмы аддитивной группы целых чисел).}

Рассмотрим структуру
\[
\mathbb{Z}_{adgp}=\langle \mathbb{Z},+,-,0\rangle.
\]

Найдём все автоморфизмы
\[
\pi:\mathbb{Z}\to\mathbb{Z}.
\]

Так как $0$ является константой языка,
\[
\pi(0)=0.
\]

Пусть
\[
\pi(1)=n\in\mathbb{Z}.
\]

Тогда для любого натурального числа $m$
\[
m=\underbrace{1+\cdots+1}_{m},
\]
поэтому
\[
\pi(m)
=
\pi(1+\cdots+1)
=
\pi(1)+\cdots+\pi(1)
=
mn.
\]

Для отрицательных чисел
\[
\pi(-m)
=
-\pi(m)
=
-mn.
\]

Следовательно,
\[
\pi(\mathbb{Z})
=
\{0,\pm n,\pm 2n,\ldots\}
=
n\mathbb{Z}.
\]

Так как автоморфизм обратим, отображение $\pi$ сюръективно. Поэтому
\[
\pi(\mathbb{Z})=\mathbb{Z},
\]
откуда
\[
n=\pm1.
\]

Получаем два автоморфизма:
\[
\pi(x)=x,
\qquad
\pi(x)=-x.
\]

\subsection{Термы}

\textbf{Определение:} Термы — это строки символов, которые относятся к константам $L$-структуры или функциям.

\textbf{Примеры для $L_{ring}$:} 
$0, 1+1, -((1+1)+1), (x \cdot (1+1)), x+y, \dots$ \\
В структуре $\mathbb{Q}$ это соответствует значениям: $0, 2, -3, \dots$ (используются функции удвоения, бинарная функция сложения и т.д.).

\textbf{Переменные:} $t, x, y, z, \dots$ 
(Замечание: символы $a, b, c, d, n$ обычно используются как параметры или индексы, но не как переменные терма).

\subsubsection{Индуктивное определение терма языка $L$}
\begin{enumerate}
    \item[I.] Каждая переменная — терм.
    \item[II.] Каждый константный символ — терм.
    \item[III.] Если $f$ — функциональный символ арности $n$, а $t_1, \dots, t_n$ — термы, то $f(t_1, \dots, t_n)$ — терм.
    \item[IV.] Строка является термом только в том случае, если она получена по пунктам I–III.
\end{enumerate}

\textbf{Замкнутый терм:} Терм, не содержащий переменных.

\subsection{Интерпретация термов}

Пусть $\mathfrak{A}$ — $L$-структура, $t$ — $L$-терм. \\
Пусть $\bar{x} = (x_1, \dots, x_r)$ — список переменных, в который входят все переменные, встречающиеся в терме $t$.

\textbf{Интерпретация терма} — это функция $t^{\mathfrak{A}}: A^r \to A$, определяемая индуктивно:

\begin{enumerate}
    \item[I.] Если $t$ является константным символом $c$, то $t^{\mathfrak{A}} = c^{\mathfrak{A}} \in A$.
    \item[II.] Если $t$ является переменной $x_i$ из списка $\bar{x}$, то $t^{\mathfrak{A}}$ — это $i$-тая координатная функция:
    \[ t^{\mathfrak{A}}(x_1, \dots, x_r) = x_i \]
    \item[III.] Если $t$ имеет вид $f(t_1, \dots, t_n)$, то интерпретация $t^{\mathfrak{A}}$ есть композиция:
    \[ t^{\mathfrak{A}}(\bar{x}) = f^{\mathfrak{A}}(t_1^{\mathfrak{A}}(\bar{x}), \dots, t_n^{\mathfrak{A}}(\bar{x})) \]
\end{enumerate}

\textbf{Пример.}

Рассмотрим структуру
\[
\mathfrak{A}=\langle \mathbb{Z}, +^{\mathbb{Z}}, 0^{\mathbb{Z}}, 1^{\mathbb{Z}} \rangle
\]
и терм
\[
t=(x+1)+1.
\]

Тогда
\[
t^{\mathfrak{A}}(3)
=
(3+1)+1
=
5.
\]

Следовательно, интерпретация терма $t$ в структуре $\mathfrak{A}$ является функцией
\[
t^{\mathfrak{A}}:\mathbb{Z}\to\mathbb{Z},
\]
задаваемой правилом
\[
t^{\mathfrak{A}}(x)=x+2.
\]

\textbf{Пример.}

Пусть
\[
t=x+y,
\]
а структура
\[
\mathfrak{A}=\langle \mathbb{Z}, +^{\mathbb{Z}} \rangle.
\]

Тогда интерпретация терма
\[
t^{\mathfrak{A}}:\mathbb{Z}^2\to\mathbb{Z}
\]
имеет вид
\[
t^{\mathfrak{A}}(a,b)=a+b.
\]

Например,
\[
t^{\mathfrak{A}}(2,3)=2+3=5.
\]

% --- Продолжение конспекта: Истинность и Модели ---


% --- Продолжение конспекта: Истинность и Модели ---

\subsection{Истинность формул в структуре}

\textbf{Лемма о термах}

Пусть
\[
\pi:\mathfrak A\to\mathfrak B
\]
— $L$-вложение. Тогда для любого $L$-терма $t$ и любого набора
\[
\bar a=(a_1,\ldots,a_r)\in A^r
\]
выполнено
\[
\pi\bigl(t^{\mathfrak A}(\bar a)\bigr)
=
t^{\mathfrak B}\bigl(\pi(\bar a)\bigr),
\]
где
\[
\pi(\bar a)=(\pi(a_1),\ldots,\pi(a_r)).
\]

\textbf{Доказательство.}
$\square$

Докажем утверждение индукцией по построению терма $t$.

Пусть сначала $t$ является константным символом $c$. Тогда
\[
t^{\mathfrak A}=c^{\mathfrak A}.
\]
Так как $\pi$ является вложением, оно сохраняет константы, поэтому
\[
\pi(c^{\mathfrak A})=c^{\mathfrak B}.
\]
Следовательно,
\[
\pi\bigl(t^{\mathfrak A}(\bar a)\bigr)
=
t^{\mathfrak B}\bigl(\pi(\bar a)\bigr).
\]

Пусть теперь $t$ является переменной $x_i$. Тогда значение этого терма на наборе $\bar a$ равно $a_i$, поэтому
\[
\pi\bigl(t^{\mathfrak A}(\bar a)\bigr)=\pi(a_i).
\]
С другой стороны, на наборе $\pi(\bar a)$ терм $x_i$ выбирает $i$-ю координату, то есть
\[
t^{\mathfrak B}\bigl(\pi(\bar a)\bigr)=\pi(a_i).
\]
Значит, в этом случае утверждение также выполнено.

Осталось рассмотреть случай, когда терм имеет вид
\[
t=f(t_1,\ldots,t_s).
\]
По предположению индукции для термов $t_1,\ldots,t_s$ уже выполнено
\[
\pi\bigl(t_i^{\mathfrak A}(\bar a)\bigr)
=
t_i^{\mathfrak B}\bigl(\pi(\bar a)\bigr)
\]
для всех $i=1,\ldots,s$.

Тогда, используя сохранение функции $f$ вложением $\pi$, получаем
\[
\pi\bigl(t^{\mathfrak A}(\bar a)\bigr)
=
\pi\bigl(f^{\mathfrak A}(t_1^{\mathfrak A}(\bar a),\ldots,t_s^{\mathfrak A}(\bar a))\bigr).
\]
По сохранению функционального символа $f$ это равно
\[
f^{\mathfrak B}\bigl(
\pi(t_1^{\mathfrak A}(\bar a)),\ldots,
\pi(t_s^{\mathfrak A}(\bar a))
\bigr).
\]
По предположению индукции получаем
\[
f^{\mathfrak B}\bigl(
t_1^{\mathfrak B}(\pi(\bar a)),\ldots,
t_s^{\mathfrak B}(\pi(\bar a))
\bigr).
\]
А это ровно
\[
t^{\mathfrak B}\bigl(\pi(\bar a)\bigr).
\]
Следовательно, утверждение доказано для всех $L$-термов $t$.
$\blacksquare$ \vspace{5mm}
 \\
Пусть $\varphi$ — $L$-формула, а $\mathfrak{A}$ — $L$-структура. 
Для формулы со свободными переменными $\bar{x} = (x_1, \dots, x_n)$ мы используем кортеж параметров $\bar{a} = (a_1, \dots, a_n)$, где каждый $a_i \in A$ (подставляем вместо х).

\textbf{Определение:} Запись $\mathfrak{A} \models \varphi[\bar{a}]$ означает, что формула $\varphi$ \textbf{истинна} в структуре $\mathfrak{A}$ при данной интерпретации переменных (или что $\mathfrak{A}$ является \textbf{моделью} для $\varphi$).

Определение строится индуктивно по структуре формулы:

\subsubsection{Условия истинности}

\begin{enumerate}
    \item[\textbf{I.}] \textbf{Равенство термов:}
    \[ \mathfrak{A} \models (t_1 = t_2)[\bar{a}] \iff t_1^{\mathfrak{A}}(\bar{a}) = t_2^{\mathfrak{A}}(\bar{a}) \]
    Утверждение о равенстве истинно, если после вычисления (интерпретации) термов в структуре мы получили один и тот же элемент носителя.
    \textit{Пример:} $1 + 1 = 2$ в $\mathbb{Q}$.

    \item[\textbf{II.}] \textbf{Атомарная формула (отношение):}
    \[ \mathfrak{A} \models R(t_1, \dots, t_s)[\bar{a}] \iff (t_1^{\mathfrak{A}}(\bar{a}), \dots, t_s^{\mathfrak{A}}(\bar{a})) \in R^{\mathfrak{A}} \]
    Отношение истинно, если кортеж значений термов попадает в подмножество $R^{\mathfrak{A}}$ (интерпретацию предиката).
    \textit{Пример:} $x < y$ истинно в $\mathbb{Z}_{<}$ при подстановке $(2, 5)$.

    \item[\textbf{III.}] \textbf{Отрицание ($\neg$):}
    \[ \mathfrak{A} \models \neg \varphi[\bar{a}] \iff \mathfrak{A} \not\models \varphi[\bar{a}] \]
    Формула $\neg \varphi$ истинна тогда и только тогда, когда $\varphi$ ложна в этой структуре.

    \item[\textbf{IV.}] \textbf{Конъюнкция ($\wedge$):}
    \[ \mathfrak{A} \models (\varphi_1 \wedge \varphi_2)[\bar{a}] \iff \mathfrak{A} \models \varphi_1[\bar{a}] \text{ и } \mathfrak{A} \models \varphi_2[\bar{a}] \]
    Сложное утверждение истинно, если истинны обе его простые части одновременно.

    \item[\textbf{V.}] \textbf{Квантор существования ($\exists$):}
    \[ \mathfrak{A} \models \exists x \, \varphi(x, \bar{a}) \iff \text{существует } b \in A, \text{ такой что } \mathfrak{A} \models \varphi(b, \bar{a}) \]
    Утверждение с квантором существования истинно, если в носителе $A$ найдется хотя бы один элемент $b$, при подстановке которого формула станет верной.
\end{enumerate}

\begin{itemize}
    \item \textbf{Свободные переменные:} Формула $\varphi(\bar{x})$ превращается в конкретное высказывание только после подстановки элементов $\bar{a}$ из носителя.
    \item \textbf{Замкнутые формулы (предложения):} Если в $\varphi$ нет свободных переменных, то её истинность зависит только от структуры $\mathfrak{A}$, но не от выбора параметров $\bar{a}$.
\end{itemize}
\par\medskip
\textbf{Предложения и теории.}

Пусть $\psi$ — $L$-предложение, то есть $L$-формула без свободных переменных.

Если $\mathfrak A$ — $L$-структура, то запись
\[
\mathfrak A\models \psi
\]
означает, что предложение $\psi$ истинно в структуре $\mathfrak A$.

Если $\Sigma$ — множество $L$-предложений, то запись
\[
\mathfrak A\models \Sigma
\]
означает, что структура $\mathfrak A$ удовлетворяет каждому предложению из $\Sigma$.

Иначе говоря,
\[
\mathfrak A\models \Sigma
\qquad
\Longleftrightarrow
\qquad
\mathfrak A\models \psi
\text{ для любого } \psi\in\Sigma.
\]

В этом случае говорят, что $\mathfrak A$ является моделью множества предложений $\Sigma$.

\textbf{Пример.}

Пусть $\Sigma_1$ — множество аксиом групп. Тогда
\[
\mathfrak A\models \Sigma_1
\]
означает, что структура $\mathfrak A$ является группой.

Пусть $\Sigma_2$ — множество аксиом абелевых групп. Тогда
\[
\mathfrak A\models \Sigma_2
\]
означает, что структура $\mathfrak A$ является абелевой группой.

Например,
\[
S_3\not\models \Sigma_2,
\]
поскольку группа $S_3$ не является абелевой.
\par\medskip
\textbf{Лемма.}

Пусть
\[
\pi:\mathfrak A\to\mathfrak B
\]
— $L$-вложение. Тогда для любой атомарной $L$-формулы $\varphi(\bar x)$ и любого набора
\[
\bar a\in A^n
\]
выполнено
\[
\mathfrak A\models \varphi(\bar a)
\qquad
\Longleftrightarrow
\qquad
\mathfrak B\models \varphi(\pi(\bar a))
\]

\textbf{Доказательство.} $\square$

Достаточно рассмотреть два вида атомарных формул.

Пусть сначала
\[
\varphi(\bar x)
\]
имеет вид
\[
R(t_1(\bar x),\ldots,t_r(\bar x)).
\]

Обозначим
\[
\alpha_i=t_i^{\mathfrak A}(\bar a),
\qquad
\beta_i=t_i^{\mathfrak B}(\pi(\bar a)).
\]

По лемме о термах
\[
\pi(\alpha_i)=\beta_i
\]
для всех $i=1,\ldots,r$. Поэтому
\[
\pi(\bar \alpha)=\bar \beta,
\]
где
\[
\bar \alpha=(\alpha_1,\ldots,\alpha_r),
\qquad
\bar \beta=(\beta_1,\ldots,\beta_r).
\]

По определению истинности атомарной формулы
\[
\mathfrak A\models R(t_1(\bar a),\ldots,t_r(\bar a))
\]
тогда и только тогда, когда
\[
\bar \alpha\in R^{\mathfrak A}.
\]

Так как $\pi$ является вложением, оно сохраняет отношения в обе стороны. Поэтому
\[
\bar \alpha\in R^{\mathfrak A}
\qquad
\Longleftrightarrow
\qquad
\bar \beta\in R^{\mathfrak B}.
\]

А последнее равносильно тому, что
\[
\mathfrak B\models R(t_1(\pi(\bar a)),\ldots,t_r(\pi(\bar a))).
\]

Следовательно,
\[
\mathfrak A\models \varphi(\bar a)
\qquad
\Longleftrightarrow
\qquad
\mathfrak B\models \varphi(\pi(\bar a)).
\]

Пусть теперь $\varphi(\bar x)$ имеет вид
\[
t_1(\bar x)=t_2(\bar x).
\]

По определению истинности равенства
\[
\mathfrak A\models (t_1=t_2)(\bar a)
\]
тогда и только тогда, когда
\[
t_1^{\mathfrak A}(\bar a)=t_2^{\mathfrak A}(\bar a).
\]

Применяя $\pi$ и используя инъективность $\pi$, получаем равносильное условие
\[
\pi(t_1^{\mathfrak A}(\bar a))
=
\pi(t_2^{\mathfrak A}(\bar a)).
\]

По лемме о термах это равносильно
\[
t_1^{\mathfrak B}(\pi(\bar a))
=
t_2^{\mathfrak B}(\pi(\bar a)).
\]

То есть
\[
\mathfrak B\models (t_1=t_2)(\pi(\bar a)).
\]

Следовательно, утверждение верно и для атомарных формул вида равенства.

Значит, для любой атомарной $L$-формулы $\varphi(\bar x)$ выполнено
\[
\mathfrak A\models \varphi(\bar a)
\qquad
\Longleftrightarrow
\qquad
\mathfrak B\models \varphi(\pi(\bar a)).
\]

$\blacksquare$

\begin{lemma}
Пусть
\[
\pi:\mathfrak A\to\mathfrak B
\]
— $L$-изоморфизм. Тогда для любой $L$-формулы $\varphi(\bar x)$ и любого набора
\[
\bar a=(a_1,\ldots,a_n)\in A^n
\]
выполнено
\[
\mathfrak A\models \varphi(\bar a)
\qquad
\Longleftrightarrow
\qquad
\mathfrak B\models \varphi(\pi(\bar a)),
\]
где
\[
\pi(\bar a)=(\pi(a_1),\ldots,\pi(a_n)).
\]
\end{lemma}

\begin{proof}
Докажем утверждение структурной индукцией по построению формулы $\varphi$.

База индукции — атомарные формулы. Если $\varphi$ атомарна, то утверждение уже доказано в предыдущей лемме:
\[
\mathfrak A\models \varphi(\bar a)
\qquad
\Longleftrightarrow
\qquad
\mathfrak B\models \varphi(\pi(\bar a)).
\]

Теперь докажем индукционный переход. Индукционное предположение состоит в том, что утверждение уже верно для непосредственных подформул рассматриваемой формулы.

Пусть сначала
\[
\varphi=\neg\psi.
\]
Тогда по определению истинности отрицания
\[
\mathfrak A\models \neg\psi(\bar a)
\qquad
\Longleftrightarrow
\qquad
\mathfrak A\not\models \psi(\bar a).
\]
По индукционному предположению это равносильно
\[
\mathfrak B\not\models \psi(\pi(\bar a)).
\]
Снова используя определение истинности отрицания, получаем
\[
\mathfrak B\models \neg\psi(\pi(\bar a)).
\]
Следовательно,
\[
\mathfrak A\models \neg\psi(\bar a)
\qquad
\Longleftrightarrow
\qquad
\mathfrak B\models \neg\psi(\pi(\bar a)).
\]

Пусть теперь
\[
\varphi=\psi\wedge\theta.
\]
Тогда по определению истинности конъюнкции
\[
\mathfrak A\models(\psi\wedge\theta)(\bar a)
\]
тогда и только тогда, когда
\[
\mathfrak A\models\psi(\bar a)
\qquad\text{и}\qquad
\mathfrak A\models\theta(\bar a).
\]
По индукционному предположению это равносильно тому, что
\[
\mathfrak B\models\psi(\pi(\bar a))
\qquad\text{и}\qquad
\mathfrak B\models\theta(\pi(\bar a)).
\]
Следовательно, по определению истинности конъюнкции
\[
\mathfrak B\models(\psi\wedge\theta)(\pi(\bar a)).
\]
Значит,
\[
\mathfrak A\models(\psi\wedge\theta)(\bar a)
\qquad
\Longleftrightarrow
\qquad
\mathfrak B\models(\psi\wedge\theta)(\pi(\bar a)).
\]

Аналогично рассматриваются связки $\vee$ и $\to$, если они используются как основные логические связки языка.

Осталось рассмотреть квантор существования. Пусть
\[
\varphi(\bar y)=\exists x\,\psi(x,\bar y).
\]

Докажем сначала импликацию слева направо. Пусть
\[
\mathfrak A\models \exists x\,\psi(x,\bar a).
\]
По определению истинности квантора существования существует элемент $c\in A$, такой что
\[
\mathfrak A\models \psi(c,\bar a).
\]
По индукционному предположению, применённому к формуле $\psi(x,\bar y)$ и набору $(c,\bar a)$, получаем
\[
\mathfrak B\models \psi(\pi(c),\pi(\bar a)).
\]
Так как $\pi(c)\in B$, то в структуре $\mathfrak B$ существует свидетель для квантора. Следовательно,
\[
\mathfrak B\models \exists x\,\psi(x,\pi(\bar a)).
\]

Докажем теперь обратную импликацию. Пусть
\[
\mathfrak B\models \exists x\,\psi(x,\pi(\bar a)).
\]
Тогда существует элемент $d\in B$, такой что
\[
\mathfrak B\models \psi(d,\pi(\bar a)).
\]
Поскольку $\pi$ является изоморфизмом, оно сюръективно. Поэтому существует элемент $c\in A$, такой что
\[
d=\pi(c).
\]
Тогда
\[
\mathfrak B\models \psi(\pi(c),\pi(\bar a)).
\]
По индукционному предположению получаем
\[
\mathfrak A\models \psi(c,\bar a).
\]
Следовательно,
\[
\mathfrak A\models \exists x\,\psi(x,\bar a).
\]

Итак,
\[
\mathfrak A\models \exists x\,\psi(x,\bar a)
\qquad
\Longleftrightarrow
\qquad
\mathfrak B\models \exists x\,\psi(x,\pi(\bar a)).
\]

Квантор всеобщности можно рассмотреть аналогично либо выразить через отрицание и квантор существования:
\[
\forall x\,\psi \equiv \neg\exists x\,\neg\psi.
\]

Следовательно, утверждение верно для всех $L$-формул $\varphi$.
\end{proof}

\begin{remark}
Смысл леммы состоит в том, что изоморфные $L$-структуры неразличимы средствами языка $L$.
Если структуры отличаются только переименованием элементов, то любая $L$-формула имеет на соответствующих наборах одни и те же истинностные значения.
\end{remark}

\begin{remark}
Условие, что $\pi$ является именно изоморфизмом, существенно.
Для произвольного вложения утверждение для всех формул может быть неверно, потому что кванторы могут увидеть новые элементы в большей структуре.
Именно поэтому в доказательстве для квантора существования используется сюръективность отображения $\pi$.
\end{remark}

\subsection{ Синтаксические сокращения}

Во многих случаях удобно считать, что некоторые логические связки и обозначения являются сокращениями через более базовые символы.

Например, неравенство термов определяется как сокращение:
\[
t_1\ne t_2
\quad:\Longleftrightarrow\quad
\neg(t_1=t_2).
\]

Дизъюнкцию можно выразить через отрицание и конъюнкцию:
\[
\varphi\vee\psi
\quad:\Longleftrightarrow\quad
\neg(\neg\varphi\wedge\neg\psi).
\]

Импликацию можно выразить через отрицание и дизъюнкцию:
\[
\varphi\to\psi
\quad:\Longleftrightarrow\quad
\neg\varphi\vee\psi.
\]

Эквивалентность формул определяется как взаимная импликация:
\[
\varphi\leftrightarrow\psi
\quad:\Longleftrightarrow\quad
(\varphi\to\psi)\wedge(\psi\to\varphi).
\]

Квантор всеобщности можно выразить через отрицание и квантор существования:
\[
\forall x\,\varphi
\quad:\Longleftrightarrow\quad
\neg\exists x\,\neg\varphi.
\]

Также используют сокращение для нескольких кванторов подряд:
\[
\forall x_1\ldots x_n\,\varphi
\quad:\Longleftrightarrow\quad
\forall x_1\cdots\forall x_n\,\varphi.
\]

Аналогично,
\[
\exists x_1\ldots x_n\,\varphi
\quad:\Longleftrightarrow\quad
\exists x_1\cdots\exists x_n\,\varphi.
\]

\begin{remark}
Смысл этих сокращений состоит в том, что формально язык можно задавать с небольшим числом логических символов, например
\[
\neg,\qquad \wedge,\qquad \exists,
\]
а остальные связки и кванторы вводить как удобные обозначения.
\end{remark}


\section{Определимые множества}

Пусть $\mathfrak A$ — $L$-структура, а $S\subseteq A^n$.

\begin{definition}
Множество $S$ называется \textbf{определимым} в структуре $\mathfrak A$, если существует $L$-формула
\[
\varphi(x_1,\ldots,x_n)
\]
такая, что
\[
S
=
\{\bar a\in A^n\mid \mathfrak A\models \varphi(\bar a)\}.
\]
В этом случае также пишут
\[
S=\varphi(\mathfrak A).
\]
\end{definition}

Иными словами, формула $\varphi(x_1,\ldots,x_n)$ задаёт множество всех наборов элементов структуры $\mathfrak A$, на которых эта формула истинна.

Если формула имеет одну свободную переменную, то она задаёт подмножество носителя:
\[
\varphi(\mathfrak A)\subseteq A.
\]
Если формула имеет две свободные переменные, то она задаёт подмножество декартова квадрата:
\[
\varphi(\mathfrak A)\subseteq A^2.
\]

\begin{example}
Рассмотрим структуру
\[
\mathfrak A=\langle \{1,2,3,4,5\},<\rangle.
\]
Пусть
\[
\varphi(x)\equiv x<4.
\]
Тогда
\[
\varphi(\mathfrak A)
=
\{a\in A\mid \mathfrak A\models a<4\}.
\]
Поскольку
\[
1<4,\qquad 2<4,\qquad 3<4,
\]
но
\[
4\not<4,\qquad 5\not<4,
\]
получаем
\[
\varphi(\mathfrak A)=\{1,2,3\}.
\]
Таким образом, формула $x<4$ определяет множество $\{1,2,3\}$ в структуре $\mathfrak A$.
\end{example}

\begin{example}
Рассмотрим структуру упорядоченного кольца вещественных чисел
\[
\mathbb R_{o\text{-}ring}
=
\langle \mathbb R,+,\cdot,-,0,1,<\rangle.
\]
Множество
\[
(0,2)\subseteq\mathbb R
\]
определяется формулой
\[
(0<x)\wedge(x<1+1).
\]
Действительно, терм $1+1$ интерпретируется как число $2$, поэтому эта формула задаёт ровно условие
\[
0<x<2.
\]
\end{example}

\begin{example}
В той же структуре
\[
\mathbb R_{o\text{-}ring}
=
\langle \mathbb R,+,\cdot,-,0,1,<\rangle
\]
множество
\[
[0,2)\subseteq\mathbb R
\]
определяется формулой
\[
(0=x\vee 0<x)\wedge(x<1+1).
\]
Здесь выражение
\[
0=x\vee 0<x
\]
задаёт условие $x\ge 0$, а выражение
\[
x<1+1
\]
задаёт условие $x<2$.
\end{example}

\begin{example}
Множество
\[
\{-1,3\}\subseteq\mathbb R
\]
определяется формулой
\[
(x+1=0)\vee(x=1+1+1).
\]
Действительно,
\[
x+1=0
\]
задаёт условие $x=-1$, а
\[
x=1+1+1
\]
задаёт условие $x=3$.
Следовательно, множество решений этой формулы равно $\{-1,3\}$.
\end{example}

\begin{example}
В структуре
\[
\mathbb R_{o\text{-}ring}
=
\langle \mathbb R,+,\cdot,-,0,1,<\rangle
\]
множество
\[
[0,2)\times(-2,0)\subseteq\mathbb R^2
\]
определяется формулой
\[
(0=x\vee 0<x)\wedge(x<1+1)\wedge(y<0)\wedge(-(1+1)<y).
\]
Первые две части формулы задают условие
\[
0\le x<2.
\]
Последние две части задают условие
\[
-2<y<0.
\]
Следовательно, вся формула определяет множество
\[
[0,2)\times(-2,0).
\]
\end{example}

\begin{remark}
Определимые множества — это множества, которые можно описать средствами языка $L$ внутри фиксированной структуры.
Если нужного символа нет в языке, то соответствующее свойство может перестать быть выразимым.
Например, порядок $<$ нельзя использовать в языке колец
\[
L_{ring}=\langle +,\cdot,-,0,1\rangle,
\]
если он не добавлен в сигнатуру.
\end{remark}

\begin{remark}
Из леммы о сохранении формул изоморфизмами следует, что изоморфизмы переводят определимые множества в определимые множества.

Если
\[
\pi:\mathfrak A\to\mathfrak B
\]
— $L$-изоморфизм и
\[
S=\varphi(\mathfrak A),
\]
то
\[
\pi(S)=\varphi(\mathfrak B).
\]
Именно поэтому определимые множества зависят не от имён элементов, а от самой структуры.
\end{remark}


\subsection{Пример: конечный порядок элемента группы}

Рассмотрим язык аддитивных групп
\[
L_{adgr}=\langle +,-,0\rangle
\]
и абелеву группу
\[
G_{ab}.
\]

Пусть $g\in G_{ab}$. В обычной алгебре говорят, что $g$ имеет конечный порядок, если существует натуральное число $n>0$ такое, что
\[
\underbrace{g+\cdots+g}_{n\text{ раз}}=0.
\]

Для каждого фиксированного $n$ свойство
\[
\underbrace{g+\cdots+g}_{n\text{ раз}}=0
\]
задаётся одной $L_{adgr}$-формулой.

Например, свойство порядка, делящего $2$, задаётся формулой
\[
g+g=0.
\]
Свойство порядка, делящего $3$, задаётся формулой
\[
g+g+g=0.
\]

Однако свойство ``$g$ имеет конечный порядок'' выглядело бы как бесконечная дизъюнкция:
\[
(g=0)\vee(g+g=0)\vee(g+g+g=0)\vee\ldots
\]
Но такая запись не является формулой первого порядка, потому что формула должна быть конечной строкой символов.

\begin{remark}
Свойство ``быть элементом конечного порядка'' не задаётся одной формулой первого порядка в языке аддитивных групп.
При этом каждое отдельное свойство вида $ng=0$ при фиксированном $n$ формулой первого порядка задаётся.
\end{remark}

\subsection{Операции над определимыми множествами}

Пусть $\mathfrak A$ — $L$-структура.

\begin{proposition}
Пусть
\[
S_1,S_2\subseteq A^n
\]
— определимые множества. Тогда множества
\[
S_1\cap S_2,\qquad S_1\cup S_2,\qquad A^n\setminus S_1
\]
также определимы.
\end{proposition}

\begin{proof}
Пусть множества $S_1$ и $S_2$ заданы формулами $\varphi(\bar x)$ и $\psi(\bar x)$ соответственно, то есть
\[
S_1=\{\bar a\in A^n\mid \mathfrak A\models\varphi(\bar a)\},
\]
\[
S_2=\{\bar a\in A^n\mid \mathfrak A\models\psi(\bar a)\}.
\]

Тогда пересечение $S_1\cap S_2$ задаётся формулой
\[
\varphi(\bar x)\wedge\psi(\bar x).
\]
Действительно,
\[
\bar a\in S_1\cap S_2
\]
тогда и только тогда, когда
\[
\mathfrak A\models\varphi(\bar a)
\qquad\text{и}\qquad
\mathfrak A\models\psi(\bar a),
\]
что равносильно
\[
\mathfrak A\models(\varphi\wedge\psi)(\bar a).
\]

Объединение $S_1\cup S_2$ задаётся формулой
\[
\varphi(\bar x)\vee\psi(\bar x),
\]
а дополнение $A^n\setminus S_1$ задаётся формулой
\[
\neg\varphi(\bar x).
\]

Следовательно, все три множества определимы.
\end{proof}

\begin{remark}
Определимые множества замкнуты относительно булевых операций. Это соответствует тому, что формулы можно соединять логическими связками:
\[
\cap \longleftrightarrow \wedge,
\qquad
\cup \longleftrightarrow \vee,
\qquad
\text{дополнение} \longleftrightarrow \neg.
\]
\end{remark}

\begin{proposition}
Пусть
\[
S\subseteq A^{n+m}
\]
— определимое множество. Тогда его проекция на первые $n$ координат
\[
\operatorname{pr}(S)
=
\{\bar a\in A^n\mid \text{существует } \bar b\in A^m,\;(\bar a,\bar b)\in S\}
\]
также определима.
\end{proposition}

\begin{proof}
Пусть множество $S$ задано формулой
\[
\psi(\bar x,\bar y),
\]
где
\[
\bar x=(x_1,\ldots,x_n),
\qquad
\bar y=(y_1,\ldots,y_m).
\]
То есть
\[
S=
\{(\bar a,\bar b)\in A^{n+m}\mid
\mathfrak A\models\psi(\bar a,\bar b)\}.
\]

Тогда проекция $\operatorname{pr}(S)$ задаётся формулой
\[
\exists y_1\ldots \exists y_m\,\psi(\bar x,\bar y).
\]
Действительно,
\[
\bar a\in\operatorname{pr}(S)
\]
тогда и только тогда, когда существует $\bar b\in A^m$, такой что
\[
(\bar a,\bar b)\in S.
\]
По определению множества $S$ это равносильно тому, что существует $\bar b\in A^m$, такой что
\[
\mathfrak A\models\psi(\bar a,\bar b).
\]
А это ровно означает
\[
\mathfrak A\models\exists y_1\ldots \exists y_m\,\psi(\bar a,\bar y).
\]

Следовательно, проекция определима.
\end{proof}

\begin{remark}
Квантор существования имеет геометрический смысл проекции. Если множество точек $(x,y)$ задано формулой $\psi(x,y)$, то формула
\[
\exists y\,\psi(x,y)
\]
задаёт множество всех $x$, для которых найдётся подходящий $y$.
\end{remark}

\subsection{Определимые функции}

Пусть
\[
D\subseteq A^n,
\qquad
C\subseteq A^m,
\]
и пусть
\[
f:D\to C
\]
— функция.

\begin{definition}
Функция $f:D\to C$ называется \textbf{определимой} в структуре $\mathfrak A$, если её график
\[
\Gamma_f=
\{(\bar a,f(\bar a))\mid \bar a\in D\}
\subseteq A^{n+m}
\]
является определимым множеством.
\end{definition}

Иными словами, функция $f$ определима, если существует $L$-формула
\[
\psi(\bar x,\bar y)
\]
такая, что
\[
\Gamma_f
=
\{(\bar a,\bar b)\in A^{n+m}\mid
\mathfrak A\models\psi(\bar a,\bar b)\}.
\]

\begin{remark}
Функция $f$ не обязана быть символом языка $L$.

Если $f\in L$ — функциональный символ языка, то в каждой $L$-структуре $\mathfrak A$ у него уже есть интерпретация
\[
f^{\mathfrak A}:A^n\to A.
\]
Например, в языке колец
\[
L_{ring}=\langle +,\cdot,-,0,1\rangle
\]
символ $+$ принадлежит языку, поэтому в структуре $\mathbb R_{ring}$ он интерпретируется как обычное сложение.

Но определимая функция может не входить в язык как символ. Достаточно, чтобы её график задавался некоторой $L$-формулой.
\end{remark}

\begin{remark}
Если график функции задан формулой $\psi(\bar x,\bar y)$, то область определения и образ этой функции также задаются с помощью квантора существования.

Область определения задаётся формулой
\[
\exists \bar y\,\psi(\bar x,\bar y),
\]
а образ задаётся формулой
\[
\exists \bar x\,\psi(\bar x,\bar y).
\]
\end{remark}

\begin{example}
Рассмотрим структуру упорядоченного кольца вещественных чисел
\[
\mathbb R_{o\text{-}ring}
=
\langle \mathbb R,+,\cdot,-,0,1,<\rangle.
\]
Функция
\[
f:\mathbb R\to\mathbb R,
\qquad
f(x)=|x|
\]
определима в этой структуре, хотя символа $|\cdot|$ нет в языке.

Её график задаётся формулой
\[
\bigl((x<0)\wedge(y=-x)\bigr)
\vee
\bigl((x=0\vee 0<x)\wedge(y=x)\bigr).
\]
Действительно, если $x<0$, то формула требует $y=-x$, а если $x\ge 0$, то формула требует $y=x$.
Следовательно, множество решений этой формулы есть график функции $x\mapsto |x|$.
\end{example}

\begin{example}
Рассмотрим формулу
\[
\psi(x,y):\quad (x-1)^2+(y-3)^2\le 4
\]
в структуре
\[
\mathbb R_{o\text{-}ring}
=
\langle \mathbb R,+,\cdot,-,0,1,<\rangle.
\]
Эта формула задаёт диск на плоскости с центром в точке $(1,3)$ и радиусом $2$.

Формула
\[
\exists y\,\psi(x,y)
\]
задаёт проекцию этого диска на ось $x$. Поскольку центр диска имеет $x$-координату $1$, а радиус равен $2$, получаем множество
\[
[-1,3].
\]

Формула
\[
\exists x\,\psi(x,y)
\]
задаёт проекцию этого диска на ось $y$. Поскольку центр диска имеет $y$-координату $3$, а радиус равен $2$, получаем множество
\[
[1,5].
\]
\end{example}

\begin{remark}
Таким образом, в языке теории моделей выполняются следующие соответствия:
\[
\text{булевы операции над множествами}
\quad\longleftrightarrow\quad
\neg,\wedge,\vee,
\]
\[
\text{проекции}
\quad\longleftrightarrow\quad
\exists,
\]
\[
\text{определимые функции}
\quad\longleftrightarrow\quad
\text{определимые графики}.
\]
\end{remark}

\subsection{Автоморфизмы, параметры и определимые множества}

Пусть $\mathfrak A$ — $L$-структура, и пусть
\[
S\subseteq A^n.
\]

Прежде чем связывать определимые множества с автоморфизмами, уточним смысл слова ``определимое''.

\begin{definition}
Множество $S\subseteq A^n$ называется \textbf{определимым без параметров} в структуре $\mathfrak A$, если существует $L$-формула
\[
\varphi(x_1,\ldots,x_n)
\]
такая, что
\[
S=
\{\bar a\in A^n\mid \mathfrak A\models\varphi(\bar a)\}.
\]
\end{definition}

Здесь формула $\varphi$ может использовать только символы языка $L$ и логические символы первого порядка. Нельзя просто взять конкретный элемент структуры и написать его имя, если этот элемент не является константным символом языка.

Например, если
\[
\mathfrak A=\langle\mathbb Q,+,<\rangle,
\]
то символов $0$ и $1$ в языке формально нет. Поэтому запись
\[
x=1
\]
не является формулой языка $\langle +,<\rangle$, если $1$ не разрешён как параметр.

\begin{definition}
Пусть $B\subseteq A$. Множество $S\subseteq A^n$ называется \textbf{определимым с параметрами из $B$}, если существуют элементы
\[
\bar b=(b_1,\ldots,b_m)\in B^m
\]
и $L$-формула
\[
\varphi(\bar x,\bar y)
\]
такие, что
\[
S=
\{\bar a\in A^n\mid \mathfrak A\models\varphi(\bar a,\bar b)\}.
\]
\end{definition}

Иными словами, параметры — это фиксированные элементы структуры, которые мы разрешаем использовать в формуле как дополнительные имена.

Определимость без параметров — это частный случай определимости с параметрами, когда никаких дополнительных элементов структуры использовать нельзя.

\begin{remark}
Множество $S$ само по себе не является формулой и не является ``свойством'' в строгом синтаксическом смысле. Точнее, $S$ — это множество наборов, на которых некоторое условие может быть истинным.

Если формула $\varphi(\bar x)$ задаёт множество $S$, то можно говорить, что элементы $\bar a\in S$ обладают свойством, выраженным формулой $\varphi$. Формально это записывается так:
\[
S=\varphi(\mathfrak A)
=
\{\bar a\in A^n\mid \mathfrak A\models\varphi(\bar a)\}.
\]
\end{remark}

\begin{remark}
Нужно различать символы языка и логические символы.

Например, в языке упорядоченных колец
\[
L_{o\text{-}ring}
=
\langle +,\cdot,-,0,1,<\rangle
\]
перечислены только нелогические символы структуры:
\[
+,\quad \cdot,\quad -,\quad 0,\quad 1,\quad <.
\]

Но при построении формул первого порядка мы также можем использовать логические символы:
\[
=,\quad \neg,\quad \wedge,\quad \vee,\quad \to,\quad \exists,\quad \forall.
\]
Они не перечисляются в сигнатуре $L$, потому что принадлежат самой логике первого порядка, а не конкретной структуре.
\end{remark}

\begin{proposition}
Если множество $S$ определимо в структуре $\mathfrak A$ без параметров, то для любого автоморфизма
\[
\pi\in Aut(\mathfrak A)
\]
выполнено
\[
\pi(S)=S.
\]
\end{proposition}

\begin{proof}
Пусть множество $S$ определимо формулой $\varphi(\bar x)$, то есть
\[
S=
\{\bar a\in A^n\mid \mathfrak A\models\varphi(\bar a)\}.
\]
Пусть
\[
\pi\in Aut(\mathfrak A).
\]
Тогда $\pi:\mathfrak A\to\mathfrak A$ является $L$-изоморфизмом.

По лемме о сохранении формул изоморфизмами для любого набора $\bar a\in A^n$ имеем
\[
\mathfrak A\models\varphi(\bar a)
\qquad
\Longleftrightarrow
\qquad
\mathfrak A\models\varphi(\pi(\bar a)).
\]
Следовательно,
\[
\bar a\in S
\qquad
\Longleftrightarrow
\qquad
\pi(\bar a)\in S.
\]
Значит,
\[
\pi(S)=S.
\]
\end{proof}

\begin{remark}
Автоморфизм можно понимать как внутреннюю симметрию структуры. Он переименовывает элементы носителя, но сохраняет всё, что задано языком $L$.

Поэтому формула языка $L$ не может различать наборы, которые переводятся друг в друга автоморфизмом. Если множество задано формулой, оно обязано быть устойчивым относительно всех таких симметрий.
\end{remark}

\begin{remark}
Это утверждение даёт удобный способ доказывать неопределимость без параметров.

Если существует автоморфизм
\[
\pi\in Aut(\mathfrak A)
\]
такой, что
\[
\pi(S)\ne S,
\]
то множество $S$ не может быть определимым без параметров в структуре $\mathfrak A$.
\end{remark}

\begin{remark}
Важно, что это только необходимое условие определимости.

Если множество определимо, то оно инвариантно относительно всех автоморфизмов. Но обратное утверждение в общем случае неверно:
\[
\forall \pi\in Aut(\mathfrak A)\quad \pi(S)=S
\qquad
\not\Longrightarrow
\qquad
S\text{ определимо}.
\]
\end{remark}

\begin{proposition}
Пусть $B\subseteq A$. Если множество $S\subseteq A^n$ определимо с параметрами из $B$, то оно сохраняется всеми автоморфизмами структуры $\mathfrak A$, которые фиксируют каждый элемент из $B$.
\end{proposition}

\begin{proof}
Пусть
\[
S=
\{\bar a\in A^n\mid \mathfrak A\models\varphi(\bar a,\bar b)\},
\]
где
\[
\bar b=(b_1,\ldots,b_m)\in B^m.
\]
Пусть $\pi\in Aut(\mathfrak A)$ и пусть $\pi$ фиксирует все элементы из $B$, то есть
\[
\pi(b)=b
\]
для любого $b\in B$.

Тогда, в частности,
\[
\pi(\bar b)=\bar b.
\]
По лемме о сохранении формул изоморфизмами
\[
\mathfrak A\models\varphi(\bar a,\bar b)
\qquad
\Longleftrightarrow
\qquad
\mathfrak A\models\varphi(\pi(\bar a),\pi(\bar b)).
\]
Так как $\pi(\bar b)=\bar b$, получаем
\[
\mathfrak A\models\varphi(\bar a,\bar b)
\qquad
\Longleftrightarrow
\qquad
\mathfrak A\models\varphi(\pi(\bar a),\bar b).
\]
Следовательно,
\[
\bar a\in S
\qquad
\Longleftrightarrow
\qquad
\pi(\bar a)\in S.
\]
Значит,
\[
\pi(S)=S.
\]
\end{proof}

\begin{remark}
Для определимости с параметрами автоморфизмный тест становится слабее. Теперь нужно рассматривать не все автоморфизмы структуры, а только те, которые фиксируют выбранные параметры.

Если множество $S$ определимо с параметрами из $B$, то его не обязаны сохранять все автоморфизмы из $Aut(\mathfrak A)$; его обязаны сохранять только автоморфизмы, фиксирующие $B$ поточечно.
\end{remark}

\begin{example}
Рассмотрим структуру
\[
\langle\mathbb Z,<\rangle.
\]
Отображение
\[
\pi:\mathbb Z\to\mathbb Z,
\qquad
\pi(x)=x+1
\]
является автоморфизмом этой структуры, поскольку для любых $a,b\in\mathbb Z$
\[
a<b
\qquad
\Longleftrightarrow
\qquad
a+1<b+1.
\]

Покажем, что множество
\[
\{0\}\subseteq\mathbb Z
\]
не определимо в структуре $\langle\mathbb Z,<\rangle$ без параметров.

Действительно,
\[
\pi(\{0\})=\{1\}\ne\{0\}.
\]
Если бы множество $\{0\}$ было определимо без параметров, то оно сохранялось бы любым автоморфизмом. Получили противоречие.

Следовательно, $\{0\}$ не определимо без параметров в структуре $\langle\mathbb Z,<\rangle$.
\end{example}

\begin{remark}
При этом множество $\{0\}$ определимо в структуре $\langle\mathbb Z,<\rangle$ с параметром $0$.

Действительно, если разрешить использовать элемент $0$ как параметр, то $\{0\}$ задаётся формулой
\[
x=0.
\]
Здесь $0$ не является символом языка $\langle <\rangle$, а используется как внешний параметр.
\end{remark}

\begin{example}
Рассмотрим структуру
\[
\mathbb R_{adgp}
=
\langle\mathbb R,+,-,0\rangle
\]
и множество
\[
S=\{(x,y)\in\mathbb R^2\mid x<y\}.
\]
Множество $S$ является бинарным отношением на $\mathbb R$. Его можно понимать как множество всех пар, в которых первая координата меньше второй.

В языке аддитивных групп нет символа порядка $<$.

Отображение
\[
\pi:\mathbb R\to\mathbb R,
\qquad
\pi(x)=-x
\]
является автоморфизмом структуры $\mathbb R_{adgp}$. Действительно,
\[
\pi(x+y)=-(x+y)=(-x)+(-y)=\pi(x)+\pi(y),
\]
и
\[
\pi(0)=0.
\]

Однако множество $S$ не сохраняется этим автоморфизмом. Например,
\[
(0,1)\in S,
\]
поскольку
\[
0<1.
\]
Но
\[
\pi(0)=0,
\qquad
\pi(1)=-1,
\]
и
\[
(0,-1)\notin S,
\]
поскольку неверно, что $0<-1$.

Следовательно,
\[
\pi(S)\ne S.
\]
Значит, множество
\[
S=\{(x,y)\in\mathbb R^2\mid x<y\}
\]
не определимо в структуре
\[
\langle\mathbb R,+,-,0\rangle
\]
без параметров.
\end{example}

\begin{remark}
Содержательно это означает, что порядок $<$ нельзя выразить средствами одной только аддитивной группы вещественных чисел.

Язык
\[
\langle +,-,0\rangle
\]
не видит положительное направление на прямой. Автоморфизм $x\mapsto -x$ является симметрией этой структуры и меняет местами ``правое'' и ``левое'' направление.

Если бы порядок был выразим формулой языка $\langle +,-,0\rangle$, то он должен был бы сохраняться этой симметрией. Но он не сохраняется.
\end{remark}

\begin{example}
Рассмотрим структуру
\[
\langle\mathbb Q,+,<\rangle.
\]
Отображение
\[
\pi:\mathbb Q\to\mathbb Q,
\qquad
\pi(x)=2x
\]
является автоморфизмом этой структуры. Оно сохраняет сложение:
\[
\pi(x+y)=2(x+y)=2x+2y=\pi(x)+\pi(y),
\]
и сохраняет порядок:
\[
x<y
\qquad
\Longleftrightarrow
\qquad
2x<2y.
\]

При этом
\[
\pi(0)=0,
\qquad
\pi(1)=2.
\]

Множество $\{0\}$ определимо в структуре $\langle\mathbb Q,+,<\rangle$ без параметров. Например, оно задаётся формулой
\[
\psi(z):\quad \forall x\,(x+z=x\wedge z+x=x).
\]
Эта формула говорит, что $z$ является нейтральным элементом сложения. В структуре $\langle\mathbb Q,+,<\rangle$ таким элементом является ровно $0$.

Однако множество $\{1\}$ не определимо без параметров. Действительно,
\[
\pi(\{1\})=\{2\}\ne\{1\}.
\]
Если бы $\{1\}$ было определимо без параметров, оно должно было бы сохраняться любым автоморфизмом. Значит, $\{1\}$ не определимо без параметров в структуре $\langle\mathbb Q,+,<\rangle$.
\end{example}

\begin{remark}
Множество $\{1\}$ всё же определимо с параметром $1$.

Если разрешить использовать элемент $1\in\mathbb Q$ как параметр, то $\{1\}$ задаётся формулой
\[
x=1.
\]
Поэтому фраза ``не определимо без параметров'' не означает ``вообще никак не определимо''. Она означает только, что множество нельзя описать, используя одни лишь символы языка.
\end{remark}

\begin{example}
Рассмотрим структуру
\[
\langle\omega,S\rangle,
\]
где
\[
\omega=\{0,1,2,\ldots\},
\]
а $S$ — функция следующего элемента:
\[
S(n)=n+1.
\]

У этой структуры нет нетривиальных автоморфизмов. Действительно, элемент $0$ является единственным элементом, который не лежит в образе функции $S$:
\[
0\notin S(\omega).
\]
Любой автоморфизм обязан сохранять это свойство, поэтому он фиксирует $0$.

Тогда он фиксирует и
\[
1=S(0),
\]
затем
\[
2=S(1),
\]
и так далее. Следовательно, любой автоморфизм структуры $\langle\omega,S\rangle$ является тождественным.
\end{example}

\subsection{Ограниченность метода автоморфизмов}

Автоморфизмы дают необходимое условие определимости:
\[
S\text{ определимо без параметров}
\qquad
\Longrightarrow
\qquad
\forall \pi\in Aut(\mathfrak A)\quad \pi(S)=S.
\]

Следовательно,
\[
\exists \pi\in Aut(\mathfrak A):\pi(S)\ne S
\qquad
\Longrightarrow
\qquad
S\text{ не определимо без параметров}.
\]

Для определимости с параметрами из $B\subseteq A$ аналогично:
\[
S\text{ определимо с параметрами из }B
\qquad
\Longrightarrow
\qquad
\forall \pi\in Aut(\mathfrak A),\ 
\bigl(\pi|_B=id_B \Rightarrow \pi(S)=S\bigr).
\]

Однако обратное утверждение в общем случае неверно:
\[
\forall \pi\in Aut(\mathfrak A)\quad \pi(S)=S
\qquad
\not\Longrightarrow
\qquad
S\text{ определимо}.
\]

Следующий пример показывает, почему метод автоморфизмов не всегда позволяет доказать неопределимость.

\begin{proposition}
У структуры упорядоченного кольца вещественных чисел
\[
\mathbb R_{o\text{-}ring}
=
\langle\mathbb R,+,\cdot,-,0,1,<\rangle
\]
нет нетривиальных автоморфизмов.
\end{proposition}

\begin{proof}
Пусть
\[
\pi\in Aut(\mathbb R_{o\text{-}ring}).
\]
Покажем, что $\pi$ является тождественным отображением.

Так как $0$ и $1$ являются константными символами языка, автоморфизм обязан их сохранять:
\[
\pi(0)=0,
\qquad
\pi(1)=1.
\]

Поскольку $\pi$ сохраняет сложение, для любого натурального числа $n$ имеем
\[
n=\underbrace{1+\cdots+1}_{n\text{ раз}},
\]
поэтому
\[
\pi(n)
=
\pi(\underbrace{1+\cdots+1}_{n\text{ раз}})
=
\underbrace{\pi(1)+\cdots+\pi(1)}_{n\text{ раз}}
=
n.
\]
Следовательно, все натуральные числа фиксируются.

Так как $\pi$ сохраняет унарный минус,
\[
\pi(-x)=-\pi(x).
\]
Поэтому для любого $n\in\mathbb N$
\[
\pi(-n)=-\pi(n)=-n.
\]
Значит, все целые числа фиксируются.

Пусть теперь
\[
q=\frac mn\in\mathbb Q,
\qquad n\ne0.
\]
Тогда
\[
m=n\cdot q.
\]
Так как $\pi$ сохраняет умножение,
\[
\pi(m)=\pi(n\cdot q)=\pi(n)\cdot\pi(q).
\]
Но $\pi(m)=m$ и $\pi(n)=n$, поэтому
\[
m=n\cdot\pi(q).
\]
Следовательно,
\[
\pi(q)=q.
\]
Значит, все рациональные числа фиксируются.

Осталось показать, что фиксируются все вещественные числа. Здесь используется то, что $\pi$ сохраняет порядок, а $\mathbb Q$ плотно в $\mathbb R$.

Пусть $r\in\mathbb R$. Предположим, что
\[
\pi(r)\ne r.
\]

Если
\[
r<\pi(r),
\]
то по плотности $\mathbb Q$ в $\mathbb R$ существует рациональное число $q\in\mathbb Q$ такое, что
\[
r<q<\pi(r).
\]
Из $r<q$ и сохранения порядка получаем
\[
\pi(r)<\pi(q).
\]
Но $\pi(q)=q$, значит,
\[
\pi(r)<q,
\]
что противоречит $q<\pi(r)$.

Если же
\[
\pi(r)<r,
\]
то по плотности $\mathbb Q$ в $\mathbb R$ существует рациональное число $q\in\mathbb Q$ такое, что
\[
\pi(r)<q<r.
\]
Из $q<r$ и сохранения порядка получаем
\[
\pi(q)<\pi(r).
\]
Но $\pi(q)=q$, значит,
\[
q<\pi(r),
\]
что противоречит $\pi(r)<q$.

Следовательно,
\[
\pi(r)=r
\]
для любого $r\in\mathbb R$.

Итак,
\[
Aut(\mathbb R_{o\text{-}ring})=\{id\}.
\]
\end{proof}

\begin{remark}
Этот пример показывает ограниченность метода автоморфизмов.

В структуре
\[
\mathbb R_{o\text{-}ring}
\]
любой автоморфизм является тождественным. Поэтому любое множество
\[
S\subseteq\mathbb R^n
\]
автоматически инвариантно относительно всех автоморфизмов этой структуры.

Однако из этого не следует, что любое множество $S\subseteq\mathbb R^n$ определимо.
\end{remark}

\begin{remark}
В структуре
\[
\mathbb R_{o\text{-}ring}
\]
существуют подмножества $\mathbb R^n$, которые не являются определимыми.

Действительно, язык
\[
L_{o\text{-}ring}
=
\langle +,\cdot,-,0,1,<\rangle
\]
счётен, поэтому формул этого языка тоже только счётно много. Каждая формула задаёт не более одного определимого множества.

Следовательно, определимых подмножеств $\mathbb R^n$ не более чем счётно много.

Но всех подмножеств $\mathbb R^n$ гораздо больше:
\[
|\mathcal P(\mathbb R^n)|=2^{|\mathbb R|}.
\]

Значит, большинство подмножеств $\mathbb R^n$ не определимы.
\end{remark}

\begin{remark}
Итак, можно пользоваться следующими практическими правилами.

Если удалось явно построить формулу, задающую множество $S$, то $S$ определимо.

Если удалось найти автоморфизм, который не сохраняет $S$, то $S$ не определимо без параметров.

Если множество сохраняется всеми автоморфизмами, это ещё не доказывает его определимость.
\end{remark}


\subsection{Бескванторные, экзистенциальные и универсальные формулы}

Пусть $\mathfrak A$ — $L$-структура.

\begin{definition}
$L$-формула называется \textbf{бескванторной}, если в ней нет кванторов
\[
\exists,\qquad \forall.
\]
Иными словами, бескванторные формулы строятся из атомарных формул при помощи логических связок:
\[
\neg,\qquad \wedge,\qquad \vee,\qquad \to,\qquad \leftrightarrow,
\]
но без использования кванторов.
\end{definition}

\begin{example}
Формулы
\[
x+y=0,
\]
\[
R(x,y)\wedge x\ne y,
\]
\[
(x<y)\vee(y<z)
\]
являются бескванторными.
\end{example}

\begin{definition}
$L$-формула называется \textbf{экзистенциальной}, если она имеет вид
\[
\exists y_1\ldots \exists y_m\,\psi(\bar x,\bar y),
\]
где $\psi(\bar x,\bar y)$ — бескванторная $L$-формула.
\end{definition}

То есть экзистенциальная формула говорит, что существуют некоторые элементы, при подстановке которых выполняется бескванторное условие.

\begin{example}
Формулы
\[
\exists y\,(x+y=0),
\]
\[
\exists y\exists z\,(x=y+z\wedge z\ne0)
\]
являются экзистенциальными.
\end{example}

\begin{definition}
$L$-формула называется \textbf{универсальной}, если она имеет вид
\[
\forall y_1\ldots \forall y_m\,\psi(\bar x,\bar y),
\]
где $\psi(\bar x,\bar y)$ — бескванторная $L$-формула.
\end{definition}

То есть универсальная формула говорит, что бескванторное условие выполняется для всех возможных значений некоторых переменных.

\begin{example}
Формулы
\[
\forall y\,(x+y=y+x),
\]
\[
\forall y\,(y\cdot x=y)
\]
являются универсальными.
\end{example}

\begin{remark}
Экзистенциальные и универсальные формулы — это не произвольные формулы с кванторами. В данном определении все кванторы стоят в начале формулы, а после них идёт бескванторная часть.

Например,
\[
\exists x\,\forall y\,\psi(x,y)
\]
не является экзистенциальной формулой в этом смысле, потому что после квантора существования появляется квантор всеобщности.
\end{remark}

\subsection{Сохранение формул при переходе к подструктурам}

Пусть
\[
\mathfrak A\subseteq \mathfrak B
\]
— $L$-подструктура. Это означает, что $A\subseteq B$, и интерпретации всех символов языка $L$ в $\mathfrak A$ согласованы с их интерпретациями в $\mathfrak B$.

В частности, если $\bar a\in A^n$, то значения термов на наборе $\bar a$ в структурах $\mathfrak A$ и $\mathfrak B$ совпадают.

\begin{lemma}
Пусть
\[
\mathfrak A\subseteq \mathfrak B,
\qquad
\bar a\in A^n.
\]
Тогда выполнены следующие утверждения.

\begin{enumerate}
    \item Если $\varphi(\bar x)$ — бескванторная $L$-формула, то
    \[
    \mathfrak A\models\varphi(\bar a)
    \qquad
    \Longleftrightarrow
    \qquad
    \mathfrak B\models\varphi(\bar a).
    \]

    \item Если $\varphi(\bar x)$ — экзистенциальная $L$-формула, то
    \[
    \mathfrak A\models\varphi(\bar a)
    \qquad
    \Longrightarrow
    \qquad
    \mathfrak B\models\varphi(\bar a).
    \]

    \item Если $\varphi(\bar x)$ — универсальная $L$-формула, то
    \[
    \mathfrak B\models\varphi(\bar a)
    \qquad
    \Longrightarrow
    \qquad
    \mathfrak A\models\varphi(\bar a).
    \]
\end{enumerate}
\end{lemma}

\begin{proof}
Докажем все три пункта.

\textbf{1. Бескванторные формулы.}

Сначала рассмотрим атомарные формулы. Для атомарных формул утверждение следует из определения подструктуры: термы на элементах из $A$ вычисляются одинаково в $\mathfrak A$ и в $\mathfrak B$, а отношения на наборах из $A$ согласованы.

То есть для атомарной формулы $\alpha(\bar x)$ имеем
\[
\mathfrak A\models\alpha(\bar a)
\qquad
\Longleftrightarrow
\qquad
\mathfrak B\models\alpha(\bar a).
\]

Теперь докажем утверждение для всех бескванторных формул индукцией по построению формулы.

Если формула имеет вид
\[
\neg\psi,
\]
то
\[
\mathfrak A\models\neg\psi(\bar a)
\]
тогда и только тогда, когда
\[
\mathfrak A\not\models\psi(\bar a).
\]
По индукционному предположению это равносильно
\[
\mathfrak B\not\models\psi(\bar a),
\]
а значит,
\[
\mathfrak B\models\neg\psi(\bar a).
\]

Если формула имеет вид
\[
\psi\wedge\theta,
\]
то
\[
\mathfrak A\models(\psi\wedge\theta)(\bar a)
\]
тогда и только тогда, когда
\[
\mathfrak A\models\psi(\bar a)
\qquad\text{и}\qquad
\mathfrak A\models\theta(\bar a).
\]
По индукционному предположению это равносильно
\[
\mathfrak B\models\psi(\bar a)
\qquad\text{и}\qquad
\mathfrak B\models\theta(\bar a),
\]
то есть
\[
\mathfrak B\models(\psi\wedge\theta)(\bar a).
\]

Остальные логические связки рассматриваются аналогично или выражаются через $\neg$ и $\wedge$.

Следовательно, для любой бескванторной формулы $\varphi(\bar x)$ выполнено
\[
\mathfrak A\models\varphi(\bar a)
\qquad
\Longleftrightarrow
\qquad
\mathfrak B\models\varphi(\bar a).
\]

\textbf{2. Экзистенциальные формулы.}

Пусть $\varphi(\bar x)$ — экзистенциальная формула. Тогда она имеет вид
\[
\varphi(\bar x)=\exists \bar y\,\psi(\bar x,\bar y),
\]
где $\psi(\bar x,\bar y)$ — бескванторная формула.

Пусть
\[
\mathfrak A\models\varphi(\bar a).
\]
Тогда
\[
\mathfrak A\models\exists \bar y\,\psi(\bar a,\bar y).
\]
По определению истинности квантора существования существует набор
\[
\bar b\in A^m
\]
такой, что
\[
\mathfrak A\models\psi(\bar a,\bar b).
\]

Так как $\psi$ бескванторная, по первому пункту получаем
\[
\mathfrak B\models\psi(\bar a,\bar b).
\]
Поскольку $\bar b\in A^m\subseteq B^m$, этот же набор $\bar b$ является допустимым свидетелем в структуре $\mathfrak B$.

Следовательно,
\[
\mathfrak B\models\exists \bar y\,\psi(\bar a,\bar y),
\]
то есть
\[
\mathfrak B\models\varphi(\bar a).
\]

\textbf{3. Универсальные формулы.}

Пусть $\varphi(\bar x)$ — универсальная формула. Тогда она имеет вид
\[
\varphi(\bar x)=\forall \bar y\,\psi(\bar x,\bar y),
\]
где $\psi(\bar x,\bar y)$ — бескванторная формула.

Пусть
\[
\mathfrak B\models\varphi(\bar a).
\]
Тогда
\[
\mathfrak B\models\forall \bar y\,\psi(\bar a,\bar y).
\]
Это значит, что для любого набора
\[
\bar b\in B^m
\]
выполнено
\[
\mathfrak B\models\psi(\bar a,\bar b).
\]

В частности, это верно для любого набора
\[
\bar b\in A^m,
\]
поскольку $A\subseteq B$. Тогда по первому пункту, так как $\psi$ бескванторная,
\[
\mathfrak A\models\psi(\bar a,\bar b).
\]
Следовательно, для любого $\bar b\in A^m$ выполнено
\[
\mathfrak A\models\psi(\bar a,\bar b).
\]
Значит,
\[
\mathfrak A\models\forall \bar y\,\psi(\bar a,\bar y),
\]
то есть
\[
\mathfrak A\models\varphi(\bar a).
\]
\end{proof}

\begin{remark}
Главная идея леммы состоит в следующем.

Бескванторные формулы проверяют только уже выбранные элементы и отношения между ними, поэтому они имеют одинаковую истинность в подструктуре и в расширении.

Экзистенциальные формулы сохраняются при переходе к расширению: если свидетель существовал в меньшей структуре, он остаётся существовать в большей.

Универсальные формулы сохраняются при переходе к подструктуре: если утверждение было верно для всех элементов большей структуры, то оно тем более верно для всех элементов меньшей структуры.
\end{remark}

\begin{remark}
Обратные направления во втором и третьем пунктах в общем случае неверны.

Экзистенциальная формула может стать истинной в расширении из-за появления новых элементов. Универсальная формула может быть истинной в подструктуре, но перестать быть истинной в расширении из-за появления новых элементов, на которых бескванторное условие нарушается.
\end{remark}

\begin{example}
Рассмотрим структуры
\[
\mathbb Z_{ring}
=
\langle\mathbb Z,+,\cdot,-,0,1\rangle
\subseteq
\langle\mathbb Q,+,\cdot,-,0,1\rangle
=
\mathbb Q_{ring}.
\]

Экзистенциальная формула
\[
\exists y\,(y+y=1)
\]
истинна в $\mathbb Q_{ring}$, потому что можно взять
\[
y=\frac12.
\]
Но она ложна в $\mathbb Z_{ring}$, потому что целого числа $y$, удовлетворяющего равенству
\[
y+y=1,
\]
не существует.

Это показывает, что экзистенциальные формулы не обязаны сохраняться при переходе от расширения к подструктуре.
\end{example}

\begin{example}
Рассмотрим структуры
\[
\mathbb N_{ring}
\subseteq
\mathbb Z_{ring}.
\]
Формула
\[
\forall y\,\neg(y+x=0)
\]
при параметре $x=1$ истинна в $\mathbb N_{ring}$, если сложение рассматривается на натуральных числах и $0,1\in\mathbb N$.

Действительно, в натуральных числах нет такого $y$, что
\[
y+1=0.
\]

Но в $\mathbb Z_{ring}$ эта формула ложна, потому что можно взять
\[
y=-1.
\]
Тогда
\[
-1+1=0.
\]

Это показывает, что универсальные формулы не обязаны сохраняться при переходе от подструктуры к расширению.
\end{example}


\begin{definition}
Пусть $\Sigma$ — множество $L$-предложений. Множество $\Sigma$ называется \textbf{выполнимым}, если существует $L$-структура $\mathfrak A$, такая что
\[
\mathfrak A\models\Sigma.
\]
\end{definition}

Запись
\[
\mathfrak A\models\Sigma
\]
означает, что структура $\mathfrak A$ удовлетворяет каждому предложению из $\Sigma$, то есть
\[
\mathfrak A\models\sigma
\]
для любого
\[
\sigma\in\Sigma.
\]

Иными словами, множество предложений выполнимо, если у него существует модель.

Если такой модели не существует, то множество предложений называется \textbf{невыполнимым}.

\begin{example}
Множество аксиом групп выполнимо, потому что существуют группы. Например,
\[
\langle\mathbb Z,+,-,0\rangle
\]
является моделью аксиом абелевых групп.
\end{example}

\begin{example}
Множество предложений
\[
\Sigma=\{\exists x\,(x\ne x)\}
\]
невыполнимо, потому что равенство в любой структуре рефлексивно:
\[
x=x
\]
для любого элемента $x$.
\end{example}


\subsection{Теории структур и дедуктивная замкнутость}

Пусть $\mathcal C$ — класс $L$-структур.

\begin{definition}
\textbf{Теорией класса структур} $\mathcal C$ называется множество всех $L$-предложений, истинных во всех структурах из $\mathcal C$:
\[
Th(\mathcal C)
=
\{\varphi\mid \mathfrak A\models\varphi
\text{ для любой } \mathfrak A\in\mathcal C\}.
\]
\end{definition}

Здесь $\varphi$ — именно $L$-предложение, то есть $L$-формула без свободных переменных. Это важно, потому что предложение имеет самостоятельное истинностное значение в структуре: оно либо истинно, либо ложно.

Если формула имеет свободные переменные, например $\varphi(x)$, то для проверки истинности нужно дополнительно указать, какой элемент структуры подставляется вместо $x$.

\begin{definition}
Если $\mathcal C=\{\mathfrak A\}$ состоит из одной структуры, то
\[
Th(\{\mathfrak A\})
\]
обычно обозначают короче:
\[
Th(\mathfrak A).
\]
Это множество называется \textbf{теорией структуры} $\mathfrak A$:
\[
Th(\mathfrak A)
=
\{\varphi\mid \mathfrak A\models\varphi\}.
\]
\end{definition}

Иными словами, $Th(\mathfrak A)$ — это множество всех предложений языка $L$, истинных в структуре $\mathfrak A$.

\begin{example}
Пусть
\[
\mathbb R_{o\text{-}ring}
=
\langle\mathbb R,+,\cdot,-,0,1,<\rangle.
\]
Тогда
\[
Th(\mathbb R_{o\text{-}ring})
\]
— это множество всех предложений первого порядка языка упорядоченных колец, истинных в вещественных числах.
\end{example}

\begin{definition}
Пусть $\Sigma$ — множество $L$-предложений. Будем говорить, что $\Sigma$ \textbf{дедуктивно замкнуто}, если для любого $L$-предложения $\varphi$
\[
\Sigma\models\varphi
\qquad
\Longrightarrow
\qquad
\varphi\in\Sigma.
\]
\end{definition}

Здесь запись
\[
\Sigma\models\varphi
\]
означает, что $\varphi$ является семантическим следствием $\Sigma$: во всякой модели $\Sigma$ предложение $\varphi$ истинно.

То есть
\[
\Sigma\models\varphi
\]
означает:
\[
\text{для любой } L\text{-структуры } \mathfrak A,\quad
\mathfrak A\models\Sigma
\Rightarrow
\mathfrak A\models\varphi.
\]

Смысл дедуктивной замкнутости состоит в следующем: если из $\Sigma$ уже логически следует некоторое предложение, то оно уже содержится в $\Sigma$.

\begin{definition}
\textbf{Дедуктивным замыканием} множества $L$-предложений $\Sigma$ называется множество
\[
\operatorname{Ded}(\Sigma)
=
\{\varphi\mid \Sigma\models\varphi\}.
\]
\end{definition}

То есть $\operatorname{Ded}(\Sigma)$ получается из $\Sigma$ добавлением всех предложений, которые являются его семантическими следствиями.

\begin{remark}
В доказательном варианте вместо семантического следования $\models$ часто используют выводимость $\vdash$ и пишут
\[
\operatorname{Ded}(\Sigma)
=
\{\varphi\mid \Sigma\vdash\varphi\}.
\]
Для логики первого порядка, благодаря теореме полноты Гёделя, эти два подхода согласованы:
\[
\Sigma\models\varphi
\qquad
\Longleftrightarrow
\qquad
\Sigma\vdash\varphi.
\]
\end{remark}

\begin{proposition}
Для любого класса $L$-структур $\mathcal C$ множество $Th(\mathcal C)$ дедуктивно замкнуто.
\end{proposition}

\begin{proof}
Пусть
\[
Th(\mathcal C)\models\varphi.
\]
Нужно доказать, что
\[
\varphi\in Th(\mathcal C).
\]

По определению $Th(\mathcal C)$ это значит, что нужно показать: $\varphi$ истинно во всех структурах из класса $\mathcal C$.

Возьмём произвольную структуру
\[
\mathfrak A\in\mathcal C.
\]
По определению $Th(\mathcal C)$ каждое предложение из $Th(\mathcal C)$ истинно в $\mathfrak A$. Поэтому
\[
\mathfrak A\models Th(\mathcal C).
\]

Так как
\[
Th(\mathcal C)\models\varphi,
\]
то всякая модель $Th(\mathcal C)$ удовлетворяет $\varphi$. Следовательно,
\[
\mathfrak A\models\varphi.
\]

Структура $\mathfrak A\in\mathcal C$ была выбрана произвольно. Значит, $\varphi$ истинно во всех структурах из $\mathcal C$. Поэтому
\[
\varphi\in Th(\mathcal C).
\]

Следовательно, $Th(\mathcal C)$ дедуктивно замкнуто.
\end{proof}

\begin{definition}
В этом курсе \textbf{теорией} будем называть выполнимое и дедуктивно замкнутое множество $L$-предложений.
\end{definition}

То есть множество $T$ является теорией, если выполнены два условия:

\[
T\text{ выполнимо},
\]
то есть существует $L$-структура $\mathfrak A$, такая что
\[
\mathfrak A\models T,
\]
и $T$ дедуктивно замкнуто:
\[
T\models\varphi
\qquad
\Longrightarrow
\qquad
\varphi\in T
\]
для любого $L$-предложения $\varphi$.

\begin{remark}
В разных курсах термин ``теория'' может использоваться немного по-разному. Иногда теорией называют произвольное множество предложений, а дедуктивно замкнутую теорию называют отдельно дедуктивно замкнутой теорией.

Здесь мы фиксируем более узкое соглашение: теория — это уже выполнимое и дедуктивно замкнутое множество предложений.
\end{remark}

\subsection{Полные теории}

\begin{definition}
Пусть $\Sigma$ — множество $L$-предложений. Будем говорить, что $\Sigma$ \textbf{полно}, если для любого $L$-предложения $\varphi$ выполнено:
\[
\Sigma\models\varphi
\qquad
\text{или}
\qquad
\Sigma\models\neg\varphi.
\]
\end{definition}

Иными словами, полная теория для каждого предложения языка решает, истинно оно или ложно во всех её моделях.

Если $T$ — дедуктивно замкнутая теория, то полнота равносильна следующему условию:
для любого $L$-предложения $\varphi$
\[
\varphi\in T
\qquad
\text{или}
\qquad
\neg\varphi\in T.
\]

\begin{proposition}
Для любой $L$-структуры $\mathfrak A$ теория
\[
Th(\mathfrak A)
\]
полна.
\end{proposition}

\begin{proof}
Пусть $\varphi$ — произвольное $L$-предложение.

В классической логике в фиксированной структуре $\mathfrak A$ предложение $\varphi$ либо истинно, либо ложно.

Если
\[
\mathfrak A\models\varphi,
\]
то
\[
\varphi\in Th(\mathfrak A).
\]

Если же
\[
\mathfrak A\not\models\varphi,
\]
то
\[
\mathfrak A\models\neg\varphi,
\]
а значит,
\[
\neg\varphi\in Th(\mathfrak A).
\]

Следовательно, для любого $L$-предложения $\varphi$ выполнено:
\[
\varphi\in Th(\mathfrak A)
\qquad
\text{или}
\qquad
\neg\varphi\in Th(\mathfrak A).
\]
Значит, $Th(\mathfrak A)$ полна.
\end{proof}

\begin{remark}
Теория одной конкретной структуры всегда полна. Теория класса структур в общем случае может быть неполной.

Причина в том, что в одной структуре каждое предложение уже имеет конкретное истинностное значение, а в разных структурах одного класса одно и то же предложение может вести себя по-разному.
\end{remark}

\subsection{Истинная арифметика и арифметика Пеано}

Рассмотрим структуру натуральных чисел в языке полуколец:
\[
\omega_{s\text{-}ring}
=
\langle\omega,+,\cdot,0,1\rangle,
\]
где
\[
\omega=\{0,1,2,\ldots\}.
\]

\begin{definition}
Теория
\[
Th(\omega_{s\text{-}ring})
\]
называется \textbf{истинной арифметикой}.
\end{definition}

Иными словами,
\[
Th(\omega_{s\text{-}ring})
\]
— это множество всех предложений первого порядка языка
\[
L_{s\text{-}ring}
=
\langle +,\cdot,0,1\rangle,
\]
которые истинны в стандартной структуре натуральных чисел
\[
\langle\omega,+,\cdot,0,1\rangle.
\]

Поскольку истинная арифметика является теорией одной конкретной структуры, она полна:
\[
Th(\omega_{s\text{-}ring})
\text{ полна}.
\]

Пусть теперь $\Sigma$ — множество аксиом Пеано в языке
\[
L_{s\text{-}ring}
=
\langle +,\cdot,0,1\rangle.
\]

\begin{definition}
\textbf{Арифметикой Пеано} называется дедуктивное замыкание множества аксиом Пеано:
\[
PA=\operatorname{Ded}(\Sigma).
\]
\end{definition}

То есть
\[
PA
=
\{\varphi\mid \Sigma\models\varphi\}.
\]

В доказательной записи можно также писать
\[
PA
=
\{\varphi\mid \Sigma\vdash\varphi\}.
\]

Стандартная структура натуральных чисел является моделью аксиом Пеано:
\[
\omega_{s\text{-}ring}\models PA.
\]

Следовательно, всякое предложение, принадлежащее $PA$, истинно в стандартной арифметике:
\[
PA\subseteq Th(\omega_{s\text{-}ring}).
\]

\begin{remark}
Интуитивно:
\[
PA
\]
— это всё, что следует из аксиом Пеано, а
\[
Th(\omega_{s\text{-}ring})
\]
— это всё, что на самом деле истинно в стандартных натуральных числах.

Поэтому истинная арифметика содержит не только следствия аксиом Пеано, но вообще все предложения, истинные в стандартной модели.
\end{remark}

\begin{proposition}
Теория $PA$ неполна.
\end{proposition}

\begin{remark}
Это утверждение является содержанием первой теоремы Гёделя о неполноте в применении к арифметике Пеано.

Смысл состоит в том, что существуют предложения языка арифметики, для которых $PA$ не доказывает ни само предложение, ни его отрицание:
\[
PA\not\models\varphi
\qquad
\text{и}
\qquad
PA\not\models\neg\varphi.
\]

При этом истинная арифметика
\[
Th(\omega_{s\text{-}ring})
\]
полна, потому что она является теорией одной структуры.
\end{remark}

Таким образом, имеем следующую картину:
\[
\Sigma=\text{аксиомы Пеано},
\]
\[
PA=\operatorname{Ded}(\Sigma),
\]
\[
Th(\omega_{s\text{-}ring})
=
\text{истинная арифметика}.
\]

Причём
\[
PA\subseteq Th(\omega_{s\text{-}ring}),
\]
но
\[
PA\ne Th(\omega_{s\text{-}ring}).
\]

\begin{remark}
Различие между $PA$ и истинной арифметикой принципиально.

$PA$ — это эффективно заданная аксиоматическая теория. Она строится из явно заданных аксиом и всех их логических следствий.

Истинная арифметика — это полная теория стандартной структуры натуральных чисел. Она содержит все истинные арифметические предложения, но не задаётся простым эффективным списком аксиом.

Именно это различие лежит в основе феномена неполноты.
\end{remark}

\section{Алгебраически замкнутые поля и упорядоченные поля}

В этом разделе мы рассмотрим несколько важных классов полей с точки зрения теории моделей:
алгебраически замкнутые поля, упорядоченные поля и полные упорядоченные поля.

Начнём с обычного языка колец:
\[
L_{ring}=\langle +,\cdot,-,0,1\rangle.
\]
Если мы хотим говорить об упорядоченных полях, то добавляем к языку бинарный символ отношения порядка:
\[
L_{o\text{-}ring}=\langle +,\cdot,-,0,1,<\rangle.
\]

\begin{definition}
Полем называется $L_{ring}$-структура
\[
K=\langle K,+,\cdot,-,0,1\rangle,
\]
в которой выполнены обычные аксиомы поля:
сложение и умножение удовлетворяют законам коммутативного кольца с единицей, $0\ne 1$, и каждый ненулевой элемент имеет обратный по умножению.
\end{definition}

Аксиомы поля можно записать предложениями первого порядка в языке $L_{ring}$. Например, аксиома существования обратного элемента имеет вид
\[
\forall x\,(x\ne 0\to \exists y\,(x\cdot y=1)).
\]

\begin{definition}
Пусть $K$ — поле. Характеристикой поля $K$ называется наименьшее натуральное число $n>0$, такое что
\[
\underbrace{1+\cdots+1}_{n\text{ раз}}=0,
\]
если такое $n$ существует.

Если такого $n$ не существует, говорят, что поле имеет характеристику $0$.
\end{definition}

\begin{example}
Поле рациональных чисел $\mathbb Q$, поле вещественных чисел $\mathbb R$ и поле комплексных чисел $\mathbb C$ имеют характеристику $0$.
\end{example}

\begin{example}
Для простого числа $p$ поле
\[
\mathbb F_p=\mathbb Z/p\mathbb Z
\]
имеет характеристику $p$, потому что
\[
\underbrace{1+\cdots+1}_{p\text{ раз}}=0
\]
в этом поле.
\end{example}

\subsection{Алгебраически замкнутые поля}

\begin{definition}
Поле $K$ называется \textbf{алгебраически замкнутым}, если всякий непостоянный многочлен с коэффициентами из $K$ имеет корень в $K$.
\end{definition}

То есть для любого многочлена
\[
f(x)=a_nx^n+\cdots+a_1x+a_0
\]
с коэффициентами $a_0,\ldots,a_n\in K$, если
\[
n\ge 1
\qquad\text{и}\qquad
a_n\ne 0,
\]
то существует $b\in K$, такой что
\[
f(b)=0.
\]

\begin{example}
Поле комплексных чисел $\mathbb C$ алгебраически замкнуто. Это содержание основной теоремы алгебры.
\end{example}

\begin{example}
Поле вещественных чисел $\mathbb R$ не является алгебраически замкнутым, потому что многочлен
\[
x^2+1
\]
не имеет корней в $\mathbb R$.
\end{example}

\begin{remark}
Свойство быть алгебраически замкнутым можно аксиоматизировать в языке колец бесконечной схемой предложений.

Для каждого $n\ge 1$ добавляется предложение:
\[
\forall a_0\ldots \forall a_n
\left(
a_n\ne 0
\to
\exists x\,
(a_nx^n+a_{n-1}x^{n-1}+\cdots+a_1x+a_0=0)
\right).
\]
Это предложение говорит, что всякий многочлен степени ровно $n$ имеет корень.
\end{remark}

\begin{definition}
Теория алгебраически замкнутых полей обозначается
\[
ACF.
\]
Она состоит из аксиом поля и схемы аксиом, утверждающей, что всякий непостоянный многочлен имеет корень.
\end{definition}

\begin{remark}
Иногда также фиксируют характеристику. Например,
\[
ACF_0
\]
обозначает теорию алгебраически замкнутых полей характеристики $0$.

Условие характеристики $0$ задаётся бесконечным набором предложений:
\[
1\ne 0,
\]
\[
1+1\ne 0,
\]
\[
1+1+1\ne 0,
\]
и так далее. То есть для каждого $n>0$ добавляется аксиома
\[
\underbrace{1+\cdots+1}_{n\text{ раз}}\ne 0.
\]
\end{remark}

\subsection{Упорядоченные поля}

Теперь добавим к языку поля символ порядка $<$.

\begin{definition}
\textbf{Упорядоченным полем} называется $L_{o\text{-}ring}$-структура
\[
K=\langle K,+,\cdot,-,0,1,<\rangle,
\]
такая что:
\begin{enumerate}
    \item $\langle K,+,\cdot,-,0,1\rangle$ является полем;
    \item $<$ является линейным порядком на $K$;
    \item порядок согласован со сложением:
    \[
    \forall x\forall y\forall z\,
    (x<y\to x+z<y+z);
    \]
    \item порядок согласован с умножением на положительные элементы:
    \[
    \forall x\forall y\,
    (0<x\wedge 0<y\to 0<x\cdot y).
    \]
\end{enumerate}
\end{definition}

Смысл этих условий такой: порядок должен вести себя так же, как обычный порядок на числовом поле. Если к обеим частям неравенства прибавить один и тот же элемент, неравенство сохраняется. А произведение положительных элементов положительно.

\begin{example}
Структура
\[
\langle\mathbb Q,+,\cdot,-,0,1,<\rangle
\]
является упорядоченным полем.
\end{example}

\begin{example}
Структура
\[
\langle\mathbb R,+,\cdot,-,0,1,<\rangle
\]
является упорядоченным полем.
\end{example}

\begin{example}
Поле комплексных чисел $\mathbb C$ нельзя сделать упорядоченным полем так, чтобы порядок был согласован с операциями поля.
\end{example}

\begin{proposition}
Всякое упорядоченное поле имеет характеристику $0$.
\end{proposition}

\begin{proof}
Пусть
\[
K=\langle K,+,\cdot,-,0,1,<\rangle
\]
— упорядоченное поле.

В упорядоченном поле обязательно выполнено
\[
0<1.
\]
Действительно, так как $1\ne 0$, по линейности порядка либо
\[
0<1,
\]
либо
\[
1<0.
\]
Если бы было $1<0$, то, прибавив $-1$ к обеим частям, получили бы
\[
0<-1.
\]
Тогда $-1$ был бы положительным элементом. Произведение положительных элементов положительно, значит
\[
(-1)\cdot(-1)>0.
\]
Но
\[
(-1)\cdot(-1)=1,
\]
то есть снова получаем $1>0$. Поэтому в любом случае $1$ положителен.

Теперь для любого $n>0$ сумма
\[
\underbrace{1+\cdots+1}_{n\text{ раз}}
\]
является положительным элементом, так как сумма положительных элементов положительна. Следовательно,
\[
\underbrace{1+\cdots+1}_{n\text{ раз}}\ne 0
\]
для любого $n>0$.

Значит, не существует положительного натурального $n$, такого что
\[
n\cdot 1=0.
\]
Следовательно, характеристика поля $K$ равна $0$.
\end{proof}

\begin{remark}
Из этого следует, что конечное поле не может быть упорядоченным полем. Всякое конечное поле имеет положительную характеристику, а всякое упорядоченное поле имеет характеристику $0$.
\end{remark}

\subsection{Полные упорядоченные поля}

Среди упорядоченных полей особую роль играет поле вещественных чисел. Его отличает свойство полноты.

\begin{definition}
Пусть
\[
K=\langle K,+,\cdot,-,0,1,<\rangle
\]
— упорядоченное поле.

Подмножество $S\subseteq K$ называется \textbf{ограниченным сверху}, если существует элемент $b\in K$, такой что
\[
s\le b
\]
для любого $s\in S$.

Такой элемент $b$ называется верхней гранью множества $S$.
\end{definition}

\begin{definition}
Элемент $c\in K$ называется \textbf{наименьшей верхней гранью} множества $S\subseteq K$, если:
\begin{enumerate}
    \item $c$ является верхней гранью $S$, то есть
    \[
    s\le c
    \]
    для любого $s\in S$;
    \item если $b$ — любая верхняя грань $S$, то
    \[
    c\le b.
    \]
\end{enumerate}
\end{definition}

Наименьшую верхнюю грань множества $S$ часто обозначают
\[
\sup S.
\]

\begin{definition}
Упорядоченное поле $K$ называется \textbf{полным}, если каждое непустое ограниченное сверху подмножество
\[
S\subseteq K
\]
имеет наименьшую верхнюю грань в $K$.
\end{definition}

Иначе говоря, если
\[
S\ne\varnothing
\]
и существует верхняя грань $S$, то существует элемент
\[
\sup S\in K.
\]

\begin{remark}
Это свойство называется аксиомой полноты или свойством полноты по Дедекинду.

Важно: это не предложение первого порядка в языке
\[
L_{o\text{-}ring}.
\]
Причина в том, что в формулировке квантор идёт по всем подмножествам $S\subseteq K$, а язык первого порядка позволяет квантифицировать только по элементам структуры, но не по её подмножествам.

Поэтому полнота упорядоченного поля является свойством второго порядка.
\end{remark}

\begin{example}
Поле рациональных чисел
\[
\mathbb Q
\]
с обычным порядком не является полным.

Например, множество
\[
S=\{q\in\mathbb Q\mid q^2<2,\ q>0\}
\]
непусто и ограничено сверху в $\mathbb Q$, но не имеет наименьшей верхней грани в $\mathbb Q$.

Его наименьшая верхняя грань в вещественных числах равна
\[
\sqrt 2,
\]
но
\[
\sqrt 2\notin\mathbb Q.
\]
\end{example}

\begin{example}
Поле вещественных чисел
\[
\mathbb R
\]
с обычным порядком является полным упорядоченным полем.
\end{example}

\subsection{Архимедовость}

\begin{definition}
Упорядоченное поле $K$ называется \textbf{архимедовым}, если для любого элемента $x\in K$ существует натуральное число $n$ такое, что
\[
x<n.
\]
\end{definition}

Здесь натуральное число $n$ рассматривается как элемент поля:
\[
n=\underbrace{1+\cdots+1}_{n\text{ раз}}.
\]

Интуитивно архимедовость означает, что в поле нет бесконечно больших элементов. Каким бы большим ни был элемент $x$, его можно превзойти некоторым натуральным числом.

\begin{lemma}
Если упорядоченное поле полно, то оно архимедово.
\end{lemma}

\begin{proof}
Пусть $K$ — полное упорядоченное поле. Докажем, что $K$ архимедово.

Предположим противное. Тогда существует элемент $x\in K$, такой что
\[
n\le x
\]
для любого натурального числа $n$.

Рассмотрим множество натуральных элементов внутри поля:
\[
N_K=\{1,2,3,\ldots\}\subseteq K.
\]
По предположению это множество ограничено сверху элементом $x$.

Множество $N_K$ непусто и ограничено сверху. Так как поле $K$ полно, существует его наименьшая верхняя грань:
\[
c=\sup N_K.
\]

Так как $c$ — наименьшая верхняя грань, элемент
\[
c-1
\]
не может быть верхней гранью множества $N_K$. Поэтому существует натуральное число $n\in N_K$, такое что
\[
c-1<n.
\]
Прибавляя $1$ к обеим частям, получаем
\[
c<n+1.
\]

Но если $n$ — натуральное число, то $n+1$ тоже натуральное число, то есть
\[
n+1\in N_K.
\]
Так как $c$ является верхней гранью $N_K$, должно быть
\[
n+1\le c.
\]
Получили противоречие:
\[
c<n+1\le c.
\]

Следовательно, наше предположение было неверным. Значит, для любого $x\in K$ существует натуральное число $n$, такое что
\[
x<n.
\]
То есть поле $K$ архимедово.
\end{proof}

\begin{remark}
Архимедовость является следствием полноты, но не наоборот.

Например, поле рациональных чисел $\mathbb Q$ архимедово, но не полно.
\end{remark}

\subsection{Единственность полного упорядоченного поля}

Следующий факт является одним из фундаментальных результатов о вещественных числах.

\begin{theorem}
Существует ровно одно полное упорядоченное поле с точностью до изоморфизма.
\end{theorem}

Иначе говоря, если $K$ и $F$ — полные упорядоченные поля, то существует изоморфизм упорядоченных полей
\[
K\cong F.
\]

\begin{remark}
Этот факт означает, что вещественные числа можно охарактеризовать как полное упорядоченное поле.

Любые две конструкции вещественных чисел, например через сечения Дедекинда или через классы фундаментальных последовательностей рациональных чисел, дают изоморфные упорядоченные поля.
\end{remark}

\begin{remark}
С точки зрения теории моделей важно помнить, что эта характеристика вещественных чисел использует полноту, а полнота не является свойством первого порядка в языке упорядоченных полей.

Поэтому выражение ``$\mathbb R$ — единственное полное упорядоченное поле'' не означает, что $\mathbb R$ единственно задаётся некоторым набором предложений первого порядка языка $L_{o\text{-}ring}$.
\end{remark}

\subsection{Конечная выполнимость}

Теперь перейдём к одному из ключевых результатов теории моделей — теореме компактности.

\begin{definition}
Пусть $\Sigma$ — множество $L$-предложений. Говорят, что $\Sigma$ \textbf{конечно выполнимо}, если любое конечное подмножество
\[
\Sigma_0\subseteq\Sigma
\]
выполнимо.
\end{definition}

То есть для любого конечного набора предложений
\[
\sigma_1,\ldots,\sigma_n\in\Sigma
\]
существует $L$-структура $\mathfrak A$, такая что
\[
\mathfrak A\models\sigma_1,\ldots,\sigma_n.
\]

Иначе говоря, всякая конечная часть $\Sigma$ имеет модель.

\begin{remark}
Если всё множество $\Sigma$ выполнимо, то оно, конечно, конечно выполнимо.

Действительно, если существует структура $\mathfrak A$, такая что
\[
\mathfrak A\models\Sigma,
\]
то эта же структура удовлетворяет любому конечному подмножеству $\Sigma_0\subseteq\Sigma$.
\end{remark}

Обратное утверждение совсем не очевидно. Оно и является содержанием теоремы компактности.

\begin{theorem}[Теорема компактности, локальная теорема Мальцева]
Пусть $\Sigma$ — множество $L$-предложений. Тогда $\Sigma$ выполнимо тогда и только тогда, когда $\Sigma$ конечно выполнимо.
\end{theorem}

То есть
\[
\Sigma\text{ имеет модель}
\qquad
\Longleftrightarrow
\qquad
\text{любое конечное подмножество }\Sigma\text{ имеет модель}.
\]

\begin{remark}
Название ``компактность'' связано с тем, что бесконечная система требований имеет модель, если модель имеет каждая её конечная часть.

Это один из центральных принципов теории моделей: если противоречие есть, то оно уже обнаруживается на конечном фрагменте теории.
\end{remark}

\begin{remark}
Теорема компактности является чисто первопорядковым результатом. Она верна для логики первого порядка, но не сохраняется в таком виде для более сильных логик, например для полной логики второго порядка.
\end{remark}

\subsection{Алгебраические замыкания полей}

Теперь применим идею существования моделей к полям.

\begin{definition}
Пусть $K$ — поле. Расширение поля
\[
L/K
\]
называется \textbf{алгебраическим}, если каждый элемент $a\in L$ является алгебраическим над $K$, то есть существует ненулевой многочлен
\[
f(x)\in K[x],
\]
такой что
\[
f(a)=0.
\]
\end{definition}

\begin{definition}
Поле $\overline K$ называется \textbf{алгебраическим замыканием} поля $K$, если:
\begin{enumerate}
    \item $\overline K$ является алгебраическим расширением $K$;
    \item $\overline K$ алгебраически замкнуто.
\end{enumerate}
\end{definition}

\begin{theorem}
Для любого поля $K$ существует его алгебраическое замыкание.
\end{theorem}

\begin{proof}
Дадим модельно-теоретическую идею доказательства через компактность.

Рассмотрим язык колец, расширенный константными символами для элементов поля $K$:
\[
L_K=L_{ring}\cup\{c_a\mid a\in K\}.
\]

Мы хотим построить алгебраически замкнутое поле, содержащее копию поля $K$.

Рассмотрим множество предложений $\Sigma$, состоящее из следующих частей.

Во-первых, добавим аксиомы полей.

Во-вторых, добавим аксиомы алгебраической замкнутости: каждый непостоянный многочлен имеет корень.

В-третьих, добавим диаграмму поля $K$, то есть предложения, фиксирующие сложение, умножение и неравенства между константами из $K$:
\[
c_{a+b}=c_a+c_b,
\]
\[
c_{a\cdot b}=c_a\cdot c_b,
\]
\[
c_{-a}=-c_a,
\]
\[
c_0=0,
\qquad
c_1=1,
\]
а также
\[
c_a\ne c_b
\]
для всех различных $a,b\in K$.

Если у $\Sigma$ есть модель, то эта модель является алгебраически замкнутым полем, в которое поле $K$ вкладывается по правилу
\[
a\mapsto c_a.
\]

Покажем, что $\Sigma$ конечно выполнимо.

Возьмём конечное подмножество $\Sigma_0\subseteq\Sigma$. В нём участвует только конечное число требований вида ``такой-то многочлен имеет корень''. Эти требования можно удовлетворить, последовательно присоединяя к $K$ корни конечного числа многочленов.

То есть существует некоторое расширение поля $K$, в котором выполнены все предложения из $\Sigma_0$.

Следовательно, всякое конечное подмножество $\Sigma$ выполнимо. Значит, по теореме компактности всё множество $\Sigma$ выполнимо.

Пусть $M$ — модель $\Sigma$. Тогда $M$ является алгебраически замкнутым полем, содержащим копию поля $K$.

Теперь внутри $M$ рассмотрим множество всех элементов, алгебраических над $K$. Обозначим это множество через
\[
\overline K.
\]
Стандартный факт алгебры говорит, что элементы, алгебраические над $K$ внутри расширения поля, образуют поле.

Кроме того, так как $M$ алгебраически замкнуто, поле $\overline K$ является алгебраически замкнутым: корни многочленов с коэффициентами из $\overline K$ лежат в $M$, а затем являются алгебраическими над $\overline K$ и, следовательно, алгебраическими над $K$.

Значит,
\[
\overline K
\]
является алгебраическим и алгебраически замкнутым расширением поля $K$.

Следовательно, $\overline K$ — алгебраическое замыкание поля $K$.
\end{proof}

\begin{remark}
Существование алгебраического замыкания можно доказывать и чисто алгебраически, например с помощью леммы Цорна.

Модельно-теоретический подход показывает, как теорема компактности позволяет строить модели бесконечных систем требований, проверяя только конечные фрагменты.
\end{remark}


\section{Фильтры}

\subsection{Идея фильтра}

Фильтры возникают как способ формально говорить о том, какие подмножества некоторого множества $I$ мы считаем \textbf{большими}.

Пусть $I$ — непустое множество. Обычно элементы $I$ можно понимать как индексы. Например, если
\[
I=\omega=\{0,1,2,\ldots\},
\]
то функции
\[
x:\omega\to\mathbb R
\]
являются обычными вещественными последовательностями.

Фильтр на $I$ позволяет сказать, какие множества индексов считаются большими, существенными или происходящими ``почти всюду''.

Идея такая:
\[
A\in\mathcal F
\]
означает, что множество $A\subseteq I$ считается большим.

Тогда дополнения к большим множествам можно считать малыми:
\[
A\in\mathcal F
\qquad
\Longrightarrow
\qquad
I\setminus A \text{ мало}.
\]

Например, при обычной сходимости последовательности большими считаются множества индексов, содержащие все достаточно большие номера. То есть множество индексов может не быть всем $\omega$, но если оно пропускает только конечное число номеров, то его всё равно считают большим.


\begin{definition}
Пусть $I$ — непустое множество. \textbf{Фильтром} на $I$ называется семейство подмножеств
\[
\mathcal F\subseteq \mathcal P(I),
\]
удовлетворяющее следующим условиям:
\begin{enumerate}
    \item
    \[
    I\in\mathcal F,
    \qquad
    \varnothing\notin\mathcal F;
    \]

    \item если
    \[
    A\in\mathcal F
    \qquad
    \text{и}
    \qquad
    A\subseteq B\subseteq I,
    \]
    то
    \[
    B\in\mathcal F;
    \]

    \item если
    \[
    A\in\mathcal F
    \qquad
    \text{и}
    \qquad
    B\in\mathcal F,
    \]
    то
    \[
    A\cap B\in\mathcal F.
    \]
\end{enumerate}
\end{definition}


\begin{definition}
Пусть $I$ — бесконечное множество. \textbf{Фильтром Фреше (Кофинитным фильтром)} на $I$ называется семейство
\[
\operatorname{cof}(I)
=
\{A\subseteq I\mid I\setminus A \text{ конечно}\}.
\]
\end{definition}

То есть множество $A\subseteq I$ принадлежит кофинитному фильтру, если его дополнение конечно.

Иными словами, $A$ считается большим, если оно содержит почти все элементы $I$, кроме, возможно, конечного числа исключений.

\begin{proposition}
Если $I$ бесконечно, то
\[
\operatorname{cof}(I)
=
\{A\subseteq I\mid I\setminus A \text{ конечно}\}
\]
является фильтром на $I$.
\end{proposition}

\begin{proof}
Проверим условия фильтра.

Во-первых,
\[
I\in\operatorname{cof}(I),
\]
поскольку
\[
I\setminus I=\varnothing
\]
конечно.

Также
\[
\varnothing\notin\operatorname{cof}(I),
\]
поскольку
\[
I\setminus\varnothing=I
\]
бесконечно.

Во-вторых, пусть
\[
A\in\operatorname{cof}(I)
\]
и
\[
A\subseteq B\subseteq I.
\]
Тогда
\[
I\setminus B\subseteq I\setminus A.
\]
Так как $I\setminus A$ конечно, то и $I\setminus B$ конечно. Значит,
\[
B\in\operatorname{cof}(I).
\]

В-третьих, пусть
\[
A,B\in\operatorname{cof}(I).
\]
Тогда $I\setminus A$ и $I\setminus B$ конечны. По формуле де Моргана
\[
I\setminus(A\cap B)
=
(I\setminus A)\cup(I\setminus B).
\]
Объединение двух конечных множеств конечно. Следовательно,
\[
I\setminus(A\cap B)
\]
конечно, а значит,
\[
A\cap B\in\operatorname{cof}(I).
\]

Итак, $\operatorname{cof}(I)$ является фильтром.
\end{proof}

\begin{remark}
Кофинитный фильтр (Фреше) — это фильтр, соответствующий выражению ``почти все элементы''. В этом смысле множество считается большим, если оно отличается от всего $I$ только на конечное множество.
\end{remark}

\subsection{Фильтры и сходимость последовательностей}

Рассмотрим случай
\[
I=\omega=\{0,1,2,\ldots\}.
\]
Функция
\[
x:\omega\to\mathbb R
\]
является вещественной последовательностью
\[
(x_0,x_1,x_2,\ldots).
\]

Обычную сходимость последовательности можно переформулировать через фильтр Фреше.

Пусть
\[
x_n\to a.
\]
Это означает, что для любого $\varepsilon>0$ существует номер $N$, такой что для всех $n\ge N$
\[
|x_n-a|<\varepsilon.
\]

Иными словами, для любого $\varepsilon>0$ множество индексов
\[
A_\varepsilon
=
\{n\in\omega\mid |x_n-a|<\varepsilon\}
\]
содержит все достаточно большие номера. Следовательно,
\[
A_\varepsilon\in\operatorname{cof}(\omega).
\]

Поэтому обычную сходимость можно записать так:
\[
x_n\to a
\qquad
\Longleftrightarrow
\qquad
\forall\varepsilon>0\ 
\{n\in\omega\mid |x_n-a|<\varepsilon\}\in\operatorname{cof}(\omega).
\]

\begin{remark}
Таким образом, обычная сходимость последовательностей — это сходимость относительно кофинитного фильтра на множестве индексов $\omega$.
\end{remark}

\begin{example}
Постоянная последовательность
\[
(a,a,a,\ldots)
\]
сходится к $a$.

Действительно, для любого $\varepsilon>0$
\[
\{n\in\omega\mid |x_n-a|<\varepsilon\}=\omega.
\]
А
\[
\omega\in\operatorname{cof}(\omega).
\]
\end{example}

\begin{example}
Последовательность
\[
\left(1,\frac12,\frac13,\frac14,\ldots\right)
\]
сходится к $0$.

Для любого $\varepsilon>0$ начиная с некоторого номера $N$ выполнено
\[
\frac1n<\varepsilon.
\]
Следовательно, множество
\[
\left\{n\in\omega\mid \frac1n<\varepsilon\right\}
\]
кофинитно.
\end{example}

\begin{example}
Последовательность
\[
\left(1,\frac14,\frac19,\frac1{16},\ldots\right)
=
\left(\frac1{n^2}\right)
\]
также сходится к $0$.

Для любого $\varepsilon>0$ множество индексов, на которых
\[
\frac1{n^2}<\varepsilon,
\]
содержит все достаточно большие $n$, то есть является кофинитным.
\end{example}

\begin{example}
Последовательность
\[
(1,2,3,\ldots)
\]
не сходится в $\mathbb R$.

Если бы она сходилась к некоторому $a\in\mathbb R$, то для $\varepsilon=1$ множество
\[
\{n\in\omega\mid |n-a|<1\}
\]
должно было бы быть кофинитным. Но оно конечно. Следовательно, сходимости нет.
\end{example}

\begin{example}
Последовательность
\[
(0,1,0,1,\ldots)
\]
не имеет обычного предела.

Например, если проверять возможный предел $0$, то для $\varepsilon=\frac12$ получаем множество индексов
\[
\{n\in\omega\mid |x_n-0|<1/2\}.
\]
Это множество чётных индексов. Оно бесконечно, но не кофинитно, потому что его дополнение тоже бесконечно.

Аналогично последовательность не сходится к $1$.
\end{example}

\subsection{Примеры фильтров}

\begin{example}[фильтр Фреше]
Пусть $I$ — бесконечное множество. Тогда
\[
\operatorname{cof}(I)
=
\{A\subseteq I\mid I\setminus A \text{ конечно}\}
\]
является фильтром.

В этом фильтре большими считаются множества, дополнение к которым конечно.
\end{example}

\begin{example}[Главный фильтр, порождённый множеством]
Пусть $J\subseteq I$ и $J\ne\varnothing$. Определим
\[
\mathcal F_J
=
\{A\subseteq I\mid J\subseteq A\}.
\]
Тогда $\mathcal F_J$ является фильтром на $I$.

Действительно, $I\in\mathcal F_J$, поскольку $J\subseteq I$, а $\varnothing\notin\mathcal F_J$, поскольку $J\ne\varnothing$.

Если
\[
A\in\mathcal F_J
\qquad
\text{и}
\qquad
A\subseteq B\subseteq I,
\]
то
\[
J\subseteq A\subseteq B,
\]
значит,
\[
B\in\mathcal F_J.
\]

Если
\[
A,B\in\mathcal F_J,
\]
то
\[
J\subseteq A
\qquad
\text{и}
\qquad
J\subseteq B.
\]
Следовательно,
\[
J\subseteq A\cap B,
\]
значит,
\[
A\cap B\in\mathcal F_J.
\]

Такой фильтр называется главным фильтром, порождённым множеством $J$.
\end{example}

\begin{example}[Главный фильтр в точке]
Пусть $a\in I$. Тогда семейство
\[
\mathcal F_a
=
\{A\subseteq I\mid a\in A\}
\]
является фильтром на $I$.

Это частный случай предыдущего примера при
\[
J=\{a\}.
\]

В этом фильтре большими считаются все множества, содержащие фиксированную точку $a$.
\end{example}

\begin{example}[Фильтр окрестностей точки]
Пусть $I=\mathbb R$ и $a\in\mathbb R$. Рассмотрим семейство
\[
\mathcal N(a)
=
\{U\subseteq\mathbb R\mid U\text{ содержит некоторую окрестность точки }a\}.
\]
Тогда $\mathcal N(a)$ является фильтром.

Действительно, всё пространство $\mathbb R$ содержит окрестность точки $a$, а пустое множество не содержит никакой окрестности точки $a$.

Если $U$ содержит окрестность точки $a$ и
\[
U\subseteq V,
\]
то $V$ тоже содержит окрестность точки $a$.

Если $U$ и $V$ содержат окрестности точки $a$, то их пересечение $U\cap V$ также содержит некоторую окрестность точки $a$.

Этот фильтр называется фильтром окрестностей точки $a$.
\end{example}

\begin{example}[Фильтр множеств полной меры]
Рассмотрим множество
\[
I=\mathbb R.
\]
Определим семейство
\[
\mathcal F_{\mathrm{co\text{-}null}}
=
\{A\subseteq\mathbb R\mid \mathbb R\setminus A
\text{ имеет меру Лебега }0\}.
\]

Тогда $\mathcal F_{\mathrm{co\text{-}null}}$ является фильтром, если рассматривать измеримые множества и множества, дополнения к которым имеют меру нуль.

В этом фильтре большими считаются множества, которые имеют полную меру, то есть множества, чьё дополнение имеет меру нуль.

Этот пример соответствует выражению ``почти всюду'' в теории меры.
\end{example}


\subsection{Ультрафильтры}

Теперь усилим понятие фильтра.

Фильтр говорит, какие множества считаются большими. Но обычный фильтр может не решить вопрос о конкретном множестве $A\subseteq I$: может оказаться, что ни $A$, ни его дополнение не являются большими.

Ультрафильтр такого не допускает.

\begin{definition}
Фильтр $\mathcal U$ на множестве $I$ называется \textbf{ультрафильтром}, если для любого подмножества
\[
A\subseteq I
\]
выполнено:
\[
A\in\mathcal U
\qquad
\text{или}
\qquad
I\setminus A\in\mathcal U.
\]
\end{definition}

Иными словами, ультрафильтр для каждого разбиения
\[
I=A\sqcup(I\setminus A)
\]
выбирает одну из двух частей как большую.

\begin{remark}
В ультрафильтре не могут одновременно лежать и множество $A$, и его дополнение $I\setminus A$.

Действительно, если бы
\[
A\in\mathcal U
\qquad
\text{и}
\qquad
I\setminus A\in\mathcal U,
\]
то по замкнутости фильтра относительно пересечений получили бы
\[
A\cap(I\setminus A)\in\mathcal U.
\]
Но
\[
A\cap(I\setminus A)=\varnothing,
\]
а пустое множество не принадлежит фильтру. Противоречие.
\end{remark}

\begin{example}
Кофинитный фильтр на $\omega$ не является ультрафильтром.

Действительно, пусть
\[
E=\{0,2,4,6,\ldots\}
\]
— множество чётных чисел. Тогда
\[
E\notin\operatorname{cof}(\omega),
\]
поскольку его дополнение, множество нечётных чисел, бесконечно.

Но и
\[
\omega\setminus E\notin\operatorname{cof}(\omega),
\]
поскольку его дополнение, множество чётных чисел, тоже бесконечно.

Значит, кофинитный фильтр не выбирает ни $E$, ни $\omega\setminus E$. Следовательно, он не является ультрафильтром.
\end{example}


\begin{definition}
Пусть $a\in I$. Семейство
\[
\mathcal U_a
=
\{A\subseteq I\mid a\in A\}
\]
называется \textbf{главным ультрафильтром}, порождённым точкой $a$.
\end{definition}

\begin{proposition}
Для любого $a\in I$ семейство
\[
\mathcal U_a
=
\{A\subseteq I\mid a\in A\}
\]
является ультрафильтром на $I$.
\end{proposition}

\begin{proof}
Сначала заметим, что $\mathcal U_a$ является фильтром.

Действительно,
\[
I\in\mathcal U_a,
\]
поскольку $a\in I$, и
\[
\varnothing\notin\mathcal U_a,
\]
поскольку $a\notin\varnothing$.

Если
\[
A\in\mathcal U_a
\qquad
\text{и}
\qquad
A\subseteq B\subseteq I,
\]
то
\[
a\in A\subseteq B,
\]
значит,
\[
a\in B.
\]
Следовательно,
\[
B\in\mathcal U_a.
\]

Если
\[
A,B\in\mathcal U_a,
\]
то
\[
a\in A
\qquad
\text{и}
\qquad
a\in B.
\]
Значит,
\[
a\in A\cap B,
\]
поэтому
\[
A\cap B\in\mathcal U_a.
\]

Теперь проверим свойство ультрафильтра. Пусть
\[
A\subseteq I.
\]
Тогда либо
\[
a\in A,
\]
и тогда
\[
A\in\mathcal U_a,
\]
либо
\[
a\notin A.
\]
Во втором случае
\[
a\in I\setminus A,
\]
значит,
\[
I\setminus A\in\mathcal U_a.
\]

Следовательно, $\mathcal U_a$ является ультрафильтром.
\end{proof}

\begin{remark}
Главный ультрафильтр в точке $a$ считает большим ровно те множества, которые содержат точку $a$.

Это самый простой пример ультрафильтра.
\end{remark}

\subsection{Максимальные фильтры}

\begin{definition}
Фильтр $\mathcal F$ на множестве $I$ называется \textbf{максимальным}, если его нельзя строго расширить до большего фильтра.

То есть если $\mathcal G$ — фильтр на $I$ и
\[
\mathcal F\subseteq\mathcal G,
\]
то обязательно
\[
\mathcal F=\mathcal G.
\]
\end{definition}

Интуитивно максимальный фильтр уже содержит все множества, которые можно добавить, не разрушив свойств фильтра.

\begin{theorem}
Любой максимальный фильтр является ультрафильтром.
\end{theorem}

\begin{proof}
Пусть $\mathcal F$ — максимальный фильтр на $I$. Докажем, что $\mathcal F$ является ультрафильтром.

Нужно показать, что для любого подмножества
\[
A\subseteq I
\]
выполнено:
\[
A\in\mathcal F
\qquad
\text{или}
\qquad
I\setminus A\in\mathcal F.
\]

Предположим противное. Пусть существует множество $A\subseteq I$, такое что
\[
A\notin\mathcal F
\qquad
\text{и}
\qquad
I\setminus A\notin\mathcal F.
\]

Попробуем добавить $A$ к фильтру $\mathcal F$ и построить новый фильтр, строго содержащий $\mathcal F$. Это даст противоречие с максимальностью $\mathcal F$.

Определим семейство
\[
\mathcal G
=
\{C\subseteq I\mid \exists B\in\mathcal F \text{ такое, что } A\cap B\subseteq C\}.
\]

Интуитивно $\mathcal G$ состоит из всех множеств, которые содержат пересечение $A$ с некоторым большим множеством $B\in\mathcal F$.

Сначала покажем, что
\[
\mathcal F\subseteq\mathcal G.
\]
Пусть
\[
C\in\mathcal F.
\]
Возьмём
\[
B=C.
\]
Тогда
\[
B\in\mathcal F
\]
и
\[
A\cap B=A\cap C\subseteq C.
\]
Следовательно,
\[
C\in\mathcal G.
\]
Значит,
\[
\mathcal F\subseteq\mathcal G.
\]

Кроме того,
\[
A\in\mathcal G.
\]
Действительно, можно взять
\[
B=I.
\]
Так как $I\in\mathcal F$, имеем
\[
A\cap I=A\subseteq A.
\]
Поэтому
\[
A\in\mathcal G.
\]

Но по предположению
\[
A\notin\mathcal F.
\]
Следовательно,
\[
\mathcal F\subsetneq\mathcal G.
\]

Осталось показать, что $\mathcal G$ является фильтром.

Во-первых,
\[
I\in\mathcal G,
\]
поскольку $\mathcal F\subseteq\mathcal G$ и $I\in\mathcal F$.

Покажем, что
\[
\varnothing\notin\mathcal G.
\]
Предположим противное:
\[
\varnothing\in\mathcal G.
\]
Тогда по определению $\mathcal G$ существует $B\in\mathcal F$, такое что
\[
A\cap B\subseteq\varnothing.
\]
Значит,
\[
A\cap B=\varnothing.
\]
Отсюда
\[
B\subseteq I\setminus A.
\]
Так как
\[
B\in\mathcal F
\]
и фильтр замкнут вверх, получаем
\[
I\setminus A\in\mathcal F.
\]
Но это противоречит предположению
\[
I\setminus A\notin\mathcal F.
\]
Следовательно,
\[
\varnothing\notin\mathcal G.
\]

Теперь проверим замкнутость вверх. Пусть
\[
C\in\mathcal G
\qquad
\text{и}
\qquad
C\subseteq D\subseteq I.
\]
По определению $\mathcal G$ существует $B\in\mathcal F$, такое что
\[
A\cap B\subseteq C.
\]
Так как $C\subseteq D$, получаем
\[
A\cap B\subseteq D.
\]
Следовательно,
\[
D\in\mathcal G.
\]

Теперь проверим замкнутость относительно пересечений. Пусть
\[
C_1,C_2\in\mathcal G.
\]
Тогда существуют
\[
B_1,B_2\in\mathcal F
\]
такие, что
\[
A\cap B_1\subseteq C_1
\]
и
\[
A\cap B_2\subseteq C_2.
\]

Так как $\mathcal F$ — фильтр, то
\[
B_1\cap B_2\in\mathcal F.
\]
Кроме того,
\[
A\cap(B_1\cap B_2)
=
(A\cap B_1)\cap(A\cap B_2)
\subseteq C_1\cap C_2.
\]
Следовательно,
\[
C_1\cap C_2\in\mathcal G.
\]

Итак, $\mathcal G$ является фильтром на $I$, причём
\[
\mathcal F\subsetneq\mathcal G.
\]
Это противоречит максимальности $\mathcal F$.

Значит, наше предположение было неверным. Следовательно, для любого $A\subseteq I$
\[
A\in\mathcal F
\qquad
\text{или}
\qquad
I\setminus A\in\mathcal F.
\]

То есть $\mathcal F$ является ультрафильтром.
\end{proof}

\begin{remark}
Содержательно доказательство говорит следующее.

Если фильтр $\mathcal F$ не является ультрафильтром, то существует множество $A\subseteq I$, для которого фильтр не выбрал ни $A$, ни его дополнение. Тогда к фильтру можно добавить $A$ и породить новый фильтр, больший исходного.

Следовательно, максимальный фильтр не может ``сомневаться'' ни по одному множеству $A$. Он обязан для каждого $A\subseteq I$ выбрать либо $A$, либо $I\setminus A$.
\end{remark}

\begin{remark}
Итак, имеем следующую картину:
\[
\text{фильтр}
\quad
=
\quad
\text{система больших множеств},
\]
\[
\text{ультрафильтр}
\quad
=
\quad
\begin{gathered}
\text{максимальная система больших множеств,}\\
\text{решающая каждое разбиение}.
\end{gathered}
\]

Для любого $A\subseteq I$ ультрафильтр выбирает ровно одну сторону:
\[
A
\qquad
\text{или}
\qquad
I\setminus A.
\]
\end{remark}

\begin{lemma}[Критерий ультрафильтра через максимальность]
Пусть $I$ — непустое множество, и пусть $\mathcal F$ — фильтр на $I$. Тогда следующие условия эквивалентны:
\begin{enumerate}
    \item $\mathcal F$ является ультрафильтром;
    \item $\mathcal F$ является максимальным по включению фильтром на $I$.
\end{enumerate}
\end{lemma}

\begin{proof}
Докажем оба направления.

Сначала докажем, что любой ультрафильтр максимален.

Пусть $\mathcal U$ — ультрафильтр на $I$. Предположим, что $\mathcal G$ — фильтр на $I$, такой что
\[
\mathcal U\subseteq \mathcal G.
\]
Покажем, что тогда
\[
\mathcal G=\mathcal U.
\]

Предположим противное. Тогда существует множество
\[
A\in\mathcal G
\]
такое, что
\[
A\notin\mathcal U.
\]

Так как $\mathcal U$ — ультрафильтр, для любого подмножества $A\subseteq I$ выполнено:
\[
A\in\mathcal U
\qquad
\text{или}
\qquad
I\setminus A\in\mathcal U.
\]
Поскольку
\[
A\notin\mathcal U,
\]
получаем
\[
I\setminus A\in\mathcal U.
\]
Но
\[
\mathcal U\subseteq\mathcal G,
\]
поэтому
\[
I\setminus A\in\mathcal G.
\]

С другой стороны,
\[
A\in\mathcal G.
\]
Так как $\mathcal G$ — фильтр, он замкнут относительно пересечений. Поэтому
\[
A\cap(I\setminus A)\in\mathcal G.
\]
Но
\[
A\cap(I\setminus A)=\varnothing.
\]
Значит,
\[
\varnothing\in\mathcal G,
\]
что невозможно для фильтра.

Получили противоречие. Следовательно, не существует множества $A\in\mathcal G\setminus\mathcal U$, а значит,
\[
\mathcal G=\mathcal U.
\]
Итак, $\mathcal U$ максимален по включению среди фильтров на $I$.

Теперь докажем обратное направление: любой максимальный фильтр является ультрафильтром.

Пусть $\mathcal F$ — максимальный по включению фильтр на $I$. Докажем, что $\mathcal F$ является ультрафильтром.

Предположим противное. Тогда существует множество $A\subseteq I$, такое что
\[
A\notin\mathcal F
\qquad
\text{и}
\qquad
I\setminus A\notin\mathcal F.
\]

Построим фильтр, который строго содержит $\mathcal F$. Для этого определим семейство
\[
\mathcal G
=
\{C\subseteq I\mid \exists B\in\mathcal F\text{ такое, что } A\cap B\subseteq C\}.
\]

Интуитивно $\mathcal G$ — это фильтр, полученный из $\mathcal F$ добавлением множества $A$ и всех множеств, которые обязаны появиться после этого.

Сначала покажем, что
\[
\mathcal F\subseteq\mathcal G.
\]
Пусть
\[
C\in\mathcal F.
\]
Возьмём
\[
B=C.
\]
Тогда $B\in\mathcal F$ и
\[
A\cap B=A\cap C\subseteq C.
\]
Следовательно,
\[
C\in\mathcal G.
\]
Значит,
\[
\mathcal F\subseteq\mathcal G.
\]

Кроме того,
\[
A\in\mathcal G.
\]
Действительно, можно взять
\[
B=I.
\]
Так как $I\in\mathcal F$, имеем
\[
A\cap I=A\subseteq A.
\]
Поэтому
\[
A\in\mathcal G.
\]
Но по предположению
\[
A\notin\mathcal F.
\]
Следовательно,
\[
\mathcal F\subsetneq\mathcal G.
\]

Осталось проверить, что $\mathcal G$ действительно является фильтром.

Во-первых,
\[
I\in\mathcal G,
\]
поскольку $\mathcal F\subseteq\mathcal G$ и $I\in\mathcal F$.

Покажем, что
\[
\varnothing\notin\mathcal G.
\]
Предположим противное:
\[
\varnothing\in\mathcal G.
\]
Тогда по определению $\mathcal G$ существует $B\in\mathcal F$, такое что
\[
A\cap B\subseteq\varnothing.
\]
Значит,
\[
A\cap B=\varnothing.
\]
Отсюда следует, что
\[
B\subseteq I\setminus A.
\]
Так как $B\in\mathcal F$ и фильтр замкнут вверх, получаем
\[
I\setminus A\in\mathcal F.
\]
Но это противоречит предположению
\[
I\setminus A\notin\mathcal F.
\]
Следовательно,
\[
\varnothing\notin\mathcal G.
\]

Теперь проверим замкнутость вверх. Пусть
\[
C\in\mathcal G
\qquad
\text{и}
\qquad
C\subseteq D\subseteq I.
\]
По определению $\mathcal G$ существует $B\in\mathcal F$, такое что
\[
A\cap B\subseteq C.
\]
Так как
\[
C\subseteq D,
\]
получаем
\[
A\cap B\subseteq D.
\]
Следовательно,
\[
D\in\mathcal G.
\]

Осталось проверить замкнутость относительно пересечений. Пусть
\[
C_1,C_2\in\mathcal G.
\]
Тогда существуют
\[
B_1,B_2\in\mathcal F
\]
такие, что
\[
A\cap B_1\subseteq C_1
\]
и
\[
A\cap B_2\subseteq C_2.
\]
Так как $\mathcal F$ — фильтр, то
\[
B_1\cap B_2\in\mathcal F.
\]
Кроме того,
\[
A\cap(B_1\cap B_2)
=
(A\cap B_1)\cap(A\cap B_2)
\subseteq C_1\cap C_2.
\]
Следовательно,
\[
C_1\cap C_2\in\mathcal G.
\]

Итак, $\mathcal G$ является фильтром на $I$ и
\[
\mathcal F\subsetneq\mathcal G.
\]
Это противоречит максимальности $\mathcal F$.

Следовательно, наше предположение было неверным. Значит, для любого $A\subseteq I$ выполнено
\[
A\in\mathcal F
\qquad
\text{или}
\qquad
I\setminus A\in\mathcal F.
\]
То есть $\mathcal F$ является ультрафильтром.
\end{proof}

\begin{remark}
Таким образом, ультрафильтр можно понимать как максимальную систему больших множеств.

Если фильтр не является ультрафильтром, то существует множество $A\subseteq I$, по которому фильтр ещё не сделал выбор: он не содержит ни $A$, ни $I\setminus A$. Тогда к фильтру можно добавить $A$ и получить фильтр большего размера.

Максимальный фильтр такой возможности не допускает, поэтому он обязан быть ультрафильтром.
\end{remark}

\begin{remark}[Лемма Цорна]
В следующих рассуждениях используется лемма Цорна.

Она утверждает: если в частично упорядоченном множестве всякая цепь имеет верхнюю грань, то в этом множестве существует максимальный элемент.

Здесь цепью называется такое подмножество частично упорядоченного множества, в котором любые два элемента сравнимы.
\end{remark}

\begin{lemma}[Лемма о расширении фильтра до ультрафильтра]
Пусть $I$ — непустое множество, и пусть $\mathcal F$ — фильтр на $I$. Тогда существует ультрафильтр $\mathcal U$ на $I$, такой что
\[
\mathcal F\subseteq\mathcal U.
\]
\end{lemma}

\begin{proof}
Рассмотрим множество всех фильтров на $I$, содержащих $\mathcal F$:
\[
\mathfrak P
=
\{\mathcal G\subseteq\mathcal P(I)\mid
\mathcal G\text{ — фильтр на }I
\text{ и }
\mathcal F\subseteq\mathcal G\}.
\]
Это множество упорядочено по включению.

Множество $\mathfrak P$ непусто, поскольку
\[
\mathcal F\in\mathfrak P.
\]

Применим лемму Цорна. Для этого нужно показать, что всякая цепь в $\mathfrak P$ имеет верхнюю грань в $\mathfrak P$.

Пусть
\[
\{\mathcal G_\alpha\}_{\alpha\in J}
\]
— цепь в $\mathfrak P$. Это означает, что для любых $\alpha,\beta\in J$ выполнено одно из двух включений:
\[
\mathcal G_\alpha\subseteq\mathcal G_\beta
\qquad
\text{или}
\qquad
\mathcal G_\beta\subseteq\mathcal G_\alpha.
\]

Положим
\[
\mathcal G
=
\bigcup_{\alpha\in J}\mathcal G_\alpha.
\]
Покажем, что $\mathcal G$ является фильтром на $I$.

Во-первых,
\[
I\in\mathcal G,
\]
поскольку $I$ лежит в каждом фильтре $\mathcal G_\alpha$.

Также
\[
\varnothing\notin\mathcal G,
\]
поскольку $\varnothing$ не лежит ни в одном фильтре $\mathcal G_\alpha$.

Проверим замкнутость вверх. Пусть
\[
A\in\mathcal G
\qquad
\text{и}
\qquad
A\subseteq B\subseteq I.
\]
Так как $A\in\mathcal G$, существует $\alpha\in J$, такое что
\[
A\in\mathcal G_\alpha.
\]
Но $\mathcal G_\alpha$ — фильтр, поэтому из $A\subseteq B$ следует
\[
B\in\mathcal G_\alpha.
\]
Значит,
\[
B\in\mathcal G.
\]

Проверим замкнутость относительно пересечений. Пусть
\[
A,B\in\mathcal G.
\]
Тогда существуют $\alpha,\beta\in J$, такие что
\[
A\in\mathcal G_\alpha,
\qquad
B\in\mathcal G_\beta.
\]
Так как $\{\mathcal G_\alpha\}_{\alpha\in J}$ — цепь, фильтры $\mathcal G_\alpha$ и $\mathcal G_\beta$ сравнимы по включению.

Например, пусть
\[
\mathcal G_\alpha\subseteq\mathcal G_\beta.
\]
Тогда
\[
A\in\mathcal G_\beta
\qquad
\text{и}
\qquad
B\in\mathcal G_\beta.
\]
Поскольку $\mathcal G_\beta$ — фильтр, получаем
\[
A\cap B\in\mathcal G_\beta.
\]
Следовательно,
\[
A\cap B\in\mathcal G.
\]

Если же
\[
\mathcal G_\beta\subseteq\mathcal G_\alpha,
\]
то рассуждение аналогично.

Значит, $\mathcal G$ является фильтром.

Кроме того, для любого $\alpha\in J$ имеем
\[
\mathcal G_\alpha\subseteq\mathcal G,
\]
и также
\[
\mathcal F\subseteq\mathcal G,
\]
поскольку $\mathcal F\subseteq\mathcal G_\alpha$ для каждого $\alpha\in J$.

Следовательно,
\[
\mathcal G\in\mathfrak P
\]
и $\mathcal G$ является верхней гранью данной цепи.

По лемме Цорна в $\mathfrak P$ существует максимальный элемент. Обозначим его через
\[
\mathcal U.
\]
Тогда $\mathcal U$ — максимальный фильтр на $I$, содержащий $\mathcal F$.

По критерию ультрафильтра через максимальность, доказанному выше, $\mathcal U$ является ультрафильтром.

Следовательно, существует ультрафильтр $\mathcal U$ на $I$, такой что
\[
\mathcal F\subseteq\mathcal U.
\]
\end{proof}

\begin{lemma}[Существование неглавного ультрафильтра на бесконечном множестве]
Пусть $I$ — бесконечное множество. Тогда на $I$ существует неглавный ультрафильтр.
\end{lemma}

\begin{proof}
Рассмотрим кофинитный фильтр на $I$:
\[
\operatorname{cof}(I)
=
\{A\subseteq I\mid I\setminus A\text{ конечно}\}.
\]
Так как $I$ бесконечно, это действительно фильтр на $I$.

По лемме о расширении фильтра до ультрафильтра существует ультрафильтр $\mathcal U$ на $I$, такой что
\[
\operatorname{cof}(I)\subseteq\mathcal U.
\]

Покажем, что $\mathcal U$ неглавный.

Предположим противное. Тогда $\mathcal U$ является главным ультрафильтром. Значит, существует точка $a\in I$, такая что
\[
\mathcal U=\mathcal U_a,
\]
где
\[
\mathcal U_a=\{A\subseteq I\mid a\in A\}.
\]
В частности,
\[
\{a\}\in\mathcal U.
\]

Но так как $I$ бесконечно, множество
\[
I\setminus\{a\}
\]
кофинитно. Следовательно,
\[
I\setminus\{a\}\in\operatorname{cof}(I).
\]
А поскольку
\[
\operatorname{cof}(I)\subseteq\mathcal U,
\]
получаем
\[
I\setminus\{a\}\in\mathcal U.
\]

Теперь в ультрафильтре $\mathcal U$ лежат оба множества:
\[
\{a\}\in\mathcal U
\qquad
\text{и}
\qquad
I\setminus\{a\}\in\mathcal U.
\]
Так как $\mathcal U$ — фильтр, он замкнут относительно пересечений. Поэтому
\[
\{a\}\cap(I\setminus\{a\})\in\mathcal U.
\]
Но
\[
\{a\}\cap(I\setminus\{a\})=\varnothing.
\]
Получаем
\[
\varnothing\in\mathcal U,
\]
что невозможно для фильтра.

Противоречие. Следовательно, $\mathcal U$ не является главным ультрафильтром.

Значит, на любом бесконечном множестве существует неглавный ультрафильтр.
\end{proof}

\begin{remark}
Существование неглавного ультрафильтра на бесконечном множестве доказывается с использованием леммы Цорна, то есть некоторой формы аксиомы выбора.

Поэтому такой ультрафильтр обычно не задаётся явно. Мы доказываем его существование, но не строим конкретный список всех множеств, которые в него входят.
\end{remark}

\section{Ультрапроизведения}

\subsection{Идея ультрапроизведения}

Пусть $L$ — язык первого порядка. Пусть
\[
\{\mathcal M_i\}_{i\in I}
\]
— семейство $L$-структур, индексированное непустым множеством $I$.

Также пусть
\[
\mathcal U\subseteq\mathcal P(I)
\]
— ультрафильтр на $I$.

Ультрапроизведение — это способ построить из семейства структур
\[
\mathcal M_i,\qquad i\in I,
\]
одну новую $L$-структуру. Интуитивно она является некоторым ``усреднением'' структур $\mathcal M_i$ по ультрафильтру $\mathcal U$.

Ультрафильтр в этой конструкции отвечает за то, какие множества индексов считаются большими. Поэтому в ультрапроизведении утверждение считается истинным, если оно истинно в структурах $\mathcal M_i$ для большого множества индексов $i$.

\begin{remark}
Главная идея ультрапроизведения:
\[
\text{истинно в ультрапроизведении}
\quad
\Longleftrightarrow
\quad
\text{истинно на большом множестве координат}.
\]
Точное утверждение такого вида называется теоремой Лося.
\end{remark}

\subsection{Прямое произведение носителей}

Сначала рассмотрим обычное декартово произведение носителей:
\[
\prod_{i\in I} M_i.
\]

Элемент этого произведения — это функция
\[
f:I\to \bigcup_{i\in I}M_i
\]
такая, что
\[
f(i)\in M_i
\]
для каждого $i\in I$.

Иными словами, элемент произведения выбирает по одному элементу из каждой структуры:
\[
f=(f(i))_{i\in I}.
\]

Если $I=\omega$, то элемент произведения можно понимать как последовательность:
\[
(f(0),f(1),f(2),\ldots),
\]
где
\[
f(n)\in M_n.
\]

\begin{example}
Если для всех $i\in\omega$ выполнено
\[
M_i=\mathbb R,
\]
то
\[
\prod_{i\in\omega} M_i
=
\prod_{i\in\omega}\mathbb R
\]
есть множество всех вещественных последовательностей.
\end{example}

\subsection{Отношение эквивалентности по ультрафильтру}

На произведении
\[
\prod_{i\in I} M_i
\]
зададим отношение $\sim_{\mathcal U}$ следующим образом.

\begin{definition}
Для
\[
f,g\in\prod_{i\in I}M_i
\]
положим
\[
f\sim_{\mathcal U} g
\]
тогда и только тогда, когда
\[
\{i\in I\mid f(i)=g(i)\}\in\mathcal U.
\]
\end{definition}

То есть две функции считаются эквивалентными, если они совпадают на большом множестве индексов.

Если ультрафильтр $\mathcal U$ главный и сосредоточен в точке $i_0\in I$, то
\[
f\sim_{\mathcal U}g
\]
означает просто
\[
f(i_0)=g(i_0).
\]

Если же $\mathcal U$ — неглавный ультрафильтр на $\omega$, содержащий кофинитный фильтр, то любые две последовательности, совпадающие начиная с некоторого номера, будут эквивалентны. Однако неглавный ультрафильтр может считать большими и более сложные бесконечные множества индексов.

\begin{lemma}
Отношение $\sim_{\mathcal U}$ является отношением эквивалентности на множестве
\[
\prod_{i\in I}M_i.
\]
\end{lemma}

\begin{proof}
Проверим рефлексивность, симметричность и транзитивность.

\textbf{Рефлексивность.}
Для любой функции
\[
f\in\prod_{i\in I}M_i
\]
имеем
\[
\{i\in I\mid f(i)=f(i)\}=I.
\]
Так как $\mathcal U$ — фильтр, то
\[
I\in\mathcal U.
\]
Следовательно,
\[
f\sim_{\mathcal U} f.
\]

\textbf{Симметричность.}
Пусть
\[
f\sim_{\mathcal U} g.
\]
Тогда
\[
\{i\in I\mid f(i)=g(i)\}\in\mathcal U.
\]
Но равенство симметрично, поэтому
\[
\{i\in I\mid f(i)=g(i)\}
=
\{i\in I\mid g(i)=f(i)\}.
\]
Следовательно,
\[
\{i\in I\mid g(i)=f(i)\}\in\mathcal U,
\]
то есть
\[
g\sim_{\mathcal U} f.
\]

\textbf{Транзитивность.}
Пусть
\[
f\sim_{\mathcal U} g
\qquad
\text{и}
\qquad
g\sim_{\mathcal U} h.
\]
Тогда множества
\[
A=\{i\in I\mid f(i)=g(i)\}
\]
и
\[
B=\{i\in I\mid g(i)=h(i)\}
\]
принадлежат ультрафильтру $\mathcal U$.

Так как $\mathcal U$ — фильтр, он замкнут относительно пересечений. Поэтому
\[
A\cap B\in\mathcal U.
\]

Для любого $i\in A\cap B$ имеем
\[
f(i)=g(i)
\qquad
\text{и}
\qquad
g(i)=h(i).
\]
Следовательно,
\[
f(i)=h(i).
\]
Значит,
\[
A\cap B\subseteq \{i\in I\mid f(i)=h(i)\}.
\]
Поскольку $\mathcal U$ замкнут вверх, получаем
\[
\{i\in I\mid f(i)=h(i)\}\in\mathcal U.
\]
Следовательно,
\[
f\sim_{\mathcal U} h.
\]

Итак, $\sim_{\mathcal U}$ является отношением эквивалентности.
\end{proof}

\subsection{Носитель ультрапроизведения}

Обозначим через
\[
[f]
\]
класс эквивалентности функции $f$ относительно отношения $\sim_{\mathcal U}$.

То есть
\[
[f]
=
\left\{
g\in\prod_{i\in I}M_i
\mid
g\sim_{\mathcal U} f
\right\}.
\]

\begin{definition}
Носителем ультрапроизведения семейства структур $\{\mathcal M_i\}_{i\in I}$ по ультрафильтру $\mathcal U$ называется множество классов эквивалентности
\[
M
=
\prod_{i\in I}M_i/\mathcal U
=
\left\{[f]\mid f\in\prod_{i\in I}M_i\right\}.
\]
\end{definition}

Будущая структура на этом носителе обозначается
\[
\prod_{i\in I}\mathcal M_i/\mathcal U.
\]

Элементы ультрапроизведения — это не сами функции $f$, а их классы эквивалентности $[f]$. Поэтому все операции и отношения нужно определить так, чтобы они не зависели от выбора представителя класса.

\subsection{Интерпретация констант}

Пусть
\[
c\in L
\]
— константный символ.

В каждой структуре $\mathcal M_i$ этот символ имеет интерпретацию
\[
c^{\mathcal M_i}\in M_i.
\]

Определим функцию
\[
f_c\in\prod_{i\in I}M_i
\]
по правилу
\[
f_c(i)=c^{\mathcal M_i}.
\]

\begin{definition}
Интерпретация константного символа $c$ в ультрапроизведении задаётся формулой
\[
c^{\prod_{i\in I}\mathcal M_i/\mathcal U}
=
[f_c],
\]
где
\[
f_c(i)=c^{\mathcal M_i}.
\]
\end{definition}

То есть константа в ультрапроизведении — это класс функции, которая в каждой координате берёт интерпретацию этой константы в соответствующей структуре.

\subsection{Интерпретация функциональных символов и её корректность}

Пусть
\[
F\in L
\]
— $n$-местный функциональный символ.

В каждой структуре $\mathcal M_i$ он интерпретируется как функция
\[
F^{\mathcal M_i}:M_i^n\to M_i.
\]

Пусть
\[
[f_1],\ldots,[f_n]\in \prod_{i\in I}M_i/\mathcal U.
\]

Определим функцию
\[
h\in\prod_{i\in I}M_i
\]
покоординатно:
\[
h(i)
=
F^{\mathcal M_i}(f_1(i),\ldots,f_n(i)).
\]

\begin{definition}
Интерпретация функционального символа $F$ в ультрапроизведении задаётся формулой
\[
F^{\prod_{i\in I}\mathcal M_i/\mathcal U}
([f_1],\ldots,[f_n])
=
[h],
\]
где
\[
h(i)
=
F^{\mathcal M_i}(f_1(i),\ldots,f_n(i)).
\]
\end{definition}

Иными словами, функции в ультрапроизведении вычисляются покоординатно, а затем берётся класс эквивалентности результата.

\subsubsection{Корректность интерпретации функциональных символов}

Нужно проверить, что определение функционального символа не зависит от выбора представителей классов.

\begin{lemma}
Пусть $F$ — $n$-местный функциональный символ языка $L$. Если
\[
f_k\sim_{\mathcal U} g_k
\]
для всех
\[
k=1,\ldots,n,
\]
то
\[
F^{\prod_{i\in I}\mathcal M_i/\mathcal U}
([f_1],\ldots,[f_n])
=
F^{\prod_{i\in I}\mathcal M_i/\mathcal U}
([g_1],\ldots,[g_n]).
\]
\end{lemma}

\begin{proof}
Для каждого $k=1,\ldots,n$ положим
\[
A_k
=
\{i\in I\mid f_k(i)=g_k(i)\}.
\]
Так как
\[
f_k\sim_{\mathcal U} g_k,
\]
имеем
\[
A_k\in\mathcal U.
\]

Поскольку $\mathcal U$ — фильтр, пересечение конечного числа множеств из $\mathcal U$ снова лежит в $\mathcal U$. Поэтому
\[
A=A_1\cap\cdots\cap A_n\in\mathcal U.
\]

Для любого $i\in A$ выполнено
\[
f_1(i)=g_1(i),\ldots,f_n(i)=g_n(i).
\]
Следовательно,
\[
F^{\mathcal M_i}(f_1(i),\ldots,f_n(i))
=
F^{\mathcal M_i}(g_1(i),\ldots,g_n(i)).
\]

Пусть
\[
h_f(i)=F^{\mathcal M_i}(f_1(i),\ldots,f_n(i)),
\]
а
\[
h_g(i)=F^{\mathcal M_i}(g_1(i),\ldots,g_n(i)).
\]
Тогда для всех $i\in A$ имеем
\[
h_f(i)=h_g(i).
\]
Значит,
\[
A\subseteq\{i\in I\mid h_f(i)=h_g(i)\}.
\]
Так как $A\in\mathcal U$ и $\mathcal U$ замкнут вверх, получаем
\[
\{i\in I\mid h_f(i)=h_g(i)\}\in\mathcal U.
\]
Следовательно,
\[
h_f\sim_{\mathcal U}h_g,
\]
то есть
\[
[h_f]=[h_g].
\]

Значит, результат применения $F$ не зависит от выбора представителей классов.
\end{proof}

\subsection{Интерпретация предикатных символов и её корректность}

Пусть
\[
R\in L
\]
— $n$-местный предикатный символ.

В каждой структуре $\mathcal M_i$ он интерпретируется как отношение
\[
R^{\mathcal M_i}\subseteq M_i^n.
\]

\begin{definition}
В ультрапроизведении отношение $R$ определяется следующим образом:
\[
R^{\prod_{i\in I}\mathcal M_i/\mathcal U}
([f_1],\ldots,[f_n])
\]
выполнено тогда и только тогда, когда
\[
\{i\in I\mid
\mathcal M_i\models R(f_1(i),\ldots,f_n(i))
\}\in\mathcal U.
\]
\end{definition}

Иными словами, предикат $R$ истиннен на классах
\[
[f_1],\ldots,[f_n]
\]
в ультрапроизведении тогда и только тогда, когда он истиннен в координатах на большом множестве индексов.

Эквивалентно можно записать:
\[
R^{\prod_{i\in I}\mathcal M_i/\mathcal U}
=
\left\{
([f_1],\ldots,[f_n])
\ \middle|\
\{i\in I\mid
(f_1(i),\ldots,f_n(i))\in R^{\mathcal M_i}
\}\in\mathcal U
\right\}.
\]

\subsection{Корректность интерпретации предикатных символов}

Так как элементы ультрапроизведения являются классами эквивалентности, нужно проверить, что истинность предиката не зависит от выбора представителей.

\begin{lemma}
Пусть $R$ — $n$-местный предикатный символ языка $L$. Если
\[
f_k\sim_{\mathcal U} g_k
\]
для всех
\[
k=1,\ldots,n,
\]
то
\[
R^{\prod_{i\in I}\mathcal M_i/\mathcal U}
([f_1],\ldots,[f_n])
\]
выполнено тогда и только тогда, когда
\[
R^{\prod_{i\in I}\mathcal M_i/\mathcal U}
([g_1],\ldots,[g_n])
\]
выполнено.
\end{lemma}

\begin{proof}
Для каждого $k=1,\ldots,n$ положим
\[
A_k
=
\{i\in I\mid f_k(i)=g_k(i)\}.
\]
Из условия
\[
f_k\sim_{\mathcal U} g_k
\]
следует, что
\[
A_k\in\mathcal U.
\]

Пусть
\[
A=A_1\cap\cdots\cap A_n.
\]
Так как $\mathcal U$ — фильтр, то
\[
A\in\mathcal U.
\]

Для любого $i\in A$ имеем
\[
(f_1(i),\ldots,f_n(i))
=
(g_1(i),\ldots,g_n(i)).
\]
Следовательно, для любого $i\in A$ выполнена эквивалентность
\[
\mathcal M_i\models R(f_1(i),\ldots,f_n(i))
\]
тогда и только тогда, когда
\[
\mathcal M_i\models R(g_1(i),\ldots,g_n(i)).
\]

Обозначим
\[
B_f
=
\{i\in I\mid
\mathcal M_i\models R(f_1(i),\ldots,f_n(i))
\}
\]
и
\[
B_g
=
\{i\in I\mid
\mathcal M_i\models R(g_1(i),\ldots,g_n(i))
\}.
\]

Мы показали, что на множестве $A$ множества $B_f$ и $B_g$ совпадают:
\[
A\cap B_f=A\cap B_g.
\]

Нужно доказать, что
\[
B_f\in\mathcal U
\qquad
\Longleftrightarrow
\qquad
B_g\in\mathcal U.
\]

Пусть сначала
\[
B_f\in\mathcal U.
\]
Тогда
\[
A\cap B_f\in\mathcal U.
\]
Но
\[
A\cap B_f=A\cap B_g
\subseteq B_g.
\]
Так как $\mathcal U$ замкнут вверх, получаем
\[
B_g\in\mathcal U.
\]

Обратно, если
\[
B_g\in\mathcal U,
\]
то аналогично
\[
A\cap B_g\in\mathcal U.
\]
Но
\[
A\cap B_g=A\cap B_f
\subseteq B_f.
\]
Следовательно,
\[
B_f\in\mathcal U.
\]

Итак,
\[
B_f\in\mathcal U
\qquad
\Longleftrightarrow
\qquad
B_g\in\mathcal U.
\]
Значит, истинность предиката $R$ в ультрапроизведении не зависит от выбора представителей классов.
\end{proof}

\subsection{Итоговое определение ультрапроизведения}

\begin{definition}
Пусть $L$ — язык первого порядка, пусть
\[
\{\mathcal M_i\}_{i\in I}
\]
— семейство $L$-структур, и пусть $\mathcal U$ — ультрафильтр на $I$.

\textbf{Ультрапроизведением} структур $\mathcal M_i$ по ультрафильтру $\mathcal U$ называется $L$-структура
\[
\prod_{i\in I}\mathcal M_i/\mathcal U,
\]
построенная следующим образом.

Её носитель:
\[
\prod_{i\in I}M_i/\mathcal U
=
\left\{[f]\mid f\in\prod_{i\in I}M_i\right\},
\]
где
\[
f\sim_{\mathcal U} g
\]
тогда и только тогда, когда
\[
\{i\in I\mid f(i)=g(i)\}\in\mathcal U.
\]

Константы интерпретируются покоординатно:
\[
c^{\prod_{i\in I}\mathcal M_i/\mathcal U}
=
[i\mapsto c^{\mathcal M_i}].
\]

Функциональные символы интерпретируются покоординатно:
\[
F^{\prod_{i\in I}\mathcal M_i/\mathcal U}
([f_1],\ldots,[f_n])
=
\left[
i\mapsto
F^{\mathcal M_i}(f_1(i),\ldots,f_n(i))
\right].
\]

Предикатные символы интерпретируются через ультрафильтр:
\[
\prod_{i\in I}\mathcal M_i/\mathcal U
\models
R([f_1],\ldots,[f_n])
\]
тогда и только тогда, когда
\[
\{i\in I\mid
\mathcal M_i\models R(f_1(i),\ldots,f_n(i))
\}\in\mathcal U.
\]
\end{definition}

\begin{remark}
Если все структуры $\mathcal M_i$ равны одной и той же структуре $\mathcal M$, то ультрапроизведение
\[
\prod_{i\in I}\mathcal M/\mathcal U
\]
называется \textbf{ультрастепенью} структуры $\mathcal M$ по ультрафильтру $\mathcal U$.
\end{remark}

\begin{remark}
Ультрапроизведение можно понимать как структуру, в которой элементы являются обобщёнными последовательностями элементов исходных структур, а равенство и истинность отношений проверяются ``почти всюду'' относительно ультрафильтра.

Позже это будет выражено точно в теореме Лося:
\[
\prod_{i\in I}\mathcal M_i/\mathcal U
\models
\varphi([f_1],\ldots,[f_n])
\]
тогда и только тогда, когда
\[
\{i\in I\mid
\mathcal M_i\models
\varphi(f_1(i),\ldots,f_n(i))
\}\in\mathcal U.
\]
\end{remark}

Перед доказательством теоремы Лося нужно зафиксировать, как ведут себя термы.

\begin{lemma}[Покоординатное вычисление термов]
Пусть $t(x_1,\ldots,x_n)$ — $L$-терм. Пусть
\[
f_1,\ldots,f_n\in\prod_{i\in I}M_i.
\]
Определим функцию
\[
g_t\in\prod_{i\in I}M_i
\]
по правилу
\[
g_t(i)=
t^{\mathcal M_i}(f_1(i),\ldots,f_n(i)).
\]
Тогда
\[
t^{\mathcal M}([f_1],\ldots,[f_n])
=
[g_t].
\]
\end{lemma}

\begin{proof}
Докажем утверждение индукцией по построению терма $t$.

Если $t$ является переменной $x_j$, то
\[
t^{\mathcal M}([f_1],\ldots,[f_n])
=
[f_j].
\]
С другой стороны,
\[
g_t(i)=f_j(i).
\]
Значит,
\[
[g_t]=[f_j].
\]
Утверждение верно.

Если $t$ является константным символом $c$, то по определению интерпретации константы в ультрапроизведении
\[
c^{\mathcal M}
=
[i\mapsto c^{\mathcal M_i}].
\]
Но именно такая функция и есть $g_t$. Значит,
\[
t^{\mathcal M}=[g_t].
\]

Пусть теперь терм имеет вид
\[
t=F(t_1,\ldots,t_m),
\]
где $F$ — $m$-местный функциональный символ, а $t_1,\ldots,t_m$ — более простые термы.

По индукционному предположению для каждого $j=1,\ldots,m$ имеем
\[
t_j^{\mathcal M}([f_1],\ldots,[f_n])
=
[g_j],
\]
где
\[
g_j(i)=
t_j^{\mathcal M_i}(f_1(i),\ldots,f_n(i)).
\]

Тогда по определению интерпретации функционального символа в ультрапроизведении
\[
t^{\mathcal M}([f_1],\ldots,[f_n])
=
F^{\mathcal M}([g_1],\ldots,[g_m])
=
[h],
\]
где
\[
h(i)=
F^{\mathcal M_i}(g_1(i),\ldots,g_m(i)).
\]
Подставляя определения $g_1,\ldots,g_m$, получаем
\[
h(i)
=
F^{\mathcal M_i}
\left(
t_1^{\mathcal M_i}(\bar f(i)),
\ldots,
t_m^{\mathcal M_i}(\bar f(i))
\right).
\]
Но правая часть — это ровно
\[
t^{\mathcal M_i}(f_1(i),\ldots,f_n(i)).
\]
Следовательно,
\[
h(i)=g_t(i)
\]
для всех $i\in I$. Поэтому
\[
[h]=[g_t].
\]

Лемма доказана.
\end{proof}


\subsection{Теорема Лося}

Пусть $L$ — язык первого порядка. Пусть
\[
\{\mathcal M_i\}_{i\in I}
\]
— семейство $L$-структур, и пусть $\mathcal U$ — ультрафильтр на $I$.

Рассмотрим ультрапроизведение
\[
\mathcal M=
\prod_{i\in I}\mathcal M_i/\mathcal U.
\]

Напомним, что элементы структуры $\mathcal M$ имеют вид
\[
[f],
\]
где
\[
f\in\prod_{i\in I}M_i.
\]
При этом
\[
[f]=[g]
\]
тогда и только тогда, когда
\[
\{i\in I\mid f(i)=g(i)\}\in\mathcal U.
\]

Теорема Лося утверждает, что истинность любой формулы первого порядка в ультрапроизведении определяется покоординатно на большом множестве индексов.

\begin{theorem}[Теорема Лося]
Пусть $\varphi(x_1,\ldots,x_n)$ — $L$-формула. Тогда для любых
\[
f_1,\ldots,f_n\in\prod_{i\in I}M_i
\]
выполнено:
\[
\mathcal M\models\varphi([f_1],\ldots,[f_n])
\]
тогда и только тогда, когда
\[
\{i\in I\mid
\mathcal M_i\models
\varphi(f_1(i),\ldots,f_n(i))
\}\in\mathcal U.
\]
\end{theorem}

\begin{remark}
Содержательно теорема Лося говорит:
\[
\text{формула истинна в ультрапроизведении}
\]
тогда и только тогда, когда
\[
\text{она истинна в координатах на большом множестве индексов}.
\]

Здесь слово ``большое'' понимается в смысле ультрафильтра $\mathcal U$.
\end{remark}

Для доказательства удобно ввести обозначение. Если
\[
\bar f=(f_1,\ldots,f_n),
\]
то для формулы $\varphi(\bar x)$ положим
\[
X_\varphi(\bar f)
=
\{i\in I\mid
\mathcal M_i\models\varphi(f_1(i),\ldots,f_n(i))
\}.
\]
Тогда утверждение теоремы можно записать короче:
\[
\mathcal M\models\varphi([f_1],\ldots,[f_n])
\qquad
\Longleftrightarrow
\qquad
X_\varphi(\bar f)\in\mathcal U.
\]

Доказательство проводится индукцией по построению формулы $\varphi$.

\begin{proof}
Докажем теорему Лося индукцией по построению формулы $\varphi$.

Достаточно рассмотреть атомарные формулы, отрицание, конъюнкцию и квантор существования. Остальные логические связки и квантор всеобщности выражаются через них.

\medskip

\textbf{1. Атомарная формула равенства.}

Пусть формула $\varphi(\bar x)$ имеет вид
\[
t_1(\bar x)=t_2(\bar x).
\]

По лемме о термах имеем
\[
t_1^{\mathcal M}([\bar f])=[g_1],
\qquad
t_2^{\mathcal M}([\bar f])=[g_2],
\]
где
\[
g_1(i)=t_1^{\mathcal M_i}(\bar f(i)),
\qquad
g_2(i)=t_2^{\mathcal M_i}(\bar f(i)).
\]

Тогда
\[
\mathcal M\models t_1([\bar f])=t_2([\bar f])
\]
тогда и только тогда, когда
\[
[g_1]=[g_2].
\]
По определению равенства классов это равносильно
\[
\{i\in I\mid g_1(i)=g_2(i)\}\in\mathcal U.
\]
Но
\[
g_1(i)=g_2(i)
\]
тогда и только тогда, когда
\[
\mathcal M_i\models
t_1(\bar f(i))=t_2(\bar f(i)).
\]
Следовательно,
\[
\mathcal M\models t_1([\bar f])=t_2([\bar f])
\]
тогда и только тогда, когда
\[
\{i\in I\mid
\mathcal M_i\models
t_1(\bar f(i))=t_2(\bar f(i))
\}\in\mathcal U.
\]
То есть теорема верна для атомарных формул равенства.

\medskip

\textbf{2. Атомарная формула с предикатным символом.}

Пусть формула $\varphi(\bar x)$ имеет вид
\[
R(t_1(\bar x),\ldots,t_m(\bar x)),
\]
где $R$ — $m$-местный предикатный символ.

По лемме о термах для каждого $j=1,\ldots,m$ имеем
\[
t_j^{\mathcal M}([\bar f])=[g_j],
\]
где
\[
g_j(i)=t_j^{\mathcal M_i}(\bar f(i)).
\]

Тогда
\[
\mathcal M\models
R(t_1([\bar f]),\ldots,t_m([\bar f]))
\]
тогда и только тогда, когда
\[
\mathcal M\models R([g_1],\ldots,[g_m]).
\]
По определению интерпретации предикатного символа в ультрапроизведении это равносильно
\[
\{i\in I\mid
\mathcal M_i\models
R(g_1(i),\ldots,g_m(i))
\}\in\mathcal U.
\]
Но по определению функций $g_j$ это то же самое, что
\[
\{i\in I\mid
\mathcal M_i\models
R(t_1^{\mathcal M_i}(\bar f(i)),\ldots,t_m^{\mathcal M_i}(\bar f(i)))
\}\in\mathcal U.
\]
Иначе говоря,
\[
\{i\in I\mid
\mathcal M_i\models
R(t_1(\bar f(i)),\ldots,t_m(\bar f(i)))
\}\in\mathcal U.
\]
Значит, теорема верна для атомарных формул с предикатным символом.

\medskip

\textbf{3. Отрицание.}

Пусть
\[
\varphi(\bar x)=\neg\psi(\bar x).
\]
Предположим, что теорема уже доказана для $\psi$.

Тогда
\[
\mathcal M\models\neg\psi([\bar f])
\]
тогда и только тогда, когда
\[
\mathcal M\not\models\psi([\bar f]).
\]
По индукционному предположению
\[
\mathcal M\models\psi([\bar f])
\]
тогда и только тогда, когда
\[
X_\psi(\bar f)\in\mathcal U.
\]
Следовательно,
\[
\mathcal M\not\models\psi([\bar f])
\]
тогда и только тогда, когда
\[
X_\psi(\bar f)\notin\mathcal U.
\]

Так как $\mathcal U$ — ультрафильтр, для любого $A\subseteq I$ выполнено
\[
A\notin\mathcal U
\qquad
\Longleftrightarrow
\qquad
I\setminus A\in\mathcal U.
\]
Поэтому
\[
X_\psi(\bar f)\notin\mathcal U
\]
тогда и только тогда, когда
\[
I\setminus X_\psi(\bar f)\in\mathcal U.
\]

Но
\[
I\setminus X_\psi(\bar f)=X_{\neg\psi}(\bar f).
\]
Следовательно,
\[
\mathcal M\models\neg\psi([\bar f])
\]
тогда и только тогда, когда
\[
X_{\neg\psi}(\bar f)\in\mathcal U.
\]

Значит, теорема верна для отрицания.

\medskip

\textbf{4. Конъюнкция.}

Пусть
\[
\varphi(\bar x)=\psi(\bar x)\wedge\theta(\bar x).
\]
Предположим, что теорема уже доказана для $\psi$ и $\theta$.

Тогда
\[
\mathcal M\models(\psi\wedge\theta)([\bar f])
\]
тогда и только тогда, когда одновременно
\[
\mathcal M\models\psi([\bar f])
\]
и
\[
\mathcal M\models\theta([\bar f]).
\]

По индукционному предположению это равносильно условиям
\[
X_\psi(\bar f)\in\mathcal U
\]
и
\[
X_\theta(\bar f)\in\mathcal U.
\]

Так как $\mathcal U$ — фильтр, это равносильно
\[
X_\psi(\bar f)\cap X_\theta(\bar f)\in\mathcal U.
\]

Но
\[
X_{\psi\wedge\theta}(\bar f)
=
X_\psi(\bar f)\cap X_\theta(\bar f).
\]
Следовательно,
\[
\mathcal M\models(\psi\wedge\theta)([\bar f])
\]
тогда и только тогда, когда
\[
X_{\psi\wedge\theta}(\bar f)\in\mathcal U.
\]

Значит, теорема верна для конъюнкции.

\medskip

\textbf{5. Квантор существования.}

Пусть
\[
\varphi(\bar x)=\exists y\,\psi(\bar x,y).
\]
Предположим, что теорема уже доказана для $\psi$.

Докажем оба направления.

Сначала пусть
\[
\mathcal M\models\exists y\,\psi([\bar f],y).
\]
Тогда существует элемент ультрапроизведения
\[
[g]\in M
\]
такой, что
\[
\mathcal M\models\psi([\bar f],[g]).
\]

По индукционному предположению
\[
\{i\in I\mid
\mathcal M_i\models \psi(\bar f(i),g(i))
\}\in\mathcal U.
\]

Если для индекса $i$ выполнено
\[
\mathcal M_i\models \psi(\bar f(i),g(i)),
\]
то тем более
\[
\mathcal M_i\models \exists y\,\psi(\bar f(i),y).
\]
Следовательно,
\[
\{i\in I\mid
\mathcal M_i\models \psi(\bar f(i),g(i))
\}
\subseteq
X_{\exists y\,\psi}(\bar f).
\]
Так как левое множество принадлежит $\mathcal U$, а ультрафильтр замкнут вверх, получаем
\[
X_{\exists y\,\psi}(\bar f)\in\mathcal U.
\]

Теперь докажем обратное направление. Пусть
\[
X_{\exists y\,\psi}(\bar f)\in\mathcal U.
\]
То есть
\[
\{i\in I\mid
\mathcal M_i\models \exists y\,\psi(\bar f(i),y)
\}\in\mathcal U.
\]

Для каждого
\[
i\in X_{\exists y\,\psi}(\bar f)
\]
выберем элемент
\[
g(i)\in M_i
\]
такой, что
\[
\mathcal M_i\models\psi(\bar f(i),g(i)).
\]

Для индексов
\[
i\notin X_{\exists y\,\psi}(\bar f)
\]
выберем $g(i)\in M_i$ произвольно. Так мы получаем функцию
\[
g\in\prod_{i\in I}M_i.
\]

По построению множество
\[
\{i\in I\mid
\mathcal M_i\models\psi(\bar f(i),g(i))
\}
\]
содержит множество
\[
X_{\exists y\,\psi}(\bar f).
\]
Так как
\[
X_{\exists y\,\psi}(\bar f)\in\mathcal U
\]
и $\mathcal U$ замкнут вверх, получаем
\[
\{i\in I\mid
\mathcal M_i\models\psi(\bar f(i),g(i))
\}\in\mathcal U.
\]

По индукционному предположению для формулы $\psi$ имеем
\[
\mathcal M\models\psi([\bar f],[g]).
\]
Следовательно,
\[
\mathcal M\models\exists y\,\psi([\bar f],y).
\]

Итак, теорема верна для квантора существования.

\medskip

Остальные логические связки выражаются через $\neg$ и $\wedge$, например
\[
\psi\vee\theta
\equiv
\neg(\neg\psi\wedge\neg\theta),
\]
а квантор всеобщности выражается через отрицание и квантор существования:
\[
\forall y\,\psi
\equiv
\neg\exists y\,\neg\psi.
\]

Следовательно, утверждение теоремы верно для любой $L$-формулы $\varphi$.
\end{proof}

\begin{remark}
В доказательстве отрицания существенно используется то, что $\mathcal U$ является ультрафильтром, а не просто фильтром.

Для обычного фильтра из
\[
A\notin\mathcal F
\]
не следует, что
\[
I\setminus A\in\mathcal F.
\]
Именно ультрафильтр умеет выбирать одну из двух частей разбиения
\[
I=A\sqcup(I\setminus A).
\]
\end{remark}

\begin{remark}
В доказательстве квантора существования используется выбор свидетелей $g(i)$ в тех структурах $\mathcal M_i$, где существует подходящий элемент. Такое рассуждение обычно опирается на стандартную форму аксиомы выбора.
\end{remark}


\begin{corollary}
Пусть $\sigma$ — $L$-предложение. Тогда
\[
\prod_{i\in I}\mathcal M_i/\mathcal U\models\sigma
\]
тогда и только тогда, когда
\[
\{i\in I\mid \mathcal M_i\models\sigma\}\in\mathcal U.
\]
\end{corollary}

\begin{proof}
Это частный случай теоремы Лося для формулы без свободных переменных.
\end{proof}

\begin{corollary}
Если $\sigma$ истинно во всех структурах $\mathcal M_i$, то
\[
\prod_{i\in I}\mathcal M_i/\mathcal U\models\sigma.
\]
\end{corollary}

\begin{proof}
Если
\[
\mathcal M_i\models\sigma
\]
для всех $i\in I$, то
\[
\{i\in I\mid \mathcal M_i\models\sigma\}=I.
\]
Так как
\[
I\in\mathcal U,
\]
по теореме Лося получаем
\[
\prod_{i\in I}\mathcal M_i/\mathcal U\models\sigma.
\]
\end{proof}

\begin{remark}
Таким образом, всякое предложение первого порядка, истинное во всех координатных структурах, истинно и в их ультрапроизведении.

Более того, по теореме Лося достаточно, чтобы предложение было истинно не во всех координатах, а на большом множестве индексов в смысле ультрафильтра.
\end{remark}


\section{Компактность и нестандартные модели}

\subsection{Компактность для семантического следования}

Теорему компактности можно переформулировать не только через выполнимость множеств предложений, но и через семантическое следование.

Напомним, что запись
\[
\Sigma\models\varphi
\]
означает: всякая модель множества предложений $\Sigma$ является моделью предложения $\varphi$.

Иными словами,
\[
\Sigma\models\varphi
\]
тогда и только тогда, когда для любой $L$-структуры $\mathcal A$
\[
\mathcal A\models\Sigma
\qquad
\Longrightarrow
\qquad
\mathcal A\models\varphi.
\]

\begin{theorem}[Компактность для семантического следования]
Пусть $\Sigma$ — множество $L$-предложений, а $\varphi$ — $L$-предложение. Если
\[
\Sigma\models\varphi,
\]
то существует конечное подмножество
\[
\Sigma_0\subseteq\Sigma
\]
такое, что
\[
\Sigma_0\models\varphi.
\]
\end{theorem}

\begin{proof}
Пусть
\[
\Sigma\models\varphi.
\]
Докажем, что существует конечное подмножество $\Sigma_0\subseteq\Sigma$, такое что
\[
\Sigma_0\models\varphi.
\]

Предположим противное. Тогда для любого конечного подмножества
\[
\Sigma_0\subseteq\Sigma
\]
неверно, что
\[
\Sigma_0\models\varphi.
\]
То есть
\[
\Sigma_0\not\models\varphi.
\]

По определению семантического следования это означает, что для любого конечного
\[
\Sigma_0\subseteq\Sigma
\]
существует $L$-структура $\mathcal A_{\Sigma_0}$, такая что
\[
\mathcal A_{\Sigma_0}\models\Sigma_0,
\]
но
\[
\mathcal A_{\Sigma_0}\not\models\varphi.
\]
Так как $\varphi$ — предложение, то
\[
\mathcal A_{\Sigma_0}\not\models\varphi
\]
равносильно
\[
\mathcal A_{\Sigma_0}\models\neg\varphi.
\]

Следовательно, для любого конечного подмножества
\[
\Sigma_0\subseteq\Sigma
\]
множество предложений
\[
\Sigma_0\cup\{\neg\varphi\}
\]
выполнимо.

Покажем, что отсюда следует конечная выполнимость множества
\[
\Sigma\cup\{\neg\varphi\}.
\]

Действительно, пусть
\[
\Gamma_0\subseteq \Sigma\cup\{\neg\varphi\}
\]
— произвольное конечное подмножество. Тогда оно имеет вид
\[
\Gamma_0\subseteq \Sigma_0\cup\{\neg\varphi\}
\]
для некоторого конечного
\[
\Sigma_0\subseteq\Sigma.
\]
А множество
\[
\Sigma_0\cup\{\neg\varphi\}
\]
выполнимо. Следовательно, и его подмножество $\Gamma_0$ выполнимо.

Значит, каждое конечное подмножество
\[
\Sigma\cup\{\neg\varphi\}
\]
выполнимо. По теореме компактности всё множество
\[
\Sigma\cup\{\neg\varphi\}
\]
выполнимо.

Следовательно, существует $L$-структура $\mathcal A$, такая что
\[
\mathcal A\models\Sigma\cup\{\neg\varphi\}.
\]
То есть одновременно
\[
\mathcal A\models\Sigma
\]
и
\[
\mathcal A\models\neg\varphi.
\]

Но из исходного предположения
\[
\Sigma\models\varphi
\]
и из
\[
\mathcal A\models\Sigma
\]
следует
\[
\mathcal A\models\varphi.
\]

Получили противоречие:
\[
\mathcal A\models\varphi
\qquad
\text{и}
\qquad
\mathcal A\models\neg\varphi.
\]

Значит, наше предположение было неверным. Следовательно, существует конечное подмножество
\[
\Sigma_0\subseteq\Sigma
\]
такое, что
\[
\Sigma_0\models\varphi.
\]
\end{proof}

\begin{remark}
Эта форма компактности говорит, что если предложение $\varphi$ логически следует из, возможно, бесконечного множества предпосылок $\Sigma$, то на самом деле для этого следования достаточно конечного числа предпосылок.

Иными словами, семантическое следование в логике первого порядка является компактным:
\[
\Sigma\models\varphi
\qquad
\Longrightarrow
\qquad
\exists \Sigma_0\subseteq\Sigma,\ |\Sigma_0|<\infty:
\Sigma_0\models\varphi.
\]
\end{remark}

\begin{remark}
Эта формулировка эквивалентна обычной теореме компактности.

Действительно, условие
\[
\Sigma\models\varphi
\]
означает, что не существует модели, в которой одновременно истинны все предложения из $\Sigma$ и предложение $\neg\varphi$. То есть множество
\[
\Sigma\cup\{\neg\varphi\}
\]
невыполнимо.

По компактности, если оно невыполнимо, то уже некоторое конечное подмножество невыполнимо. Такое конечное подмножество имеет вид
\[
\Sigma_0\cup\{\neg\varphi\},
\]
где
\[
\Sigma_0\subseteq\Sigma
\]
конечно.

Невыполнимость множества
\[
\Sigma_0\cup\{\neg\varphi\}
\]
ровно означает, что
\[
\Sigma_0\models\varphi.
\]
\end{remark}


\subsection{Нестандартные модели истинной арифметики}

Рассмотрим язык полуколец
\[
L_{s\text{-}ring}
=
\langle +,\cdot,0,1\rangle.
\]
Стандартную структуру натуральных чисел в этом языке будем обозначать через
\[
\omega_{s\text{-}ring}
=
\langle \omega,+,\cdot,0,1\rangle,
\]
где
\[
\omega=\{0,1,2,\ldots\}.
\]

Полная теория этой структуры
\[
Th(\omega_{s\text{-}ring})
\]
называется истинной арифметикой. Она состоит из всех предложений языка
\[
L_{s\text{-}ring}
\]
истинных в стандартных натуральных числах.

\begin{theorem}[Существование нестандартной модели истинной арифметики]
Существует модель
\[
\mathcal M\models Th(\omega_{s\text{-}ring}),
\]
не изоморфная стандартной модели
\[
\omega_{s\text{-}ring}.
\]
\end{theorem}

\begin{proof}
Идея доказательства состоит в том, чтобы с помощью компактности добавить в модель новый элемент, который не равен ни одному стандартному натуральному числу.

Расширим язык
\[
L_{s\text{-}ring}
\]
новым константным символом $c$:
\[
L^+
=
L_{s\text{-}ring}\cup\{c\}.
\]

Рассмотрим множество $L^+$-предложений
\[
\Sigma
=
Th(\omega_{s\text{-}ring})
\cup
\{c\ne 0,\ c\ne 1,\ c\ne 1+1,\ c\ne 1+1+1,\ldots\}.
\]

То есть в $\Sigma$ входят все предложения истинной арифметики, а также бесконечный набор условий, говорящих, что новый элемент $c$ не равен ни одному стандартному натуральному числу.

Покажем, что $\Sigma$ конечно выполнимо.

Пусть
\[
\Sigma_0\subseteq \Sigma
\]
— конечное подмножество. Тогда $\Sigma_0$ содержит только конечное число условий вида
\[
c\ne n,
\]
где
\[
n=\underbrace{1+\cdots+1}_{n\text{ раз}}
\]
— стандартный натуральный терм.

Следовательно, среди этих условий есть только конечное число запретов:
\[
c\ne 0,\quad c\ne 1,\quad \ldots,\quad c\ne N
\]
для некоторого натурального $N$.

Рассмотрим стандартную структуру натуральных чисел и интерпретируем новый символ $c$ как число
\[
N+1.
\]
Получаем $L^+$-структуру
\[
\omega^+
=
\langle \omega,+,\cdot,0,1,c^{\omega^+}\rangle,
\]
где
\[
c^{\omega^+}=N+1.
\]

Тогда все предложения из $Th(\omega_{s\text{-}ring})$, входящие в $\Sigma_0$, истинны в $\omega^+$, потому что редукт $\omega^+$ к исходному языку $L_{s\text{-}ring}$ есть стандартная структура
\[
\omega_{s\text{-}ring}.
\]

Кроме того, все конечные условия вида
\[
c\ne 0,\quad c\ne 1,\quad \ldots,\quad c\ne N
\]
также выполнены, поскольку
\[
c^{\omega^+}=N+1.
\]

Значит, всякое конечное подмножество $\Sigma_0\subseteq\Sigma$ имеет модель. Следовательно, $\Sigma$ конечно выполнимо.

По теореме компактности множество $\Sigma$ выполнимо. Значит, существует $L^+$-структура
\[
\mathcal M^+
\]
такая, что
\[
\mathcal M^+\models\Sigma.
\]

Пусть теперь $\mathcal M$ — редукт структуры $\mathcal M^+$ к исходному языку
\[
L_{s\text{-}ring}.
\]
Так как
\[
\mathcal M^+\models Th(\omega_{s\text{-}ring}),
\]
то
\[
\mathcal M\models Th(\omega_{s\text{-}ring}).
\]

Осталось показать, что
\[
\mathcal M\not\cong \omega_{s\text{-}ring}.
\]

В структуре $\mathcal M^+$ есть элемент
\[
c^{\mathcal M^+}
\]
такой, что
\[
c^{\mathcal M^+}\ne 0,
\]
\[
c^{\mathcal M^+}\ne 1,
\]
\[
c^{\mathcal M^+}\ne 1+1,
\]
\[
c^{\mathcal M^+}\ne 1+1+1,
\]
и так далее.

То есть этот элемент не равен ни одному стандартному натуральному числу, заданному термом
\[
\underbrace{1+\cdots+1}_{n\text{ раз}}.
\]

В стандартной модели
\[
\omega_{s\text{-}ring}
\]
каждый элемент является некоторым стандартным натуральным числом:
\[
0,\quad 1,\quad 1+1,\quad 1+1+1,\quad \ldots
\]

Если бы
\[
\mathcal M\cong \omega_{s\text{-}ring},
\]
то каждый элемент $\mathcal M$ соответствовал бы некоторому стандартному натуральному числу. Но элемент $c^{\mathcal M^+}$ не равен ни одному из стандартных чисел.

Следовательно,
\[
\mathcal M\not\cong \omega_{s\text{-}ring}.
\]

Таким образом, мы построили модель истинной арифметики, не изоморфную стандартной модели натуральных чисел.
\end{proof}

\begin{definition}
Модель
\[
\mathcal M\models Th(\omega_{s\text{-}ring})
\]
называется \textbf{нестандартной моделью истинной арифметики}, если
\[
\mathcal M\not\cong \omega_{s\text{-}ring}.
\]
\end{definition}

\begin{remark}
Элемент
\[
c^{\mathcal M^+}
\]
можно понимать как нестандартное натуральное число. Оно является элементом модели арифметики, но не совпадает ни с одним стандартным числом
\[
0,\quad 1,\quad 2,\quad 3,\quad \ldots
\]

С точки зрения самой модели оно ведёт себя как натуральное число, но извне мы видим, что это новый, нестандартный элемент.
\end{remark}

\begin{remark}
Этот результат не противоречит полноте теории
\[
Th(\omega_{s\text{-}ring}).
\]

Полнота теории означает, что для любого предложения $\varphi$ языка
\[
L_{s\text{-}ring}
\]
выполнено:
\[
\varphi\in Th(\omega_{s\text{-}ring})
\qquad
\text{или}
\qquad
\neg\varphi\in Th(\omega_{s\text{-}ring}).
\]

То есть все модели этой теории удовлетворяют одним и тем же предложениям первого порядка.

Но полнота не означает единственность модели с точностью до изоморфизма. Единственность модели с точностью до изоморфизма называется категоричностью.

Таким образом,
\[
Th(\omega_{s\text{-}ring})
\]
полна, но не категорична.
\end{remark}

\begin{remark}
Содержательно этот результат показывает ограниченность логики первого порядка.

Даже полная теория стандартных натуральных чисел не задаёт стандартную модель однозначно с точностью до изоморфизма. Первый порядок не может исключить существование нестандартных элементов, если мы разрешаем бесконечные модели и используем компактность.
\end{remark}

\subsection{Нестандартные бесконечно малые элементы}

Аналогичный приём можно применить к вещественным числам.

Рассмотрим структуру упорядоченного поля вещественных чисел
\[
\mathbb R_{o\text{-}ring}
=
\langle\mathbb R,+,\cdot,-,0,1,<\rangle.
\]

Расширим язык новым константным символом $c$ и рассмотрим множество предложений
\[
\Sigma
=
Th(\mathbb R_{o\text{-}ring})
\cup
\left\{
0<c,\ 
c<1,\ 
c<\frac12,\ 
c<\frac13,\ 
c<\frac14,\ldots
\right\}.
\]

Более формально, для каждого натурального $n>0$ добавляется условие
\[
0<c<\frac1n.
\]

Любое конечное подмножество $\Sigma_0\subseteq\Sigma$ выполнимо в стандартных вещественных числах. Действительно, если в $\Sigma_0$ встречаются только условия
\[
0<c<1,\quad 0<c<\frac12,\quad \ldots,\quad 0<c<\frac1N,
\]
то можно интерпретировать $c$ как положительное вещественное число, меньшее $\frac1N$, например
\[
c^{\mathbb R^+}=\frac1{N+1}.
\]

Значит, по теореме компактности всё множество $\Sigma$ имеет модель.

В этой модели существует положительный элемент $c$, такой что
\[
0<c<\frac1n
\]
для любого стандартного натурального $n>0$.

Такой элемент называется \textbf{положительным бесконечно малым}.

\begin{remark}
В стандартном поле вещественных чисел такого элемента нет. Действительно, если $c>0$ — обычное вещественное число, то по архимедовости существует $n\in\mathbb N$, такое что
\[
\frac1n<c.
\]
Поэтому невозможно, чтобы
\[
0<c<\frac1n
\]
для всех $n>0$.

Значит, модель, построенная по компактности, не изоморфна стандартному полю вещественных чисел.
\end{remark}


\subsection{Элементарно эквивалентные неизоморфные модели}

В этом разделе докажем ещё одно важное следствие теоремы компактности: у всякой бесконечной структуры существует элементарно эквивалентная, но неизоморфная ей структура.

\begin{definition}
Пусть $\mathcal A$ и $\mathcal B$ — $L$-структуры. Говорят, что $\mathcal A$ и $\mathcal B$ \textbf{элементарно эквивалентны}, если они удовлетворяют одним и тем же $L$-предложениям.

Обозначение:
\[
\mathcal A\equiv\mathcal B.
\]

То есть
\[
\mathcal A\equiv\mathcal B
\]
означает, что для любого $L$-предложения $\varphi$
\[
\mathcal A\models\varphi
\qquad
\Longleftrightarrow
\qquad
\mathcal B\models\varphi.
\]

Эквивалентно,
\[
Th(\mathcal A)=Th(\mathcal B).
\]
\end{definition}

\begin{remark}
Элементарная эквивалентность слабее изоморфизма.

Если
\[
\mathcal A\cong\mathcal B,
\]
то
\[
\mathcal A\equiv\mathcal B.
\]
Действительно, изоморфные структуры не различаются никакими формулами первого порядка.

Однако обратное в общем случае неверно: структуры могут удовлетворять одним и тем же предложениям первого порядка, но не быть изоморфными.
\end{remark}

\begin{theorem}[Существование элементарно эквивалентной модели большей мощности]
Пусть $\mathcal A$ — бесконечная $L$-структура. Тогда существует $L$-структура $\mathcal B$, такая что
\[
\mathcal B\equiv\mathcal A,
\]
но
\[
\mathcal B\not\cong\mathcal A.
\]
Более того, можно добиться
\[
|B|>|A|.
\]
\end{theorem}

\begin{proof}
Идея доказательства состоит в том, чтобы с помощью компактности построить модель теории
\[
Th(\mathcal A)
\]
с носителем строго большей мощности, чем $A$.

Выберем множество индексов $I$ такое, что
\[
|I|>|A|.
\]
Такое множество существует, например можно взять
\[
I=\mathcal P(A),
\]
поскольку по теореме Кантора
\[
|\mathcal P(A)|>|A|.
\]

Расширим язык $L$ новыми константными символами:
\[
L_I=L\cup\{c_i\mid i\in I\}.
\]

Рассмотрим множество $L_I$-предложений
\[
\Sigma
=
Th(\mathcal A)
\cup
\{c_i\ne c_j\mid i,j\in I,\ i\ne j\}.
\]

Первая часть,
\[
Th(\mathcal A),
\]
требует, чтобы будущая модель была элементарно эквивалентна $\mathcal A$.

Вторая часть,
\[
\{c_i\ne c_j\mid i\ne j\},
\]
требует, чтобы все новые константы интерпретировались попарно различными элементами.

Покажем, что $\Sigma$ конечно выполнимо.

Пусть
\[
\Sigma_0\subseteq\Sigma
\]
— конечное подмножество. Тогда в $\Sigma_0$ участвует только конечное число новых констант:
\[
c_{i_1},\ldots,c_{i_n}.
\]

Поскольку структура $\mathcal A$ бесконечна, в её носителе можно выбрать попарно различные элементы
\[
a_1,\ldots,a_n\in A.
\]

Расширим структуру $\mathcal A$ до $L_I$-структуры $\mathcal A^+$ следующим образом:
\[
c_{i_k}^{\mathcal A^+}=a_k
\qquad
\text{для }k=1,\ldots,n.
\]
Остальные константы $c_i$, которые не участвуют в $\Sigma_0$, интерпретируем произвольно.

Тогда все предложения из $Th(\mathcal A)$, входящие в $\Sigma_0$, истинны в $\mathcal A^+$, потому что редукт $\mathcal A^+$ к исходному языку $L$ есть сама структура $\mathcal A$.

Также в $\mathcal A^+$ истинны все неравенства между константами, входящие в $\Sigma_0$, потому что элементы
\[
a_1,\ldots,a_n
\]
были выбраны попарно различными.

Следовательно,
\[
\mathcal A^+\models\Sigma_0.
\]
То есть любое конечное подмножество $\Sigma$ имеет модель.

По теореме компактности всё множество $\Sigma$ выполнимо. Значит, существует $L_I$-структура $\mathcal B^+$ такая, что
\[
\mathcal B^+\models\Sigma.
\]

Теперь забудем новые константы $c_i$ и возьмём редукт структуры $\mathcal B^+$ к исходному языку $L$. Обозначим его через
\[
\mathcal B.
\]

Так как
\[
\mathcal B^+\models Th(\mathcal A),
\]
то
\[
\mathcal B\models Th(\mathcal A).
\]
Следовательно,
\[
\mathcal B\equiv\mathcal A.
\]

Осталось показать, что $\mathcal B$ не изоморфна $\mathcal A$.

В структуре $\mathcal B^+$ все константы $c_i$ интерпретируются попарно различными элементами:
\[
c_i^{\mathcal B^+}\ne c_j^{\mathcal B^+}
\qquad
\text{при } i\ne j.
\]
Следовательно, в носителе $B$ есть по крайней мере $|I|$ различных элементов. Поэтому
\[
|B|\ge |I|.
\]
А множество $I$ было выбрано так, что
\[
|I|>|A|.
\]
Значит,
\[
|B|>|A|.
\]

Изоморфные структуры имеют носители одинаковой мощности. Поэтому
\[
\mathcal B\not\cong\mathcal A.
\]

Итак, мы построили структуру $\mathcal B$, такую что
\[
\mathcal B\equiv\mathcal A,
\]
но
\[
\mathcal B\not\cong\mathcal A.
\]
\end{proof}

\begin{remark}
Эта теорема показывает, что логика первого порядка плохо контролирует мощности бесконечных структур.

Даже если мы берём полную теорию одной конкретной бесконечной структуры
\[
Th(\mathcal A),
\]
у этой теории может быть модель большей мощности, чем сама $\mathcal A$.
\end{remark}

\begin{remark}
Доказанная теорема является слабой формой восходящей теоремы Лёвенгейма--Сколема.

Классическая восходящая теорема Лёвенгейма--Сколема утверждает, что если теория имеет бесконечную модель, то она имеет модели всех достаточно больших мощностей.

Здесь мы доказали более простой, но очень важный частный случай: у всякой бесконечной структуры есть элементарно эквивалентная модель строго большей мощности.
\end{remark}

\begin{remark}
Это рассуждение похоже на построение нестандартной модели истинной арифметики.

Там мы добавляли одну новую константу $c$ и требовали:
\[
c\ne 0,\quad c\ne 1,\quad c\ne 1+1,\quad \ldots
\]
То есть заставляли модель иметь элемент, отличный от всех стандартных натуральных чисел.

Здесь мы добавляем много новых констант
\[
c_i,\qquad i\in I,
\]
и требуем
\[
c_i\ne c_j
\qquad
(i\ne j).
\]
То есть заставляем новую модель иметь много различных элементов.
\end{remark}

\subsection{Аксиоматизируемые классы и конечные модели}

Пусть $L$ — язык первого порядка, и пусть $\Sigma$ — множество $L$-предложений.

\begin{definition}
Классом моделей множества предложений $\Sigma$ называется класс всех $L$-структур, удовлетворяющих всем предложениям из $\Sigma$:
\[
Mod(\Sigma)
=
\{\mathcal M\mid \mathcal M\models\Sigma\}.
\]
\end{definition}

То есть
\[
\mathcal M\in Mod(\Sigma)
\]
тогда и только тогда, когда
\[
\mathcal M\models\sigma
\]
для любого
\[
\sigma\in\Sigma.
\]

\begin{definition}
Класс $L$-структур $\mathcal C$ называется \textbf{аксиоматизируемым}, если существует множество $L$-предложений $\Sigma$, такое что
\[
\mathcal C=Mod(\Sigma).
\]
\end{definition}

Иными словами, класс структур аксиоматизируем, если его можно задать некоторым набором аксиом первого порядка.

\begin{definition}
Класс $L$-структур $\mathcal C$ называется \textbf{элементарным}, если существует одно $L$-предложение $\varphi$, такое что
\[
\mathcal C=Mod(\varphi).
\]
\end{definition}

Таким образом, элементарный класс задаётся одной аксиомой, а аксиоматизируемый класс может задаваться бесконечным множеством аксиом.

\begin{remark}
Если класс задаётся конечным набором предложений
\[
\varphi_1,\ldots,\varphi_n,
\]
то он элементарен, потому что конечный набор можно заменить одной конъюнкцией:
\[
\varphi_1\wedge\cdots\wedge\varphi_n.
\]
\end{remark}

\subsection{Примеры аксиоматизируемых классов}

\begin{example}
Класс всех групп аксиоматизируем.

Например, в языке
\[
L_{grp}=\langle \cdot,{}^{-1},e\rangle
\]
аксиомы групп можно записать следующим образом.

Ассоциативность:
\[
\forall x\forall y\forall z\,
((x\cdot y)\cdot z=x\cdot(y\cdot z)).
\]

Нейтральный элемент:
\[
\forall x\,(e\cdot x=x\wedge x\cdot e=x).
\]

Обратный элемент:
\[
\forall x\,(x\cdot x^{-1}=e\wedge x^{-1}\cdot x=e).
\]

Следовательно, класс всех групп задаётся конечным набором предложений. Поэтому он является элементарным классом.
\end{example}

\begin{example}
Класс групп порядка ровно $3$ является элементарным.

Для начала запишем предложение, говорящее, что в структуре ровно три элемента:
\[
\exists x\exists y\exists z
\left(
x\ne y\wedge x\ne z\wedge y\ne z
\wedge
\forall t\,(t=x\vee t=y\vee t=z)
\right).
\]

Это предложение утверждает, что существуют три попарно различных элемента $x,y,z$, и всякий элемент структуры равен одному из них.

Если обозначить через $\varphi_{grp}$ конъюнкцию аксиом группы, а через $\varphi_{=3}$ предложение ``в структуре ровно три элемента'', то класс групп порядка $3$ задаётся одним предложением:
\[
\varphi_{grp}\wedge\varphi_{=3}.
\]

Значит, класс групп порядка $3$ элементарен.
\end{example}

\begin{example}
Для каждого фиксированного натурального числа $n\ge 1$ свойство ``в структуре есть хотя бы $n$ элементов'' выражается одним предложением:
\[
\psi_n
=
\exists x_1\ldots\exists x_n
\bigwedge_{1\le i<j\le n} x_i\ne x_j.
\]

Следовательно, для каждого фиксированного $n$ класс групп порядка хотя бы $n$ элементарен.
\end{example}

\subsection{Бесконечные и конечные модели}

Для каждого $n\ge 1$ пусть
\[
\psi_n
=
\exists x_1\ldots\exists x_n
\bigwedge_{1\le i<j\le n} x_i\ne x_j.
\]
Это предложение говорит, что в структуре есть хотя бы $n$ различных элементов.

\begin{proposition}
Класс бесконечных групп аксиоматизируем.
\end{proposition}

\begin{proof}
Пусть $\Sigma_{grp}$ — множество аксиом групп. Рассмотрим теорию
\[
\Sigma
=
\Sigma_{grp}\cup\{\psi_n\mid n\ge 1\}.
\]

Если группа $\mathcal G$ является моделью $\Sigma$, то она удовлетворяет всем предложениям $\psi_n$. Значит, для любого $n$ в ней есть хотя бы $n$ различных элементов. Следовательно, $\mathcal G$ бесконечна.

Обратно, если $\mathcal G$ — бесконечная группа, то она удовлетворяет аксиомам группы и для каждого $n$ содержит хотя бы $n$ различных элементов. Значит,
\[
\mathcal G\models\Sigma.
\]

Таким образом,
\[
Mod(\Sigma)
=
\{\text{бесконечные группы}\}.
\]
Следовательно, класс бесконечных групп аксиоматизируем.
\end{proof}

\begin{remark}
Класс бесконечных групп аксиоматизируем, но не задаётся одной формулой первого порядка. Интуитивно причина в том, что одна формула имеет конечную длину и не может одновременно потребовать существования $n$ различных элементов для всех $n$.
\end{remark}

\begin{proposition}
Класс конечных групп не аксиоматизируем в логике первого порядка.
\end{proposition}

\begin{proof}
Предположим противное. Пусть существует множество предложений $\Sigma$, такое что
\[
Mod(\Sigma)
=
\{\text{конечные группы}\}.
\]

Для каждого $n\ge 1$ рассмотрим предложение
\[
\psi_n
=
\exists x_1\ldots\exists x_n
\bigwedge_{1\le i<j\le n} x_i\ne x_j,
\]
говорящее, что в структуре есть хотя бы $n$ элементов.

Рассмотрим множество предложений
\[
\Gamma
=
\Sigma\cup\{\psi_n\mid n\ge 1\}.
\]

Покажем, что $\Gamma$ конечно выполнимо.

Пусть
\[
\Gamma_0\subseteq\Gamma
\]
— конечное подмножество. Тогда в $\Gamma_0$ входит только конечное число предложений вида $\psi_n$, скажем
\[
\psi_{n_1},\ldots,\psi_{n_k}.
\]
Пусть
\[
N=\max\{n_1,\ldots,n_k\}.
\]

Нужно найти конечную группу, имеющую хотя бы $N$ элементов. Такая группа существует, например циклическая группа
\[
\mathbb Z/m\mathbb Z
\]
при
\[
m\ge N.
\]

Она является конечной группой, значит удовлетворяет $\Sigma$, и имеет хотя бы $N$ элементов, значит удовлетворяет всем предложениям $\psi_{n_1},\ldots,\psi_{n_k}$.

Следовательно, $\Gamma_0$ выполнимо. Значит, $\Gamma$ конечно выполнимо.

По теореме компактности $\Gamma$ выполнимо. Пусть
\[
\mathcal G\models\Gamma.
\]

Так как
\[
\mathcal G\models\Sigma,
\]
то по предположению $\mathcal G$ должна быть конечной группой.

Но также
\[
\mathcal G\models\psi_n
\]
для каждого $n\ge 1$. Значит, в $\mathcal G$ есть хотя бы $n$ элементов для любого $n$. Следовательно, $\mathcal G$ бесконечна.

Получили противоречие: $\mathcal G$ одновременно конечна и бесконечна.

Следовательно, класс конечных групп не аксиоматизируем.
\end{proof}

\subsection{Компактностная лемма о больших конечных моделях}

Следующее утверждение является общей формой предыдущего рассуждения.

\begin{lemma}[Компактностная лемма о больших конечных моделях]
Пусть $\mathcal C$ — аксиоматизируемый класс $L$-структур. Предположим, что в $\mathcal C$ существуют сколь угодно большие конечные модели, то есть для любого $n\ge 1$ найдётся конечная структура
\[
\mathcal M_n\in\mathcal C
\]
такая, что
\[
|M_n|\ge n.
\]
Тогда класс $\mathcal C$ содержит бесконечную модель.
\end{lemma}

\begin{proof}
Так как $\mathcal C$ аксиоматизируем, существует множество $L$-предложений $\Sigma$, такое что
\[
\mathcal C=Mod(\Sigma).
\]

Для каждого $n\ge 1$ рассмотрим предложение
\[
\psi_n
=
\exists x_1\ldots\exists x_n
\bigwedge_{1\le i<j\le n} x_i\ne x_j.
\]
Оно утверждает, что в структуре есть хотя бы $n$ различных элементов.

Рассмотрим множество предложений
\[
\Gamma
=
\Sigma\cup\{\psi_n\mid n\ge 1\}.
\]

Покажем, что $\Gamma$ конечно выполнимо.

Пусть
\[
\Gamma_0\subseteq\Gamma
\]
— конечное подмножество. Тогда среди предложений $\psi_n$ в $\Gamma_0$ встречается только конечное число:
\[
\psi_{n_1},\ldots,\psi_{n_k}.
\]
Пусть
\[
N=\max\{n_1,\ldots,n_k\}.
\]

По предположению в классе $\mathcal C$ существует конечная модель $\mathcal M_N$ мощности хотя бы $N$:
\[
\mathcal M_N\in\mathcal C,
\qquad
|M_N|\ge N.
\]

Так как
\[
\mathcal M_N\in\mathcal C=Mod(\Sigma),
\]
имеем
\[
\mathcal M_N\models\Sigma.
\]
Кроме того, поскольку
\[
|M_N|\ge N,
\]
структура $\mathcal M_N$ удовлетворяет всем предложениям
\[
\psi_{n_1},\ldots,\psi_{n_k}.
\]
Значит,
\[
\mathcal M_N\models\Gamma_0.
\]

Следовательно, всякое конечное подмножество $\Gamma$ выполнимо. По теореме компактности всё множество $\Gamma$ выполнимо.

Пусть
\[
\mathcal M\models\Gamma.
\]
Тогда
\[
\mathcal M\models\Sigma,
\]
значит,
\[
\mathcal M\in\mathcal C.
\]
Также
\[
\mathcal M\models\psi_n
\]
для любого $n\ge 1$. Значит, в $\mathcal M$ есть хотя бы $n$ различных элементов для каждого $n$. Следовательно, $\mathcal M$ бесконечна.

Итак, класс $\mathcal C$ содержит бесконечную модель.
\end{proof}

\begin{remark}
Эта лемма показывает типичное ограничение логики первого порядка: если аксиоматизируемый класс содержит конечные модели сколь угодно большой мощности, то он не может состоять только из конечных моделей. Компактность заставляет в таком классе появиться бесконечную модель.
\end{remark}

\subsection{n-кручения и периодические абелевы группы}

Рассмотрим язык аддитивных групп
\[
L_{adgp}
=
\langle +,-,0\rangle,
\]

\[
G=\langle G,+,-,0\rangle,
\]
удовлетворяющая обычным аксиомам абелевых групп.

\begin{definition}
Элемент \(g\in G\) называется \textbf{элементом \(n\)-кручения}, если
\[
ng=0.
\]
\end{definition}

То есть элемент является элементом \(n\)-кручения, если при сложении с самим собой \(n\) раз он даёт нейтральный элемент группы.

\begin{example}
В группе
\[
\mathbb Z/6\mathbb Z
\]
элемент \(\overline{3}\) является элементом \(2\)-кручения, потому что
\[
2\overline{3}
=
\overline{3}+\overline{3}
=
\overline{6}
=
\overline{0}.
\]

Элемент \(\overline{2}\) является элементом \(3\)-кручения, потому что
\[
3\overline{2}
=
\overline{2}+\overline{2}+\overline{2}
=
\overline{6}
=
\overline{0}.
\]
\end{example}

\begin{definition}
Элемент \(g\in G\) называется \textbf{элементом кручения}, если существует натуральное число \(n\ge 1\), такое что
\[
ng=0.
\]
\end{definition}

Иными словами, элемент кручения — это элемент конечного порядка.

Для каждого фиксированного \(n\ge 1\) свойство
\[
nx=0
\]
выражается формулой первого порядка в языке \(L_{adgp}\).

Например:
\[
2x=0
\]
означает
\[
x+x=0,
\]
а
\[
3x=0
\]
означает
\[
x+x+x=0.
\]

Однако свойство
\[
x \text{ является элементом кручения}
\]
требует условия
\[
x=0
\quad\vee\quad
2x=0
\quad\vee\quad
3x=0
\quad\vee\quad
\cdots.
\]
Это бесконечная дизъюнкция, а формулы первого порядка должны быть конечными. Поэтому свойство ``быть элементом кручения'' не задаётся напрямую одной формулой первого порядка в языке групп.

\begin{definition}
Абелева группа \(G\) называется \textbf{периодической}, если каждый её элемент является элементом кручения.

То есть
\[
\forall g\in G\ \exists n\ge 1\quad ng=0.
\]
\end{definition}

\begin{example}
Любая конечная абелева группа является периодической.

Действительно, если \(G\) конечна, то любой элемент \(g\in G\) имеет конечный порядок. Следовательно, существует \(n\ge 1\), такое что
\[
ng=0.
\]
\end{example}

\begin{example}
Группа
\[
\mathbb Z
\]
с обычным сложением не является периодической.

Действительно, элемент \(1\in\mathbb Z\) имеет бесконечный порядок, потому что
\[
n\cdot 1=n\ne 0
\]
для любого \(n\ge 1\).
\end{example}


\begin{definition}
Абелева группа \(G\) называется \textbf{группой без кручения}, если для любого \(n\ge 1\) и любого \(x\in G\)
\[
nx=0
\quad\Longrightarrow\quad
x=0.
\]
\end{definition}

Для каждого фиксированного \(n\ge 1\) рассмотрим предложение
\[
\varphi_n:
\quad
\forall x\,(nx=0\to x=0).
\]

Это предложение говорит, что в группе нет ненулевых элементов \(n\)-кручения.

\begin{proposition}
Класс абелевых групп без кручения аксиоматизируем в языке \(L_{adgp}\).
\end{proposition}

\begin{proof}
Пусть \(\Sigma_{adgp}\) — множество аксиом абелевых групп. Рассмотрим множество предложений
\[
\Sigma
=
\Sigma_{adgp}\cup\{\varphi_n\mid n\ge 1\},
\]
где
\[
\varphi_n:
\quad
\forall x\,(nx=0\to x=0).
\]

Если \(G\models\Sigma\), то \(G\) является абелевой группой, и для каждого \(n\ge 1\) в ней нет ненулевых элементов \(n\)-кручения. Значит, \(G\) является группой без кручения.

Обратно, если \(G\) — абелева группа без кручения, то она удовлетворяет аксиомам абелевых групп и всем предложениям \(\varphi_n\). Следовательно,
\[
G\models\Sigma.
\]

Значит, класс абелевых групп без кручения аксиоматизируем.
\end{proof}

\subsubsection{Неаксиоматизируемость класса периодических абелевых групп}

Теперь покажем, что класс периодических абелевых групп ведёт себя иначе.

\begin{theorem}
Класс периодических абелевых групп не аксиоматизируем в логике первого порядка в языке
\[
L_{adgp}=\langle +,-,0\rangle.
\]
\end{theorem}

\begin{proof}
Предположим противное. Пусть существует множество \(L_{adgp}\)-предложений \(\Sigma\), такое что
\[
Mod(\Sigma)
=
\{\text{периодические абелевы группы}\}.
\]

Расширим язык \(L_{adgp}\) новым константным символом \(c\):
\[
L'=L_{adgp}\cup\{c\}.
\]

Мы хотим заставить элемент \(c\) иметь бесконечный порядок. Для этого добавим условия
\[
c\ne 0,
\]
\[
2c\ne 0,
\]
\[
3c\ne 0,
\]
и так далее.

Рассмотрим множество \(L'\)-предложений
\[
\Gamma
=
\Sigma
\cup
\{nc\ne 0\mid n\ge 1\}.
\]

Если \(\Gamma\) имеет модель, то эта модель удовлетворяет \(\Sigma\), а значит, по предположению, является периодической абелевой группой. Но одновременно элемент \(c\) в ней удовлетворяет условиям
\[
nc\ne 0
\]
для всех \(n\ge 1\). То есть \(c\) имеет бесконечный порядок. Это невозможно в периодической группе.

Остаётся показать, что \(\Gamma\) действительно имеет модель. Для этого используем компактность.

Пусть
\[
\Gamma_0\subseteq \Gamma
\]
— конечное подмножество. Тогда в \(\Gamma_0\) встречается только конечное число условий вида
\[
nc\ne 0.
\]
Пусть среди них встречаются только условия с номерами
\[
n_1,\ldots,n_k.
\]
Выберем натуральное число \(N\), такое что
\[
N>n_1,\ldots,N>n_k.
\]

Рассмотрим циклическую группу порядка \(N\):
\[
C_N=\mathbb Z/N\mathbb Z.
\]
Пусть
\[
c^{C_N}=\overline{1}.
\]
Тогда элемент \(c^{C_N}\) имеет порядок \(N\). Следовательно, для любого \(m\) такого, что
\[
1\le m<N,
\]
имеем
\[
m c^{C_N}\ne 0.
\]

В частности, все условия
\[
n_1c\ne 0,\ldots,n_kc\ne 0
\]
выполнены в \(C_N\).

Кроме того, \(C_N\) — конечная абелева группа, значит она периодическая. По предположению \(\Sigma\) аксиоматизирует класс всех периодических абелевых групп, поэтому
\[
C_N\models\Sigma.
\]

Следовательно,
\[
C_N\models\Gamma_0.
\]
То есть всякое конечное подмножество \(\Gamma\) выполнимо.

По теореме компактности всё множество \(\Gamma\) выполнимо. Значит, существует \(L'\)-структура \(G\), такая что
\[
G\models\Gamma.
\]

Так как
\[
G\models\Sigma,
\]
то \(G\) является периодической абелевой группой.

Но так как
\[
G\models nc\ne 0
\]
для любого \(n\ge 1\), элемент \(c^G\) имеет бесконечный порядок. Значит, \(G\) не является периодической группой.

Получили противоречие.

Следовательно, класс периодических абелевых групп не аксиоматизируем в логике первого порядка.
\end{proof}

\begin{remark}
Важно различать два свойства.

Первое свойство:
\[
\text{в группе нет ненулевых элементов конечного порядка}.
\]
Оно аксиоматизируемо схемой предложений
\[
\forall x\,(nx=0\to x=0),
\qquad
n\ge 1.
\]

Второе свойство:
\[
\text{каждый элемент имеет конечный порядок}.
\]
Оно требовало бы бесконечной дизъюнкции
\[
x=0
\vee
2x=0
\vee
3x=0
\vee
\cdots.
\]
Но такая бесконечная дизъюнкция не является формулой первого порядка.

Теорема компактности показывает, что это не просто техническая проблема записи: класса периодических абелевых групп действительно нельзя задать никаким множеством предложений первого порядка.
\end{remark}

\begin{remark}
Это ещё один пример того, как компактность выявляет ограничения выразимости первого порядка.

Каждое конечное приближение к условию ``элемент \(c\) имеет бесконечный порядок'' можно реализовать внутри некоторой конечной циклической группы. Компактность позволяет собрать все эти конечные приближения в одну модель, где появляется элемент бесконечного порядка.

Именно это противоречит периодичности.
\end{remark}









\subsection{Конечная аксиоматизируемость и дополнения классов}

Пусть $L$ — язык первого порядка.

\begin{definition}
Класс $L$-структур $\mathcal C$ называется \textbf{конечно аксиоматизируемым}, если существует конечное множество $L$-предложений
\[
\Sigma=\{\varphi_1,\ldots,\varphi_n\}
\]
такое, что
\[
\mathcal C=Mod(\Sigma).
\]
\end{definition}

Так как конечное множество предложений можно заменить одной конъюнкцией
\[
\varphi_1\wedge\cdots\wedge\varphi_n,
\]
то класс $\mathcal C$ конечно аксиоматизируем тогда и только тогда, когда существует одно $L$-предложение $\varphi$, такое что
\[
\mathcal C=Mod(\varphi).
\]

\begin{remark}
Если класс $\mathcal C$ конечно аксиоматизируем, то его можно задать одной аксиомой.

Действительно, если
\[
\mathcal C=Mod(\varphi_1,\ldots,\varphi_n),
\]
то
\[
\mathcal C=Mod(\varphi_1\wedge\cdots\wedge\varphi_n).
\]
\end{remark}

\begin{example}
Класс всех групп фиксированного конечного порядка $n$ конечно аксиоматизируем.

Аксиомы группы задаются конечным набором предложений. Условие ``в структуре ровно $n$ элементов'' также выражается одним предложением:
\[
\exists x_1\ldots\exists x_n
\left(
\bigwedge_{1\le i<j\le n}x_i\ne x_j
\wedge
\forall y\,(y=x_1\vee\cdots\vee y=x_n)
\right).
\]
Следовательно, класс групп порядка $n$ задаётся конечным набором аксиом.
\end{example}

\begin{example}
Класс всех конечных групп не является конечно аксиоматизируемым. Более того, как было показано ранее, этот класс вообще не является аксиоматизируемым в логике первого порядка.
\end{example}

\subsubsection{Лемма о конечном фрагменте аксиоматизации}

\begin{lemma}[Лемма о конечном фрагменте аксиоматизации]
Пусть
\[
\mathcal C=Mod(\Sigma),
\]
где $\Sigma$ — множество $L$-предложений.

Если класс $\mathcal C$ конечно аксиоматизируем, то существует конечное подмножество
\[
\Sigma_0\subseteq\Sigma
\]
такое, что
\[
\mathcal C=Mod(\Sigma_0).
\]
\end{lemma}

\begin{proof}
Пусть
\[
\mathcal C=Mod(\Sigma)
\]
и пусть $\mathcal C$ конечно аксиоматизируем.

Тогда существует одно $L$-предложение $\varphi$, такое что
\[
\mathcal C=Mod(\varphi).
\]
Следовательно,
\[
Mod(\Sigma)=Mod(\varphi).
\]

Докажем, что существует конечное подмножество
\[
\Sigma_0\subseteq\Sigma
\]
такое, что
\[
Mod(\Sigma_0)=Mod(\Sigma).
\]

Предположим противное: никакое конечное подмножество $\Sigma_0\subseteq\Sigma$ не аксиоматизирует класс $\mathcal C$.

Тогда для любого конечного
\[
\Sigma_0\subseteq\Sigma
\]
имеем
\[
Mod(\Sigma_0)\ne Mod(\Sigma).
\]

Поскольку
\[
\Sigma_0\subseteq\Sigma,
\]
всякая модель $\Sigma$ является моделью $\Sigma_0$. Поэтому
\[
Mod(\Sigma)\subseteq Mod(\Sigma_0).
\]
Значит, если
\[
Mod(\Sigma_0)\ne Mod(\Sigma),
\]
то включение строгое:
\[
Mod(\Sigma)\subsetneq Mod(\Sigma_0).
\]

Следовательно, для каждого конечного $\Sigma_0\subseteq\Sigma$ существует структура $\mathcal A$, такая что
\[
\mathcal A\models\Sigma_0,
\]
но
\[
\mathcal A\not\models\Sigma.
\]

Так как
\[
Mod(\Sigma)=Mod(\varphi),
\]
из условия
\[
\mathcal A\not\models\Sigma
\]
следует
\[
\mathcal A\not\models\varphi.
\]
Поскольку $\varphi$ — предложение, это равносильно
\[
\mathcal A\models\neg\varphi.
\]

Итак, для любого конечного подмножества
\[
\Sigma_0\subseteq\Sigma
\]
множество предложений
\[
\Sigma_0\cup\{\neg\varphi\}
\]
выполнимо.

Значит, множество
\[
\Sigma\cup\{\neg\varphi\}
\]
конечно выполнимо. Действительно, любое его конечное подмножество содержится в некотором множестве вида
\[
\Sigma_0\cup\{\neg\varphi\},
\]
где $\Sigma_0\subseteq\Sigma$ конечно, а такое множество выполнимо по доказанному выше.

По теореме компактности множество
\[
\Sigma\cup\{\neg\varphi\}
\]
выполнимо.

Следовательно, существует $L$-структура $\mathcal B$, такая что
\[
\mathcal B\models\Sigma\cup\{\neg\varphi\}.
\]
То есть одновременно
\[
\mathcal B\models\Sigma
\]
и
\[
\mathcal B\models\neg\varphi.
\]

Но из
\[
\mathcal B\models\Sigma
\]
следует
\[
\mathcal B\in Mod(\Sigma).
\]
А так как
\[
Mod(\Sigma)=Mod(\varphi),
\]
получаем
\[
\mathcal B\models\varphi.
\]

Итак,
\[
\mathcal B\models\varphi
\qquad
\text{и}
\qquad
\mathcal B\models\neg\varphi,
\]
что невозможно.

Получили противоречие. Следовательно, наше предположение было неверным.

Значит, существует конечное подмножество
\[
\Sigma_0\subseteq\Sigma
\]
такое, что
\[
Mod(\Sigma_0)=Mod(\Sigma).
\]
То есть
\[
\mathcal C=Mod(\Sigma_0).
\]
\end{proof}

\begin{remark}
Смысл леммы такой: если класс задан, возможно, бесконечной системой аксиом $\Sigma$, но при этом сам класс в принципе можно задать конечным числом аксиом, то нужные конечные аксиомы можно выбрать прямо из $\Sigma$.
\end{remark}

\subsubsection{Аксиоматизируемость дополнения класса}

Пусть $\mathcal C$ — класс $L$-структур. Его дополнение обозначим через
\[
\mathcal D.
\]
То есть
\[
\mathcal D=\{\mathcal A\mid \mathcal A\text{ — }L\text{-структура и }\mathcal A\notin\mathcal C\}.
\]

\begin{lemma}[Критерий конечной аксиоматизируемости через дополнение]
Пусть $\mathcal C$ — аксиоматизируемый класс $L$-структур, и пусть $\mathcal D$ — его дополнение.

Тогда
\[
\mathcal D \text{ аксиоматизируем}
\]
тогда и только тогда, когда
\[
\mathcal C \text{ конечно аксиоматизируем}.
\]

Более того, если $\mathcal C$ и $\mathcal D$ оба аксиоматизируемы, то оба они конечно аксиоматизируемы.
\end{lemma}

\begin{proof}
Докажем оба направления.

Сначала пусть $\mathcal C$ конечно аксиоматизируем. Тогда существует одно $L$-предложение $\varphi$, такое что
\[
\mathcal C=Mod(\varphi).
\]
Тогда дополнение $\mathcal D$ задаётся предложением
\[
\neg\varphi.
\]
Действительно,
\[
\mathcal A\in\mathcal D
\]
тогда и только тогда, когда
\[
\mathcal A\notin\mathcal C,
\]
то есть
\[
\mathcal A\not\models\varphi.
\]
Поскольку $\varphi$ — предложение, это равносильно
\[
\mathcal A\models\neg\varphi.
\]
Значит,
\[
\mathcal D=Mod(\neg\varphi).
\]
Следовательно, $\mathcal D$ аксиоматизируем. Более того, он конечно аксиоматизируем.

Теперь докажем обратное направление.

Пусть $\mathcal C$ аксиоматизируем и $\mathcal D$ аксиоматизируем. Тогда существуют множества $L$-предложений $\Sigma$ и $\Phi$ такие, что
\[
\mathcal C=Mod(\Sigma)
\]
и
\[
\mathcal D=Mod(\Phi).
\]

Так как $\mathcal C$ и $\mathcal D$ являются дополнениями друг друга, не существует структуры, которая одновременно принадлежит $\mathcal C$ и $\mathcal D$.

Значит, множество предложений
\[
\Sigma\cup\Phi
\]
невыполнимо.

По теореме компактности существует конечное подмножество
\[
\Sigma_0\cup\Phi_0\subseteq \Sigma\cup\Phi,
\]
где
\[
\Sigma_0\subseteq\Sigma,
\qquad
\Phi_0\subseteq\Phi,
\]
такое, что
\[
\Sigma_0\cup\Phi_0
\]
невыполнимо.

Докажем, что
\[
Mod(\Sigma_0)=Mod(\Sigma).
\]

Так как
\[
\Sigma_0\subseteq\Sigma,
\]
то всякая модель $\Sigma$ является моделью $\Sigma_0$. Поэтому
\[
Mod(\Sigma)\subseteq Mod(\Sigma_0).
\]

Теперь докажем обратное включение. Пусть
\[
\mathcal A\models\Sigma_0.
\]
Нужно показать, что
\[
\mathcal A\models\Sigma.
\]

Предположим противное:
\[
\mathcal A\not\models\Sigma.
\]
Тогда
\[
\mathcal A\notin Mod(\Sigma).
\]
Но
\[
Mod(\Sigma)=\mathcal C.
\]
Значит,
\[
\mathcal A\notin\mathcal C.
\]
Следовательно,
\[
\mathcal A\in\mathcal D.
\]
А так как
\[
\mathcal D=Mod(\Phi),
\]
получаем
\[
\mathcal A\models\Phi.
\]
В частности,
\[
\mathcal A\models\Phi_0.
\]

Но вместе с
\[
\mathcal A\models\Sigma_0
\]
это даёт
\[
\mathcal A\models\Sigma_0\cup\Phi_0.
\]
А это невозможно, потому что
\[
\Sigma_0\cup\Phi_0
\]
невыполнимо.

Получили противоречие. Значит,
\[
\mathcal A\models\Sigma.
\]

Следовательно,
\[
Mod(\Sigma_0)\subseteq Mod(\Sigma).
\]

Итак,
\[
Mod(\Sigma_0)=Mod(\Sigma).
\]

Так как $\Sigma_0$ конечно, класс $\mathcal C$ конечно аксиоматизируем.

Аналогично, поменяв местами $\Sigma$ и $\Phi$, получаем конечное подмножество
\[
\Phi_0\subseteq\Phi
\]
такое, что
\[
Mod(\Phi_0)=Mod(\Phi).
\]
Значит, класс $\mathcal D$ тоже конечно аксиоматизируем.
\end{proof}

\begin{remark}
Эта лемма говорит, что для аксиоматизируемого класса $\mathcal C$ аксиоматизируемость дополнения очень сильна. Если дополнение $\mathcal D$ тоже можно задать аксиомами первого порядка, то сам класс $\mathcal C$ уже задаётся конечным числом аксиом.
\end{remark}

\begin{remark}
На практике эту лемму часто используют так: если класс $\mathcal C$ аксиоматизируем, но мы можем показать, что его дополнение не аксиоматизируемо, то $\mathcal C$ не является конечно аксиоматизируемым.
\end{remark}








\section{Элементарные подструктуры и диаграммы}

\subsection{Элементарные подструктуры}

Пусть $\mathcal A$ и $\mathcal B$ — $L$-структуры.

Напомним, что запись
\[
\mathcal A\subseteq \mathcal B
\]
означает, что $\mathcal A$ является $L$-подструктурой структуры $\mathcal B$. То есть
\[
A\subseteq B,
\]
и интерпретации всех символов языка $L$ на $A$ согласованы с их интерпретациями в $\mathcal B$.

Однако обычная подструктура не обязана сохранять истинность всех формул первого порядка. Формулы с кванторами могут вести себя по-разному в подструктуре и в большой структуре.

\begin{definition}
Пусть $\mathcal A\subseteq\mathcal B$. Говорят, что $\mathcal A$ является \textbf{элементарной подструктурой} структуры $\mathcal B$, если для любой $L$-формулы
\[
\varphi(x_1,\ldots,x_n)
\]
и любых элементов
\[
a_1,\ldots,a_n\in A
\]
выполнено
\[
\mathcal A\models \varphi(a_1,\ldots,a_n)
\]
тогда и только тогда, когда
\[
\mathcal B\models \varphi(a_1,\ldots,a_n).
\]
\end{definition}

Обозначение:
\[
\mathcal A\preccurlyeq\mathcal B.
\]

То есть
\[
\mathcal A\preccurlyeq\mathcal B
\]
означает, что $\mathcal A\subseteq\mathcal B$, и $\mathcal A$ и $\mathcal B$ одинаково отвечают на все формулы первого порядка с параметрами из $A$.

\begin{remark}
Элементарная подструктура — это гораздо более сильное понятие, чем обычная подструктура.

Обычная подструктура гарантирует хорошее поведение термов и атомарных формул. Элементарная подструктура гарантирует сохранение истинности всех формул первого порядка, включая формулы с кванторами.
\end{remark}

\begin{example}
Рассмотрим язык аддитивных групп
\[
L_{adgp}=\langle +,-,0\rangle.
\]
Пусть
\[
\mathcal B=\mathbb Z_{adgp}
=
\langle \mathbb Z,+,-,0\rangle
\]
и пусть
\[
\mathcal A=\langle 2\mathbb Z,+,-,0\rangle.
\]
Тогда
\[
\mathcal A\subseteq \mathcal B,
\]
поскольку $2\mathbb Z$ является подгруппой группы $\mathbb Z$.

Однако
\[
\mathcal A\not\preccurlyeq \mathcal B.
\]

Действительно, рассмотрим формулу
\[
\psi(x):\quad \exists y\,(y+y=x).
\]
Эта формула говорит, что элемент $x$ делится на $2$ внутри рассматриваемой группы.

Возьмём элемент
\[
2\in 2\mathbb Z.
\]
В структуре $\mathcal B=\mathbb Z_{adgp}$ имеем
\[
\mathcal B\models \psi(2),
\]
потому что можно взять
\[
y=1,
\]
и тогда
\[
1+1=2.
\]

Но в структуре $\mathcal A=\langle 2\mathbb Z,+,-,0\rangle$ формула $\psi(2)$ ложна:
\[
\mathcal A\not\models \psi(2).
\]
Действительно, в $\mathcal A$ переменная $y$ пробегает только чётные целые числа. Чтобы было
\[
y+y=2,
\]
нужно было бы взять
\[
y=1,
\]
но
\[
1\notin 2\mathbb Z.
\]

Следовательно, существует формула $\psi(x)$ и параметр $2\in A$, такие что
\[
\mathcal B\models\psi(2),
\]
но
\[
\mathcal A\not\models\psi(2).
\]
Значит,
\[
\mathcal A\not\preccurlyeq\mathcal B.
\]
\end{example}

\subsection{Элементарные вложения}

Понятие элементарной подструктуры является частным случаем более общего понятия элементарного вложения.

\begin{definition}
Пусть $\mathcal A$ и $\mathcal B$ — $L$-структуры. Отображение
\[
\pi:A\to B
\]
называется \textbf{элементарным вложением} из $\mathcal A$ в $\mathcal B$, если для любой $L$-формулы
\[
\varphi(x_1,\ldots,x_n)
\]
и любого набора
\[
a_1,\ldots,a_n\in A
\]
выполнено
\[
\mathcal A\models\varphi(a_1,\ldots,a_n)
\]
тогда и только тогда, когда
\[
\mathcal B\models\varphi(\pi(a_1),\ldots,\pi(a_n)).
\]
\end{definition}

Кратко это можно записать так:
\[
\mathcal A\models\varphi(\bar a)
\qquad
\Longleftrightarrow
\qquad
\mathcal B\models\varphi(\pi(\bar a)).
\]

\begin{remark}
Любое элементарное вложение является $L$-вложением.

Действительно, элементарное вложение сохраняет истинность всех формул, в частности атомарных формул. А сохранение атомарных формул означает сохранение равенства, функций, отношений и констант языка.
\end{remark}

\begin{remark}
Если
\[
\mathcal A\subseteq\mathcal B,
\]
то включение
\[
\iota:A\hookrightarrow B
\]
является элементарным вложением тогда и только тогда, когда
\[
\mathcal A\preccurlyeq\mathcal B.
\]
\end{remark}

\subsection{Диаграммы}

Пусть $\mathcal A$ — $L$-структура. Для каждого элемента
\[
a\in A
\]
добавим в язык новый константный символ $c_a$.

Полученный язык обозначим через
\[
L_A.
\]
То есть
\[
L_A=L\cup\{c_a\mid a\in A\}.
\]

\begin{remark}
Через $\mathcal A^+$ обозначается естественное расширение структуры $\mathcal A$ до языка $L_A$, в котором каждый новый константный символ $c_a$ интерпретируется самим элементом $a$:
\[
c_a^{\mathcal A^+}=a.
\]
\end{remark}

Идея этого расширения языка состоит в том, что мы даём имена всем элементам структуры $\mathcal A$. После этого можно записывать предложения, которые говорят не только о структуре вообще, но и о конкретных элементах $a\in A$ через их имена $c_a$.


\begin{definition}
\textbf{Полной диаграммой} структуры $\mathcal A$ называется множество всех $L_A$-предложений, истинных в структуре $\mathcal A^+$:
\[
CDiag(\mathcal A)=Th_{L_A}(\mathcal A^+).
\]
\end{definition}

Иными словами,
\[
CDiag(\mathcal A)
=
\{\sigma\mid \sigma \text{ — }L_A\text{-предложение и } \mathcal A^+\models\sigma\}.
\]

Полная диаграмма называется полной потому, что она содержит все предложения расширенного языка $L_A$, истинные в $\mathcal A^+$, а не только атомарные факты.

\begin{remark}
Полная диаграмма структуры $\mathcal A$ полностью описывает $\mathcal A$ в языке, где каждый элемент $a\in A$ имеет своё имя $c_a$.
\end{remark}


Следующая лемма объясняет, зачем нужна полная диаграмма.

\begin{lemma}[Полная диаграмма и элементарные вложения]
Пусть $\mathcal A$ и $\mathcal B$ — $L$-структуры. Тогда существует элементарное вложение
\[
\pi:\mathcal A\to\mathcal B
\]
тогда и только тогда, когда структуру $\mathcal B$ можно расширить до $L_A$-структуры $\mathcal B^+$ так, что
\[
\mathcal B^+\models CDiag(\mathcal A).
\]
\end{lemma}

\begin{proof}
Докажем оба направления.

Сначала пусть существует элементарное вложение
\[
\pi:\mathcal A\to\mathcal B.
\]
Нужно построить расширение $\mathcal B^+$ структуры $\mathcal B$ до языка $L_A$, которое будет моделью полной диаграммы $CDiag(\mathcal A)$.

Определим интерпретацию новых констант в $\mathcal B^+$ следующим образом:
\[
c_a^{\mathcal B^+}=\pi(a)
\]
для каждого
\[
a\in A.
\]

Покажем, что
\[
\mathcal B^+\models CDiag(\mathcal A).
\]

Пусть
\[
\sigma\in CDiag(\mathcal A).
\]
Тогда
\[
\mathcal A^+\models\sigma.
\]

Так как $\sigma$ — предложение языка $L_A$, в нём встречается только конечное число новых констант:
\[
c_{a_1},\ldots,c_{a_n}.
\]
Поэтому $\sigma$ можно записать в виде
\[
\psi(c_{a_1},\ldots,c_{a_n}),
\]
где $\psi(x_1,\ldots,x_n)$ — некоторая $L$-формула.

Из
\[
\mathcal A^+\models\psi(c_{a_1},\ldots,c_{a_n})
\]
получаем
\[
\mathcal A\models\psi(a_1,\ldots,a_n).
\]

Так как $\pi$ — элементарное вложение,
\[
\mathcal B\models\psi(\pi(a_1),\ldots,\pi(a_n)).
\]

Но в $\mathcal B^+$ по определению
\[
c_{a_k}^{\mathcal B^+}=\pi(a_k)
\]
для каждого $k=1,\ldots,n$. Поэтому
\[
\mathcal B^+\models\psi(c_{a_1},\ldots,c_{a_n}).
\]
То есть
\[
\mathcal B^+\models\sigma.
\]

Так как $\sigma\in CDiag(\mathcal A)$ было произвольным, получаем
\[
\mathcal B^+\models CDiag(\mathcal A).
\]

Теперь докажем обратное направление.

Пусть структуру $\mathcal B$ можно расширить до $L_A$-структуры $\mathcal B^+$ так, что
\[
\mathcal B^+\models CDiag(\mathcal A).
\]

Определим отображение
\[
\pi:A\to B
\]
по правилу
\[
\pi(a)=c_a^{\mathcal B^+}.
\]

Покажем, что $\pi$ является элементарным вложением.

Пусть $\psi(x_1,\ldots,x_n)$ — произвольная $L$-формула и пусть
\[
a_1,\ldots,a_n\in A.
\]
Нужно доказать:
\[
\mathcal A\models\psi(a_1,\ldots,a_n)
\]
тогда и только тогда, когда
\[
\mathcal B\models\psi(\pi(a_1),\ldots,\pi(a_n)).
\]

Рассмотрим $L_A$-предложение
\[
\psi(c_{a_1},\ldots,c_{a_n}).
\]

Если
\[
\mathcal A\models\psi(a_1,\ldots,a_n),
\]
то
\[
\mathcal A^+\models\psi(c_{a_1},\ldots,c_{a_n}).
\]
Значит,
\[
\psi(c_{a_1},\ldots,c_{a_n})\in CDiag(\mathcal A).
\]
Так как
\[
\mathcal B^+\models CDiag(\mathcal A),
\]
получаем
\[
\mathcal B^+\models\psi(c_{a_1},\ldots,c_{a_n}).
\]
По определению отображения $\pi$ это означает
\[
\mathcal B\models\psi(\pi(a_1),\ldots,\pi(a_n)).
\]

Обратно, пусть
\[
\mathcal A\not\models\psi(a_1,\ldots,a_n).
\]
Тогда
\[
\mathcal A\models\neg\psi(a_1,\ldots,a_n).
\]
Следовательно,
\[
\mathcal A^+\models\neg\psi(c_{a_1},\ldots,c_{a_n}).
\]
Значит,
\[
\neg\psi(c_{a_1},\ldots,c_{a_n})\in CDiag(\mathcal A).
\]
Так как
\[
\mathcal B^+\models CDiag(\mathcal A),
\]
получаем
\[
\mathcal B^+\models\neg\psi(c_{a_1},\ldots,c_{a_n}).
\]
По определению $\pi$ это означает
\[
\mathcal B\models\neg\psi(\pi(a_1),\ldots,\pi(a_n)).
\]
То есть
\[
\mathcal B\not\models\psi(\pi(a_1),\ldots,\pi(a_n)).
\]

Таким образом,
\[
\mathcal A\models\psi(a_1,\ldots,a_n)
\]
тогда и только тогда, когда
\[
\mathcal B\models\psi(\pi(a_1),\ldots,\pi(a_n)).
\]

Следовательно, $\pi$ является элементарным вложением.

Отдельно заметим, что $\pi$ инъективно. Если $a\ne b$ в $A$, то
\[
\mathcal A^+\models c_a\ne c_b.
\]
Значит,
\[
c_a\ne c_b\in CDiag(\mathcal A).
\]
Так как
\[
\mathcal B^+\models CDiag(\mathcal A),
\]
получаем
\[
c_a^{\mathcal B^+}\ne c_b^{\mathcal B^+}.
\]
То есть
\[
\pi(a)\ne\pi(b).
\]

Итак, $\pi$ — элементарное вложение из $\mathcal A$ в $\mathcal B$.
\end{proof}

\subsection{Атомарная диаграмма}

Кроме полной диаграммы, часто рассматривают обычную, или атомарную, диаграмму.

\begin{definition}
\textbf{Диаграммой} структуры $\mathcal A$ называется множество всех атомарных и отрицаний атомарных $L_A$-предложений, истинных в структуре $\mathcal A^+$.

Обозначение:
\[
Diag(\mathcal A).
\]
\end{definition}

То есть
\[
Diag(\mathcal A)
\]
содержит предложения вида
\[
\tau_1=\tau_2,
\qquad
\tau_1\ne\tau_2,
\]
а также предложения вида
\[
R(\tau_1,\ldots,\tau_n)
\]
и
\[
\neg R(\tau_1,\ldots,\tau_n),
\]
если они истинны в $\mathcal A^+$.

\begin{remark}
Обычная диаграмма $Diag(\mathcal A)$ кодирует только атомарную информацию о структуре $\mathcal A$ с именами для всех её элементов.

Полная диаграмма $CDiag(\mathcal A)$ кодирует все формулы первого порядка, истинные в $\mathcal A^+$.
\end{remark}

\begin{lemma}[Диаграмма и вложения]
Пусть $\mathcal A$ и $\mathcal B$ — $L$-структуры. Тогда существует $L$-вложение
\[
\pi:\mathcal A\to\mathcal B
\]
тогда и только тогда, когда структуру $\mathcal B$ можно расширить до $L_A$-структуры $\mathcal B^+$ так, что
\[
\mathcal B^+\models Diag(\mathcal A).
\]
\end{lemma}

\begin{proof}
Доказательство аналогично доказательству леммы о полной диаграмме и элементарных вложениях.

Если дано $L$-вложение
\[
\pi:\mathcal A\to\mathcal B,
\]
то расширяем $\mathcal B$ до языка $L_A$, полагая
\[
c_a^{\mathcal B^+}=\pi(a)
\]
для всех $a\in A$.

Так как $\pi$ является $L$-вложением, оно сохраняет истинность атомарных формул и отрицаний атомарных формул. Поэтому все предложения из $Diag(\mathcal A)$ истинны в $\mathcal B^+$:
\[
\mathcal B^+\models Diag(\mathcal A).
\]

Обратно, пусть
\[
\mathcal B^+\models Diag(\mathcal A).
\]
Определим
\[
\pi(a)=c_a^{\mathcal B^+}.
\]

Так как в $Diag(\mathcal A)$ содержатся все истинные атомарные и отрицания атомарных предложения о константах $c_a$, отображение $\pi$ сохраняет функции, отношения, константы языка $L$ и является инъективным.

Следовательно, $\pi$ является $L$-вложением
\[
\mathcal A\to\mathcal B.
\]
\end{proof}

\begin{remark}
Итак, диаграммы дают удобный способ переводить существование вложений в существование моделей некоторых множеств предложений.

\[
Diag(\mathcal A)
\]
отвечает за обычные вложения:
\[
\mathcal A\hookrightarrow\mathcal B
\quad
\Longleftrightarrow
\quad
\mathcal B \text{ можно расширить до модели }Diag(\mathcal A).
\]

\[
CDiag(\mathcal A)
\]
отвечает за элементарные вложения:
\[
\mathcal A\preccurlyeq_{\mathrm{elem}}\mathcal B
\quad
\Longleftrightarrow
\quad
\mathcal B \text{ можно расширить до модели }CDiag(\mathcal A).
\]
\end{remark}







\subsection{Теорема Лёвенгейма--Сколема о подъёме}

В этом разделе докажем восходящую теорему Лёвенгейма--Сколема. Она показывает, что если структура бесконечна, то у неё существуют элементарные расширения сколь угодно большой мощности.

\begin{theorem}[теорема Лёвенгейма--Сколема]
Пусть $\mathcal A$ — бесконечная $L$-структура. Пусть $\kappa$ — кардинал, такой что
\[
\kappa\ge \max(|A|,|L|).
\]
Тогда существует $L$-структура $\mathcal B$, такая что
\[
\mathcal A\preccurlyeq \mathcal B
\]
и
\[
|B|=\kappa.
\]
\end{theorem}

\begin{proof}
Идея доказательства состоит в том, чтобы с помощью компактности построить элементарное расширение структуры $\mathcal A$, в котором есть как минимум $\kappa$ попарно различных новых элементов. Затем с помощью нисходящей теоремы Лёвенгейма--Сколема выделяется элементарная подструктура ровно мощности $\kappa$.

Рассмотрим язык
\[
L_A=L\cup\{c_a\mid a\in A\},
\]
то есть язык, полученный из $L$ добавлением константных символов для всех элементов структуры $\mathcal A$.

Пусть $I$ — множество мощности
\[
|I|=\kappa.
\]
Добавим к языку $L_A$ новые константные символы
\[
c_i,\qquad i\in I.
\]
Полученный язык обозначим через
\[
L_{A,I}.
\]

Рассмотрим множество $L_{A,I}$-предложений
\[
\Sigma
=
CDiag(\mathcal A)
\cup
\{c_i\ne c_j\mid i,j\in I,\ i\ne j\}.
\]

Здесь $CDiag(\mathcal A)$ — полная диаграмма структуры $\mathcal A$. Первая часть теории $\Sigma$ заставляет будущую модель содержать элементарную копию структуры $\mathcal A$. Вторая часть заставляет все новые константы $c_i$ интерпретироваться попарно различными элементами.

Покажем, что $\Sigma$ конечно выполнимо.

Пусть
\[
\Sigma_0\subseteq\Sigma
\]
— конечное подмножество. Тогда в $\Sigma_0$ встречается только конечное число новых констант
\[
c_{i_1},\ldots,c_{i_n}.
\]
Поскольку структура $\mathcal A$ бесконечна, можно выбрать попарно различные элементы
\[
a_1,\ldots,a_n\in A.
\]

Возьмём естественное расширение $\mathcal A^+$ структуры $\mathcal A$ до языка $L_A$, в котором
\[
c_a^{\mathcal A^+}=a
\]
для каждого $a\in A$.

Теперь расширим $\mathcal A^+$ до языка $L_{A,I}$, интерпретируя новые константы так:
\[
c_{i_k}^{\mathcal A^+}=a_k,
\qquad
k=1,\ldots,n.
\]
Все остальные константы $c_i$, которые не встречаются в $\Sigma_0$, можно интерпретировать произвольно.

В таком расширении выполнены все предложения из $CDiag(\mathcal A)$, входящие в $\Sigma_0$, потому что на языке $L_A$ мы имеем именно естественное расширение структуры $\mathcal A$. Кроме того, все неравенства между константами $c_{i_1},\ldots,c_{i_n}$, входящие в $\Sigma_0$, выполнены, потому что элементы
\[
a_1,\ldots,a_n
\]
выбраны попарно различными.

Следовательно, всякое конечное подмножество $\Sigma_0\subseteq\Sigma$ выполнимо. Значит, $\Sigma$ конечно выполнимо.

По теореме компактности существует $L_{A,I}$-структура $\mathcal D^+$ такая, что
\[
\mathcal D^+\models\Sigma.
\]

Пусть $\mathcal D$ — редукт структуры $\mathcal D^+$ к исходному языку $L$.

Так как
\[
\mathcal D^+\models CDiag(\mathcal A),
\]
то по лемме о полной диаграмме существует элементарное вложение
\[
\pi:\mathcal A\to\mathcal D.
\]
Отождествляя $\mathcal A$ с её образом при $\pi$, можно считать, что
\[
\mathcal A\preccurlyeq\mathcal D.
\]

Кроме того, в $\mathcal D^+$ все константы $c_i$, где $i\in I$, интерпретируются попарно различными элементами. Поэтому
\[
|D|\ge |I|=\kappa.
\]

Теперь выберем подмножество
\[
S\subseteq D
\]
такое, что
\[
A\subseteq S
\]
и
\[
|S|=\kappa.
\]
Это возможно, поскольку
\[
|A|\le\kappa
\]
и
\[
|D|\ge\kappa.
\]

По нисходящей теореме Лёвенгейма--Сколема существует элементарная подструктура
\[
\mathcal B\preccurlyeq\mathcal D
\]
такая что
\[
S\subseteq B
\]
и
\[
|B|\le \max(|S|,|L|).
\]
Так как
\[
|S|=\kappa
\]
и
\[
\kappa\ge |L|,
\]
получаем
\[
|B|\le\kappa.
\]
С другой стороны, поскольку
\[
S\subseteq B
\]
и
\[
|S|=\kappa,
\]
имеем
\[
|B|\ge\kappa.
\]
Следовательно,
\[
|B|=\kappa.
\]

Осталось проверить, что $\mathcal B$ является элементарным расширением $\mathcal A$.

Мы имеем
\[
\mathcal A\preccurlyeq\mathcal D,
\qquad
\mathcal B\preccurlyeq\mathcal D,
\qquad
A\subseteq B\subseteq D.
\]
Пусть $\varphi(\bar x)$ — любая $L$-формула и пусть $\bar a\in A^n$. Тогда
\[
\mathcal A\models\varphi(\bar a)
\]
тогда и только тогда, когда
\[
\mathcal D\models\varphi(\bar a),
\]
поскольку
\[
\mathcal A\preccurlyeq\mathcal D.
\]
А так как
\[
\mathcal B\preccurlyeq\mathcal D
\]
и $\bar a\in B^n$, имеем
\[
\mathcal D\models\varphi(\bar a)
\]
тогда и только тогда, когда
\[
\mathcal B\models\varphi(\bar a).
\]
Следовательно,
\[
\mathcal A\models\varphi(\bar a)
\qquad
\Longleftrightarrow
\qquad
\mathcal B\models\varphi(\bar a).
\]
Значит,
\[
\mathcal A\preccurlyeq\mathcal B.
\]

Итак, мы построили $L$-структуру $\mathcal B$ такую, что
\[
\mathcal A\preccurlyeq\mathcal B
\]
и
\[
|B|=\kappa.
\]
\end{proof}

\begin{remark}
В доказательстве используются две основные идеи.

Во-первых, с помощью полной диаграммы $CDiag(\mathcal A)$ мы заставляем будущую модель содержать элементарную копию структуры $\mathcal A$.

Во-вторых, с помощью новых констант
\[
c_i,\qquad i\in I,
\]
и условий
\[
c_i\ne c_j
\qquad
(i\ne j)
\]
мы заставляем будущую модель иметь мощность не меньше $\kappa$.
\end{remark}

\begin{remark}
Компактность даёт модель мощности хотя бы $\kappa$, но не обязательно ровно $\kappa$. Чтобы получить модель ровно мощности $\kappa$, используется нисходящая теорема Лёвенгейма--Сколема: из большой структуры $\mathcal D$ выбирается элементарная подструктура $\mathcal B$, содержащая нужное множество $S$ мощности $\kappa$.
\end{remark}

\begin{corollary}
Пусть $\mathcal A$ — бесконечная $L$-структура. Тогда у $\mathcal A$ существуют элементарно эквивалентные структуры сколь угодно большой мощности.
\end{corollary}

\begin{proof}
По восходящей теореме Лёвенгейма--Сколема для любого кардинала
\[
\kappa\ge\max(|A|,|L|)
\]
существует структура $\mathcal B$ такая, что
\[
\mathcal A\preccurlyeq\mathcal B
\]
и
\[
|B|=\kappa.
\]
Из
\[
\mathcal A\preccurlyeq\mathcal B
\]
следует, что
\[
\mathcal A\equiv\mathcal B.
\]
\end{proof}

\begin{example}
Рассмотрим стандартную модель арифметики
\[
\omega_{s\text{-}ring}
=
\langle\omega,+,\cdot,0,1\rangle.
\]
По восходящей теореме Лёвенгейма--Сколема для любого кардинала
\[
\kappa\ge\aleph_0
\]
существует модель $\mathcal M$ такая, что
\[
\omega_{s\text{-}ring}\preccurlyeq\mathcal M
\]
и
\[
|M|=\kappa.
\]

Если
\[
\kappa>\aleph_0,
\]
то
\[
\mathcal M\not\cong \omega_{s\text{-}ring},
\]
поскольку их носители имеют разные мощности.

Элементы множества
\[
M\setminus\omega
\]
называются нестандартными натуральными числами.
\end{example}

\begin{remark}
Восходящая теорема Лёвенгейма--Сколема показывает, что логика первого порядка не контролирует мощность бесконечных моделей. Если у структуры есть бесконечная модель, то существуют элементарно эквивалентные модели больших мощностей.
\end{remark}




















\subsection{Нестандартные натуральные числа и простые элементы}


Стандартное натуральное число \(n\) внутри языка арифметики задаётся термом
\[
\underbrace{1+\cdots+1}_{n\text{ раз}}.
\]

Порядок в языке
\[
L_{s\text{-}ring}
\]
можно понимать как определимое отношение:
\[
x\le y
\quad:\Longleftrightarrow\quad
\exists z\,(x+z=y).
\]
Тогда
\[
x<y
\quad:\Longleftrightarrow\quad
x\le y\wedge x\ne y.
\]

\begin{lemma}
Пусть
\[
\mathcal M\models Th(\omega_{s\text{-}ring})
\]
и пусть
\[
a\in M\setminus\omega.
\]
Тогда для любого стандартного натурального числа \(n\in\omega\)
\[
n<a.
\]
\end{lemma}

\begin{proof}
Докажем от противного.

Предположим, что для некоторого стандартного натурального числа \(n\in\omega\)
\[
a\le n.
\]

В стандартной модели натуральных чисел истинно предложение
\[
\omega_{s\text{-}ring}
\models
\forall x
\left(
x\le n
\to
\bigvee_{i=0}^{n}x=i
\right).
\]
Оно говорит, что всякий элемент, не превосходящий \(n\), равен одному из чисел
\[
0,1,\ldots,n.
\]

Так как
\[
\mathcal M\models Th(\omega_{s\text{-}ring}),
\]
то то же самое предложение истинно и в \(\mathcal M\):
\[
\mathcal M
\models
\forall x
\left(
x\le n
\to
\bigvee_{i=0}^{n}x=i
\right).
\]

Подставляя \(a\), из предположения \(a\le n\) получаем
\[
\mathcal M\models \bigvee_{i=0}^{n}a=i.
\]
То есть
\[
a=0
\quad\text{или}\quad
a=1
\quad\text{или}\quad
\cdots
\quad\text{или}\quad
a=n.
\]

Следовательно,
\[
a\in\omega,
\]
что противоречит выбору
\[
a\in M\setminus\omega.
\]

Значит, предположение неверно, и для любого стандартного \(n\in\omega\)
\[
n<a.
\]
\end{proof}

\begin{remark}
Эта лемма говорит, что стандартные натуральные числа образуют начальный отрезок в любой нестандартной модели арифметики:
\[
0<1<2<3<\cdots<a
\]
для любого нестандартного элемента \(a\).

То есть нестандартные элементы не вставляются между стандартными числами, а идут после всех стандартных натуральных чисел.
\end{remark}

\begin{lemma}
Пусть
\[
\mathcal M\models Th(\omega_{s\text{-}ring})
\]
и пусть
\[
a\in M\setminus\omega.
\]
Тогда в \(\mathcal M\) существует нестандартное простое число \(p\), такое что
\[
a<p.
\]
\end{lemma}

\begin{proof}
Пусть \(\psi(y)\) — формула языка арифметики, выражающая, что \(y\) является простым числом. Например, можно взять формулу
\[
\psi(y):
\quad
y\ne 0\wedge y\ne 1
\wedge
\forall x\forall z
\bigl(x\cdot z=y\to (x=1\vee x=y)\bigr).
\]
В стандартной модели натуральных чисел истинно предложение о неограниченности простых чисел:
\[
\omega_{s\text{-}ring}
\models
\forall y\,\exists x\,(\psi(x)\wedge y<x).
\]

Поскольку
\[
\mathcal M\models Th(\omega_{s\text{-}ring}),
\]
это же предложение истинно и в \(\mathcal M\):
\[
\mathcal M
\models
\forall y\,\exists x\,(\psi(x)\wedge y<x).
\]

Подставим \(a\) вместо \(y\). Тогда существует элемент \(p\in M\), такой что
\[
\mathcal M\models \psi(p)
\]
и
\[
a<p.
\]

Так как \(a\) нестандартен, по предыдущей лемме для любого стандартного \(n\in\omega\)
\[
n<a.
\]
А поскольку
\[
a<p,
\]
получаем
\[
n<p
\]
для любого стандартного \(n\in\omega\).

Следовательно, \(p\) не может быть стандартным натуральным числом. Значит,
\[
p\in M\setminus\omega.
\]

Итак, \(p\) — нестандартное простое число, большее \(a\).
\end{proof}

\begin{remark}
Здесь используется только то, что неограниченность простых чисел выражается предложением первого порядка. Поэтому это свойство переносится из стандартной арифметики во всякую модель
\[
Th(\omega_{s\text{-}ring}).
\]
\end{remark}

\subsection{Элемент с бесконечно многими простыми делителями}

Следующая лемма строится уже непосредственно с помощью компактности.

\begin{lemma}
Существует нестандартная модель
\[
\mathcal M\models Th(\omega_{s\text{-}ring})
\]
и элемент
\[
\beta\in M,
\]
который имеет бесконечно много различных простых делителей.
\end{lemma}

\begin{proof}
Расширим язык
\[
L_{s\text{-}ring}
\]
новым константным символом \(c\):
\[
L_c=L_{s\text{-}ring}\cup\{c\}.
\]

Рассмотрим множество \(L_c\)-предложений
\[
\Sigma
=
Th(\omega_{s\text{-}ring})
\cup
\left\{
\exists x\,
\left(
\underbrace{x+\cdots+x}_{n\text{ раз}}=c
\right)
\mid n\in\omega,\ n\ge 1
\right\}.
\]

Формула
\[
\exists x\,
\left(
\underbrace{x+\cdots+x}_{n\text{ раз}}=c
\right)
\]
говорит, что элемент \(c\) делится на стандартное число \(n\), то есть
\[
n\mid c.
\]

Покажем, что \(\Sigma\) конечно выполнимо.

Пусть
\[
\Sigma_0\subseteq\Sigma
\]
— конечное подмножество. Тогда в \(\Sigma_0\) встречается только конечное число условий делимости:
\[
n_1\mid c,\ldots,n_k\mid c.
\]
Пусть
\[
N=\max\{n_1,\ldots,n_k\}.
\]

В стандартной модели натуральных чисел интерпретируем \(c\) как
\[
N!.
\]
Тогда \(N!\) делится на каждое из чисел
\[
1,2,\ldots,N,
\]
а значит, в частности, делится на
\[
n_1,\ldots,n_k.
\]

Все предложения из
\[
Th(\omega_{s\text{-}ring})
\]
истинны в стандартной модели натуральных чисел. Следовательно, при интерпретации
\[
c=N!
\]
структура \(\omega_{s\text{-}ring}\) является моделью конечного подмножества \(\Sigma_0\).

Значит, всякое конечное подмножество \(\Sigma\) выполнимо.

По теореме компактности множество \(\Sigma\) выполнимо. Следовательно, существует \(L_c\)-структура
\[
\mathcal M^+
\]
такая, что
\[
\mathcal M^+\models\Sigma.
\]

Пусть \(\mathcal M\) — редукт \(\mathcal M^+\) к исходному языку
\[
L_{s\text{-}ring},
\]
и положим
\[
\beta=c^{\mathcal M^+}.
\]

Так как
\[
\mathcal M^+\models Th(\omega_{s\text{-}ring}),
\]
то
\[
\mathcal M\models Th(\omega_{s\text{-}ring}).
\]

Кроме того, для каждого стандартного натурального \(n\ge 1\)
\[
\mathcal M\models n\mid \beta.
\]
В частности, если \(p\) — стандартное простое число, то
\[
p\mid\beta.
\]

Стандартных простых чисел бесконечно много, поэтому \(\beta\) имеет бесконечно много различных простых делителей.

Следовательно, \(\beta\) не может быть стандартным натуральным числом. Поэтому построенная модель \(\mathcal M\) является нестандартной.
\end{proof}

\begin{remark}
Идея доказательства состоит в том, что любое конечное число требований делимости можно выполнить в стандартной арифметике, выбрав достаточно большой факториал.

Компактность позволяет выполнить сразу все требования:
\[
1\mid c,\quad 2\mid c,\quad 3\mid c,\quad \ldots
\]
В результате появляется нестандартный элемент, делящийся на каждое стандартное натуральное число.
\end{remark}

\subsection{Нестандартная пара простых близнецов}

Пусть снова \(\psi(x)\) обозначает формулу
\[
\psi(x):
\quad
Prime(x)\wedge Prime(x+(1+1)).
\]
Она говорит, что \(x\) и \(x+2\) являются простыми числами.

\begin{lemma}
Если гипотеза о простых близнецах верна, то существует нестандартная модель
\[
\mathcal M\models Th(\omega_{s\text{-}ring})
\]
и элемент
\[
c\in M\setminus\omega,
\]
такой что
\[
\mathcal M\models \psi(c).
\]
Иными словами, существует нестандартная пара простых близнецов.
\end{lemma}

\begin{proof}
Предположим, что гипотеза о простых близнецах верна. Тогда в стандартной модели натуральных чисел формула
\[
\psi(x)
\]
имеет бесконечно много решений.

Расширим язык новым константным символом \(c\):
\[
L_c=L_{s\text{-}ring}\cup\{c\}.
\]

Рассмотрим множество \(L_c\)-предложений
\[
\Sigma
=
Th(\omega_{s\text{-}ring})
\cup
\{\psi(c)\}
\cup
\{n<c\mid n\in\omega\}.
\]

Смысл этих условий такой:
\begin{itemize}
    \item \(Th(\omega_{s\text{-}ring})\) заставляет будущую модель быть моделью истинной арифметики;
    \item \(\psi(c)\) говорит, что \(c\) и \(c+2\) являются простыми;
    \item условия \(n<c\) для всех стандартных \(n\in\omega\) заставляют \(c\) быть нестандартным числом.
\end{itemize}

Покажем, что \(\Sigma\) конечно выполнимо.

Пусть
\[
\Sigma_0\subseteq\Sigma
\]
— конечное подмножество. Тогда в \(\Sigma_0\) встречается только конечное число условий вида
\[
n<c.
\]
Пусть среди них встречаются условия
\[
n_1<c,\ldots,n_k<c.
\]
Положим
\[
N=\max\{n_1,\ldots,n_k\}.
\]

Так как по предположению существует бесконечно много пар простых близнецов, можно выбрать стандартное натуральное число \(p>N\), такое что
\[
p
\]
простое и
\[
p+2
\]
тоже простое.

В стандартной структуре \(\omega_{s\text{-}ring}\) интерпретируем
\[
c=p.
\]
Тогда
\[
\omega_{s\text{-}ring}\models \psi(c),
\]
и все конечные условия
\[
n_1<c,\ldots,n_k<c
\]
также выполнены.

Следовательно, каждое конечное подмножество \(\Sigma\) выполнимо.

По теореме компактности \(\Sigma\) выполнимо. Значит, существует \(L_c\)-структура
\[
\mathcal M^+
\]
такая, что
\[
\mathcal M^+\models\Sigma.
\]

Пусть \(\mathcal M\) — редукт \(\mathcal M^+\) к языку \(L_{s\text{-}ring}\).

Из
\[
\mathcal M^+\models Th(\omega_{s\text{-}ring})
\]
получаем
\[
\mathcal M\models Th(\omega_{s\text{-}ring}).
\]

Из условий
\[
n<c
\qquad
(n\in\omega)
\]
следует, что
\[
c^{\mathcal M^+}
\]
больше любого стандартного натурального числа. Следовательно,
\[
c^{\mathcal M^+}\in M\setminus\omega.
\]

Наконец, из
\[
\mathcal M^+\models \psi(c)
\]
получаем
\[
\mathcal M\models
Prime(c^{\mathcal M^+})
\wedge
Prime(c^{\mathcal M^+}+(1+1)).
\]

То есть в \(\mathcal M\) существует нестандартная пара простых близнецов.
\end{proof}

\begin{remark}
Это не доказательство гипотезы о простых близнецах. Гипотеза используется как предположение.

Содержательно утверждение такое: если некоторая арифметическая формула имеет бесконечно много стандартных решений, то с помощью компактности можно построить нестандартную модель, где эта формула имеет нестандартное решение.
\end{remark}



















\section{Категоричность}

\subsection{Категоричность и К-категоричность теории}

Пусть $T$ — теория в языке первого порядка $L$.

\begin{definition}
Теория $T$ называется \textbf{категоричной}, если любые две её модели изоморфны.

То есть для любых $L$-структур $\mathcal M$ и $\mathcal N$:
\[
\mathcal M\models T
\quad
\text{и}
\quad
\mathcal N\models T
\qquad
\Longrightarrow
\qquad
\mathcal M\cong\mathcal N.
\]
\end{definition}

Иными словами, теория категорична, если она задаёт свою модель однозначно с точностью до изоморфизма.

\begin{remark}
Для теорий с бесконечными моделями полная категоричность встречается редко. Причина связана с теоремами Лёвенгейма--Сколема: если теория имеет бесконечную модель, то часто она имеет модели разных бесконечных мощностей.

Но структуры разных мощностей не могут быть изоморфны. Поэтому для бесконечных структур вводят более тонкое понятие — категоричность в фиксированной мощности.
\end{remark}


\begin{definition}
Пусть $\kappa$ — кардинал. Теория $T$ называется \textbf{$\kappa$-категоричной}, если любые две модели теории $T$ мощности $\kappa$ изоморфны.

То есть если
\[
\mathcal M\models T,
\qquad
\mathcal N\models T,
\]
и
\[
|M|=|N|=\kappa,
\]
то
\[
\mathcal M\cong\mathcal N.
\]
\end{definition}

Смысл этого определения такой: теория может иметь модели разных мощностей, но в каждой фиксированной мощности $\kappa$ модель может быть единственной с точностью до изоморфизма.

\begin{remark}
Когда говорят, что ``с точностью до изоморфизма существует лишь одна модель'', важно уточнять мощность модели.

Правильная формулировка для бесконечных теорий обычно такая:
\[
\text{с точностью до изоморфизма существует ровно одна модель мощности }\kappa.
\]

Именно это и означает $\kappa$-категоричность.
\end{remark}

\subsection{Бесконечномерные векторные пространства}

Рассмотрим один из основных примеров $\kappa$-категоричных теорий: теорию бесконечномерных векторных пространств над фиксированным полем.

Для простоты будем считать, что $F$ — конечное поле. Например, можно взять
\[
F=\mathbb F_p.
\]

Рассмотрим язык
\[
L_F=\langle +,-,0,(\lambda\cdot)_{\lambda\in F}\rangle.
\]
Здесь для каждого скаляра $\lambda\in F$ в языке есть унарный функциональный символ
\[
\lambda\cdot,
\]
который интерпретируется как умножение вектора на скаляр $\lambda$.

Модели этой теории — это векторные пространства над полем $F$.

\begin{definition}
Обозначим через
\[
T_F^\infty
\]
теорию бесконечномерных векторных пространств над полем $F$.
\end{definition}

Эта теория состоит из:
\begin{enumerate}
    \item аксиом $F$-векторного пространства;
    \item аксиом, утверждающих, что размерность пространства не меньше $n$ для каждого $n\ge 1$.
\end{enumerate}

Для каждого фиксированного $n$ условие
\[
\dim_F V\ge n
\]
означает, что существуют $n$ линейно независимых векторов:
\[
\exists v_1\ldots\exists v_n\,
\operatorname{LinInd}(v_1,\ldots,v_n).
\]

Так как поле $F$ конечно, линейная независимость конечного набора выражается одной конечной формулой:
\[
\operatorname{LinInd}(v_1,\ldots,v_n)
\]
означает, что для любого ненулевого набора коэффициентов
\[
(\lambda_1,\ldots,\lambda_n)\in F^n\setminus\{(0,\ldots,0)\}
\]
выполнено
\[
\lambda_1v_1+\cdots+\lambda_nv_n\ne 0.
\]

Поскольку $F$ конечно, таких наборов коэффициентов конечное число, поэтому это действительно можно записать одной формулой первого порядка.

\begin{remark}
Если поле $F$ бесконечно, то пример тоже стандартно рассматривается, но формальная запись линейной независимости требует большей аккуратности. Для базового примера категоричности удобнее брать конечное поле, чтобы все формулы были явно конечными.
\end{remark}

\subsubsection{Классификация моделей по размерности}

Основной факт линейной алгебры:

\[
V\cong W
\qquad
\Longleftrightarrow
\qquad
\dim_F V=\dim_F W.
\]

То есть два векторных пространства над одним и тем же полем изоморфны тогда и только тогда, когда у них одинаковая размерность.

Поэтому модели теории
\[
T_F^\infty
\]
классифицируются размерностью.

\begin{remark}
Теория $T_F^\infty$ требует, чтобы размерность была бесконечной, но не фиксирует конкретную бесконечную размерность.

Поэтому у неё есть модели разных бесконечных размерностей:
\[
\aleph_0,\quad \aleph_1,\quad 2^{\aleph_0},\quad \ldots
\]
и так далее.
\end{remark}

\subsubsection{Мощность бесконечномерного векторного пространства}

Пусть $V$ — векторное пространство над конечным полем $F$, и пусть
\[
\dim_F V=\lambda,
\]
где $\lambda$ — бесконечный кардинал.

Тогда
\[
|V|=\lambda.
\]

Действительно, каждый вектор является конечной линейной комбинацией базисных векторов. Если базис имеет мощность $\lambda$, то число конечных линейных комбинаций элементов базиса над конечным полем также равно $\lambda$.

Иначе говоря, для конечного поля $F$ и бесконечномерного пространства $V$ имеем:
\[
|V|=\dim_F V.
\]

\subsection{Категоричность бесконечномерных векторных пространств}

\begin{theorem}
Пусть $F$ — конечное поле. Тогда теория
\[
T_F^\infty
\]
бесконечномерных векторных пространств над $F$ является $\kappa$-категоричной для любого бесконечного кардинала $\kappa$.
\end{theorem}

\begin{proof}
Пусть
\[
V\models T_F^\infty
\]
и
\[
W\models T_F^\infty
\]
— две модели мощности $\kappa$:
\[
|V|=|W|=\kappa.
\]

Так как $V$ и $W$ являются бесконечномерными векторными пространствами над конечным полем $F$, имеем:
\[
\dim_F V=|V|=\kappa
\]
и
\[
\dim_F W=|W|=\kappa.
\]

Следовательно,
\[
\dim_F V=\dim_F W.
\]

По теореме линейной алгебры о классификации векторных пространств по размерности получаем:
\[
V\cong W.
\]

Значит, любые две модели теории $T_F^\infty$ мощности $\kappa$ изоморфны. Поэтому $T_F^\infty$ является $\kappa$-категоричной.
\end{proof}

\begin{remark}
Важно не говорить, что у теории $T_F^\infty$ существует только одна модель вообще.

У неё есть модели разных бесконечных мощностей. Например, можно взять векторные пространства с базисами мощностей
\[
\aleph_0,\quad \aleph_1,\quad 2^{\aleph_0},
\]
и эти пространства не будут изоморфны, потому что у них разные размерности.

Правильное утверждение:
\[
\text{для каждой бесконечной мощности }\kappa
\text{ существует ровно одна модель мощности }\kappa
\text{ с точностью до изоморфизма.}
\]
\end{remark}

\subsubsection{Связь с полнотой}

Теория
\[
T_F^\infty
\]
над фиксированным конечным полем $F$ является полной.

Интуитивно это связано с тем, что любое предложение первого порядка говорит только о конечной конфигурации векторов и линейных зависимостей между ними. В бесконечномерном пространстве всегда достаточно места, чтобы реализовать любую конечную линейно-алгебраическую конфигурацию, если она вообще совместима с аксиомами векторного пространства.

Поэтому все бесконечномерные векторные пространства над одним и тем же конечным полем элементарно эквивалентны:
\[
V\models T_F^\infty,
\qquad
W\models T_F^\infty
\qquad
\Longrightarrow
\qquad
V\equiv W.
\]

Но они не обязательно изоморфны: изоморфизм зависит от размерности.

\begin{remark}
Более продвинуто это можно объяснить через исключение кванторов для теории векторных пространств в подходящем языке. Но для текущего уровня достаточно использовать линейно-алгебраическую интуицию: все конечные свойства в бесконечномерных пространствах устроены одинаково, а изоморфный тип определяется размерностью.
\end{remark}


Теория бесконечномерных векторных пространств над конечным полем показывает главный смысл $\kappa$-категоричности:

\[
\kappa\text{-категоричность}
\quad
=
\quad
\text{единственность модели данной мощности }\kappa.
\]

В этом примере:

\[
|V|=\kappa
\]
определяет
\[
\dim_F V=\kappa,
\]
а размерность определяет изоморфный тип векторного пространства.

Поэтому
\[
T_F^\infty
\]
имеет ровно одну модель мощности $\kappa$ с точностью до изоморфизма для каждого бесконечного кардинала $\kappa$.


---------------------------------------------------------------------------------------------------\\
\textbf{ДАЛЬШЕ БОГА НЕТ!!!}


\section{Исключение кванторов}
\section{Насыщенные модели}
\section{Устойчивость}

\section{Задачник по теории моделей}

\subsection{Блок 1. Основы теории моделей}

\subsubsection*{Задача №1}
\addcontentsline{toc}{subsubsection}{Задача №1}

\par\medskip
\textbf{Условие.}
\par\smallskip
Пусть $\mathfrak A$ и $\mathfrak B$ — $L$-структуры. Гомоморфизмом $L$-структур называется отображение
\[
\pi:A\to B,
\]
такое, что оно сохраняет отношения, функции и константы языка в следующем смысле:

\begin{enumerate}
\item Для любого символа отношения $R^n$ и любых $a_1,\dots,a_n\in A$:
\[
(a_1,\dots,a_n)\in R^{\mathfrak A}
\Rightarrow
(\pi(a_1),\dots,\pi(a_n))\in R^{\mathfrak B}.
\]

\item Для любого функционального символа $f^n$ и любых $a_1,\dots,a_n\in A$:
\[
\pi(f^{\mathfrak A}(a_1,\dots,a_n))
=
f^{\mathfrak B}(\pi(a_1),\dots,\pi(a_n)).
\]

\item Для любого константного символа $c$:
\[
\pi(c^{\mathfrak A})=c^{\mathfrak B}.
\]

\end{enumerate}

Нужно проверить, что для языков групп и колец это совпадает с обычным определением гомоморфизма групп и колец. Также нужно сравнить $L$-гомоморфизмы с вложениями, доказать, что биективный гомоморфизм функционального языка является изоморфизмом, и построить пример, показывающий, что условие отсутствия символов отношений опустить нельзя.

\par\medskip
\textbf{Решение.}
\par\smallskip

\textbf{1. Случай языка групп.}

Рассмотрим язык групп:
\[
L_{gp}=\langle \cdot^2,(\blank^{-1})^1,1^0\rangle.
\]

В этом языке есть два функциональных символа и один константный символ. Символов отношений нет, поэтому условие для отношений в определении гомоморфизма не появляется.

Условие для бинарной функции $\cdot$ даёт:
\[
\pi(a\cdot b)=\pi(a)\cdot\pi(b).
\]

Условие для унарной функции взятия обратного элемента даёт:
\[
\pi(a^{-1})=(\pi(a))^{-1}.
\]

Условие для константы $1$ даёт:
\[
\pi(1)=1.
\]

Таким образом, $L_{gp}$-гомоморфизм — это обычный гомоморфизм групп. В стандартном курсе алгебры обычно достаточно требовать только сохранение умножения, а сохранение единицы и обратного элемента выводится. Здесь эти условия появляются сразу, потому что $1$ и $\blank^{-1}$ входят в язык.

\textbf{2. Случай языка колец.}

Рассмотрим язык колец:
\[
L_{ring}=\langle +^2,\cdot^2,-^1,0^0,1^0\rangle.
\]

В этом языке также нет символов отношений. Поэтому гомоморфизм должен сохранять только функции и константы.

Для сложения получаем:
\[
\pi(a+b)=\pi(a)+\pi(b).
\]

Для умножения:
\[
\pi(a\cdot b)=\pi(a)\cdot\pi(b).
\]

Для унарного минуса:
\[
\pi(-a)=-\pi(a).
\]

Для констант:
\[
\pi(0)=0,
\qquad
\pi(1)=1.
\]

Следовательно, $L_{ring}$-гомоморфизм — это обычный гомоморфизм колец с единицей.

\textbf{3. Связь гомоморфизмов и вложений.}

Сравним определение гомоморфизма с определением вложения.

Для функций условия совпадают: и гомоморфизм, и вложение должны сохранять результат применения функции. То есть если сначала применить функцию в $\mathfrak A$, а потом перенести результат с помощью $\pi$, получится то же самое, что если сначала перенести все аргументы в $\mathfrak B$, а потом применить функцию там.

Для констант условия тоже совпадают: константа из $\mathfrak A$ должна перейти в соответствующую константу из $\mathfrak B$.

Различие возникает для отношений. Гомоморфизм требует только сохранения истинности отношения:
\[
(a_1,\dots,a_n)\in R^{\mathfrak A}
\Rightarrow
(\pi(a_1),\dots,\pi(a_n))\in R^{\mathfrak B}.
\]

Вложение требует более сильного условия:
\[
(a_1,\dots,a_n)\in R^{\mathfrak A}
\iff
(\pi(a_1),\dots,\pi(a_n))\in R^{\mathfrak B}.
\]

То есть вложение сохраняет не только истинность отношения, но и его ложность. Кроме того, вложение по определению должно быть инъективным.

Следовательно:
\[
\text{вложение}
\Rightarrow
\text{инъективный гомоморфизм}.
\]

Обратное в общем случае неверно: инъективный гомоморфизм может сохранять отношение только в одну сторону, а для вложения нужна эквивалентность.

\textbf{4. Биективный гомоморфизм функционального языка.}

Пусть язык $L$ является функциональным, то есть в нём есть только функциональные символы и константы, но нет символов отношений.

Пусть
\[
\pi:\mathfrak A\to\mathfrak B
\]
— биективный гомоморфизм. Нужно показать, что $\pi$ является изоморфизмом.

Так как $\pi$ биективно, для него существует обратное отображение
\[
\sigma:B\to A.
\]

По определению изоморфизма нужно проверить, что обратное отображение $\sigma$ тоже сохраняет функции и константы.

Пусть $f$ — функциональный символ арности $n$, а $b_1,\dots,b_n\in B$. Так как $\pi$ сюръективно, можно выбрать элементы $a_1,\dots,a_n\in A$ так, что
\[
\pi(a_i)=b_i
\]
для всех $i=1,\dots,n$.

Так как $\pi$ является гомоморфизмом, имеем:
\[
\pi(f^{\mathfrak A}(a_1,\dots,a_n))
=
f^{\mathfrak B}(\pi(a_1),\dots,\pi(a_n)).
\]

После подстановки $\pi(a_i)=b_i$ получаем:
\[
\pi(f^{\mathfrak A}(a_1,\dots,a_n))
=
f^{\mathfrak B}(b_1,\dots,b_n).
\]

Теперь применим к обеим частям обратное отображение $\sigma$. Тогда левая часть вернётся обратно в $f^{\mathfrak A}(a_1,\dots,a_n)$, а справа получится:
\[
\sigma(f^{\mathfrak B}(b_1,\dots,b_n))
=
f^{\mathfrak A}(a_1,\dots,a_n).
\]

Но $a_i=\sigma(b_i)$, значит:
\[
\sigma(f^{\mathfrak B}(b_1,\dots,b_n))
=
f^{\mathfrak A}(\sigma(b_1),\dots,\sigma(b_n)).
\]

Это ровно условие сохранения функции для отображения $\sigma$.

Для констант проверка короче. Так как $\pi$ сохраняет константу $c$, имеем $\pi(c^{\mathfrak A})=c^{\mathfrak B}$. Применяя $\sigma$, получаем:
\[
\sigma(c^{\mathfrak B})=c^{\mathfrak A}.
\]

Значит, $\sigma$ также является гомоморфизмом. Следовательно, $\pi$ является обратимым гомоморфизмом, то есть изоморфизмом.

\textbf{5. Почему нельзя опустить условие отсутствия отношений.}

Построим пример, где есть символ отношения.

Пусть язык
\[
L=\langle R^1\rangle
\]
состоит из одного унарного символа отношения.

Рассмотрим две $L$-структуры с одинаковым носителем:
\[
A=B=\{1\}.
\]

Пусть в первой структуре отношение пустое:
\[
R^{\mathfrak A}=\varnothing,
\]
а во второй структуре отношение содержит единственный элемент:
\[
R^{\mathfrak B}=\{1\}.
\]

Рассмотрим отображение $\pi:A\to B$, заданное правилом $\pi(1)=1$. Это отображение биективно.

Проверим, что оно является гомоморфизмом. Для отношения нужно проверить только импликацию:
\[
a\in R^{\mathfrak A}
\Rightarrow
\pi(a)\in R^{\mathfrak B}.
\]

Но $R^{\mathfrak A}$ пусто. Значит, в первой структуре нет элементов, для которых левая часть условия выполняется. Поэтому импликация истинна автоматически.

Следовательно, $\pi$ — биективный гомоморфизм.

Однако $\pi$ не является изоморфизмом. Действительно, в первой структуре элемент $1$ не лежит в отношении, а во второй лежит:
\[
1\notin R^{\mathfrak A},
\qquad
1\in R^{\mathfrak B}.
\]

Значит, отношение не сохраняется в обратную сторону. Поэтому структуры не изоморфны.

Итак, для функциональных языков биективный гомоморфизм является изоморфизмом, но при наличии символов отношений это утверждение уже неверно.


\subsubsection*{Задача №2}
\addcontentsline{toc}{subsubsection}{Задача №2}

\par\medskip
\textbf{Условие.}
\par\smallskip

Пусть
\[
L=\langle f^1,c^0\rangle,
\]
где $f$ — унарный функциональный символ, а $c$ — константный символ.

Описать с точностью до изоморфизма все $L$-структуры на носителе
\[
A=\{1,2\}.
\]

\par\medskip
\textbf{Решение.}
\par\smallskip

Для задания $L$-структуры необходимо указать интерпретацию константы и интерпретацию функции.

Константа может интерпретироваться как один из элементов носителя:
\[
c^{\mathfrak A}=1
\qquad \text{или} \qquad
c^{\mathfrak A}=2.
\]

Функция
\[
f^{\mathfrak A}:A\to A
\]
должна каждому элементу множества $\{1,2\}$ сопоставлять элемент того же множества.

Так как для каждого из двух элементов имеется два возможных значения образа, получаем
\[
2^2=4
\]
различные функции:

\begin{enumerate}
    \item Тождественная функция:
    \[
    f(1)=1,
    \qquad
    f(2)=2.
    \]

    \item Константная функция в $1$:
    \[
    f(1)=1,
    \qquad
    f(2)=1.
    \]

    \item Перестановка элементов:
    \[
    f(1)=2,
    \qquad
    f(2)=1.
    \]

    \item Константная функция в $2$:
    \[
    f(1)=2,
    \qquad
    f(2)=2.
    \]
\end{enumerate}

Следовательно, до учёта изоморфизма имеется
\[
2\cdot4=8
\]
различных структур.

Далее необходимо учитывать возможное переименование элементов носителя.

Изоморфизм должен сохранять интерпретацию константы и функции. Поэтому некоторые из полученных структур оказываются изоморфными друг другу.

Например, структуры

\[
c^{\mathfrak A}=1,
\qquad
f(1)=1,\quad f(2)=1,
\]

и

\[
c^{\mathfrak B}=2,
\qquad
f(1)=2,\quad f(2)=2
\]

изоморфны при перестановке элементов
\[
1\leftrightarrow2.
\]

Аналогично изоморфны структуры, отличающиеся только заменой обозначений элементов $1$ и $2$.

Таким образом, задача сводится к классификации полученных восьми структур по классам изоморфизма.

\subsubsection*{Задача №3}
\addcontentsline{toc}{subsubsection}{Задача №3}

\par\medskip
\textbf{Условие.}
\par\smallskip

Найти все автоморфизмы структуры
\[
\mathbb{Q}_{adgp}=\langle \mathbb{Q},+,-,0\rangle.
\]

\par\medskip
\textbf{Решение.}
\par\smallskip

Пусть
\[
\pi:\mathbb{Q}\to\mathbb{Q}
\]
— автоморфизм структуры $\mathbb{Q}_{adgp}$.

Так как $\pi$ сохраняет сложение, для любых $x,y\in\mathbb{Q}$ выполнено
\[
\pi(x+y)=\pi(x)+\pi(y).
\]

Положим
\[
\pi(1)=q.
\]

Тогда для любого натурального $n$
\[
\pi(n)=nq.
\]

Для отрицательных целых чисел также получаем
\[
\pi(-n)=-nq.
\]

Следовательно, для любого $m\in\mathbb{Z}$
\[
\pi(m)=mq.
\]

Теперь пусть
\[
x=\frac{m}{n},
\]
где $m\in\mathbb{Z}$, $n\in\mathbb{N}$.

Тогда в аддитивной группе $\mathbb{Q}$ выполнено
\[
nx=m.
\]

Применяя $\pi$, получаем
\[
n\pi(x)=\pi(m).
\]

Так как $\pi(m)=mq$, отсюда следует
\[
\pi(x)=\frac{m}{n}q.
\]

Значит, для любого $x\in\mathbb{Q}$
\[
\pi(x)=qx.
\]

Если $q=0$, то $\pi$ не является биекцией. Поэтому
\[
q\neq0.
\]

Обратно, для любого $q\in\mathbb{Q}\setminus\{0\}$ отображение
\[
\pi_q(x)=qx
\]
является биекцией $\mathbb{Q}\to\mathbb{Q}$ и сохраняет сложение, ноль и взятие противоположного элемента.

Следовательно,
\[
Aut(\mathbb{Q}_{adgp})
=
\{\pi_q\mid \pi_q(x)=qx,\ q\in\mathbb{Q}\setminus\{0\}\}.
\]


\subsubsection*{Задача №4}
\addcontentsline{toc}{subsubsection}{Задача №4}

\par\medskip
\textbf{Условие.}
\par\smallskip

Найти автоморфизм $\pi$ структуры
\[
\langle \mathbb{R},<\rangle
\]
такой, что
\[
\pi(0)=1,
\qquad
\pi(1)=5.
\]
Является ли это отображение также автоморфизмом структуры
\[
\langle \mathbb{R},+\rangle?
\]

\par\medskip
\textbf{Решение.}
\par\smallskip

Автоморфизм структуры $\langle \mathbb{R},<\rangle$ — это биекция
\[
\pi:\mathbb{R}\to\mathbb{R},
\]
сохраняющая порядок.

Достаточно взять линейное отображение
\[
\pi(x)=4x+1.
\]

Тогда
\[
\pi(0)=1,
\qquad
\pi(1)=5.
\]

Так как коэффициент при $x$ положительный, отображение $\pi$ строго возрастает. Кроме того, оно является биекцией $\mathbb{R}\to\mathbb{R}$. Следовательно, $\pi$ является автоморфизмом структуры $\langle \mathbb{R},<\rangle$.

Теперь проверим, является ли $\pi$ автоморфизмом структуры $\langle \mathbb{R},+\rangle$.

Для этого должно выполняться равенство
\[
\pi(x+y)=\pi(x)+\pi(y)
\]
для любых $x,y\in\mathbb{R}$.

Но для $x=y=0$ получаем
\[
\pi(0+0)=\pi(0)=1,
\]
а с другой стороны
\[
\pi(0)+\pi(0)=1+1=2.
\]

Следовательно,
\[
\pi(0+0)\neq \pi(0)+\pi(0).
\]

Значит, $\pi$ не сохраняет сложение и не является автоморфизмом структуры $\langle \mathbb{R},+\rangle$.

Ответ: отображение
\[
\pi(x)=4x+1
\]
является автоморфизмом структуры $\langle \mathbb{R},<\rangle$, но не является автоморфизмом структуры $\langle \mathbb{R},+\rangle$.


\subsubsection*{Задача №5}
\addcontentsline{toc}{subsubsection}{Задача №5}

\par\medskip
\textbf{Условие.}
\par\smallskip

Пусть $\mathfrak A^+$ — обогащение структуры $\mathfrak A$, и пусть $\pi$ — автоморфизм структуры $\mathfrak A^+$.

Доказать, что $\pi$ также является автоморфизмом структуры $\mathfrak A$. Показать, что обратное утверждение неверно.

\par\medskip
\textbf{Решение.}
\par\smallskip

Обогащение $\mathfrak A^+$ получается из структуры $\mathfrak A$ добавлением новых символов в язык. При этом носитель остаётся тем же самым, а все старые символы структуры $\mathfrak A$ сохраняют свои интерпретации.

Если $\pi$ является автоморфизмом $\mathfrak A^+$, то $\pi$ сохраняет все символы расширенного языка. В частности, она сохраняет все символы исходного языка структуры $\mathfrak A$.

Следовательно, $\pi$ является автоморфизмом структуры $\mathfrak A$.

Обратное утверждение неверно.

Например, рассмотрим структуру
\[
\mathfrak A=\langle \mathbb Z,<\rangle.
\]

Отображение
\[
\pi(n)=n+1
\]
является автоморфизмом $\mathfrak A$, так как оно является биекцией $\mathbb Z\to\mathbb Z$ и сохраняет порядок.

Теперь обогатим структуру константой $0$:
\[
\mathfrak A^+=\langle \mathbb Z,<,0\rangle.
\]

Автоморфизм обогащённой структуры обязан сохранять интерпретацию константы $0$. Значит, должно выполняться
\[
\pi(0)=0.
\]

Но для отображения $\pi(n)=n+1$ имеем
\[
\pi(0)=1.
\]

Поэтому $\pi$ не является автоморфизмом структуры $\mathfrak A^+$.

Следовательно, всякий автоморфизм обогащения является автоморфизмом исходной структуры, но не всякий автоморфизм исходной структуры остаётся автоморфизмом её обогащения.

\subsubsection*{Задача №6}
\addcontentsline{toc}{subsubsection}{Задача №6}

\par\medskip
\textbf{Условие.}
\par\smallskip

Доказать, что $Aut(\mathfrak A)$ образует группу относительно композиции отображений.

\par\medskip
\textbf{Решение.}
\par\smallskip

Множество $Aut(\mathfrak A)$ состоит из всех автоморфизмов структуры $\mathfrak A$, то есть из всех изоморфизмов структуры $\mathfrak A$ в саму себя.

Покажем, что относительно композиции отображений это множество является группой.

Во-первых, композиция двух автоморфизмов снова является автоморфизмом. Действительно, если
\[
\pi,\sigma\in Aut(\mathfrak A),
\]
то отображение
\[
\sigma\circ\pi:A\to A
\]
является биекцией. Кроме того, оно сохраняет все отношения, функции и константы языка, поскольку сначала их сохраняет $\pi$, а затем их сохраняет $\sigma$.

Следовательно,
\[
\sigma\circ\pi\in Aut(\mathfrak A).
\]

Во-вторых, композиция отображений ассоциативна. Поэтому для любых
\[
\pi,\sigma,\rho\in Aut(\mathfrak A)
\]
выполнено
\[
(\rho\circ\sigma)\circ\pi=\rho\circ(\sigma\circ\pi).
\]

В-третьих, тождественное отображение
\[
id_A:A\to A
\]
является автоморфизмом структуры $\mathfrak A$, так как оно сохраняет все элементы, а значит, сохраняет все отношения, функции и константы языка.

Это отображение является нейтральным элементом относительно композиции.

В-четвёртых, у каждого автоморфизма есть обратный автоморфизм. Если
\[
\pi\in Aut(\mathfrak A),
\]
то $\pi$ является изоморфизмом структуры $\mathfrak A$ в себя. Поэтому обратное отображение
\[
\pi^{-1}:A\to A
\]
также является автоморфизмом структуры $\mathfrak A$.

Итак, множество $Aut(\mathfrak A)$ замкнуто относительно композиции, композиция ассоциативна, существует нейтральный элемент и у каждого элемента есть обратный.

Следовательно, $Aut(\mathfrak A)$ является группой относительно композиции отображений.


\chapter*{Глава 8 покрывает курс А.Станкевича "Дискретная математика s4 вторая половина"}
\addcontentsline{toc}{chapter}{Глава 8 покрывает курс А.Станкевича "Дискретная математика s4 вторая половина"}


\chapter{Теория вычислимости}

\section{Понятие алгоритма}
\section{Частично-рекурсивные функции}
\section{Машины Тьюринга}
\section{\texorpdfstring{$\lambda$}{lambda}-исчисление}
\section{Тезис Черча--Тьюринга}
\section{Разрешимые множества}
\section{Перечислимые множества}
\section{Неразрешимые задачи}
\section{Теорема Райса}
\section{Арифметическая иерархия}

\section{Задачник к главе}


\chapter*{Глава 9 покрывает курс Д.Штукенберга "Теория типов"}
\addcontentsline{toc}{chapter}{Глава 9 покрывает курс Д.Штукенберга "Теория типов"}

\chapter{Теория типов}

\section{Простая теория типов}
\section{Типизированное \texorpdfstring{$\lambda$}{lambda}-исчисление}
\section{Система F}
\section{Зависимые типы}
\section{Изоморфизм Карри--Ховарда}
\section{Типы как доказательства}
\section{Конструктивные основания математики}
\section{Coq, Agda и Lean}

\section{Задачник к главе}

\chapter{Основания математики}

\section{Логицизм}
\section{Формализм Гильберта}
\section{Интуиционизм Брауэра}
\section{Теоретико-множественные основания}
\section{Категориальные основания}
\section{Унивалентные основания}

\section{Задачник к главе}

\end{document}
